% Formuler les idées essentielles d’un texte à résumer — Français Tle Tech — Cours signé OnBuch+
% Programme officiel MINESEC. Compiler avec : tectonic N25-formuler-idees-essentielles-texte-resumer.tex
\newcommand{\DOCMATIERE}{Français}
\newcommand{\DOCNIVEAU}{Tle Tech}
\newcommand{\DOCMODULE}{Français – classe de Terminale technique}
\newcommand{\DOCLECON}{N25}
\newcommand{\DOCTITRE}{Formuler les idées essentielles d’un texte à résumer}
% OnBuch+ — préambule « cours signés » ENRICHI (compilé avec tectonic / XeLaTeX)
% Système visuel « Pop » d'OnBuch+ : fond crème (jamais blanc), bordures encre
% épaisses, ombres « dures » décalées, titres Archivo Black en capitales.
% Version MATHÉMATIQUES : graphes (pgfplots/tikz), boîtes pédagogiques
% (définition, propriété, méthode, attention, le savais-tu, exercice résolu,
% exercice, TP, bilan), page de garde et signature OnBuch+ ancrée en bas.
%
% Macros attendues dans main.tex AVANT \input{preamble} :
%   \DOCMATIERE  \DOCNIVEAU  \DOCMODULE  \DOCLECON (numéro)  \DOCTITRE
\documentclass[11pt]{article}

\usepackage{fontspec}
\usepackage[french]{babel}
\usepackage[a4paper,margin=2cm,top=3.4cm,bottom=2.8cm,headheight=1.4cm,footskip=1.3cm]{geometry}
\usepackage{amsmath,amssymb,amsfonts}
\usepackage{mathtools} % \xmapsto, \coloneqq... souvent utilisés par les modèles
\usepackage{cancel} % simplifications barrées (bilans de forces)
% \abs{z} (module, valeur absolue) et \norm{u} (norme), utilisés par les modèles
\providecommand{\abs}{}\let\abs\relax\DeclarePairedDelimiter{\abs}{\lvert}{\rvert}
\providecommand{\norm}{}\let\norm\relax\DeclarePairedDelimiter{\norm}{\lVert}{\rVert}
\usepackage[table]{xcolor} % [table] : fournit aussi \rowcolors (alternance de lignes)
\usepackage{pagecolor}
\usepackage{tikz}
\usetikzlibrary{arrows.meta,positioning,calc,shapes.geometric,shapes.misc,shapes.arrows,decorations.pathmorphing,decorations.pathreplacing,decorations.markings,patterns,shadows,mindmap,trees,fit,backgrounds,matrix,intersections,angles,quotes}
\usepackage{pgfplots}
\usepackage[version=4]{mhchem} % équations nucléaires \ce{^{235}_{92}U}
\usepackage[european,siunitx]{circuitikz} % schémas électriques
\pgfplotsset{compat=1.18}
\usepgfplotslibrary{fillbetween,groupplots} % aires entre courbes (calcul intégral)
\usepackage{enumitem}
\usepackage{fancyhdr}
% [most] charge toutes les bibliothèques tcolorbox (listings, minted, raster,
% documentation...) et gonfle inutilement la mémoire TeX sur un document long
% (~300 boîtes) : on ne charge que ce qui est réellement utilisé.
\usepackage[skins,breakable]{tcolorbox}
\usepackage{graphicx}
\usepackage{colortbl}
\usepackage{booktabs}
\usepackage{tabularx}
\usepackage{multirow}
\usepackage{diagbox} % case d'en-tête diagonale des tableaux à double entrée
% Colonnes extensibles centrée / gauche / droite (C, L, R), souvent utilisées par les modèles
\newcolumntype{C}{>{\centering\arraybackslash}X}
\newcolumntype{L}{>{\raggedright\arraybackslash}X}
\newcolumntype{R}{>{\raggedleft\arraybackslash}X}
\usepackage{array}
\usepackage{multicol}
\usepackage{siunitx}
% parse-numbers=false : accepte aussi \SI{2,5 \times 10^{-3}}{...}
\sisetup{locale=FR,output-decimal-marker={,},group-minimum-digits=4,parse-numbers=false}
\DeclareSIUnit\ohm{\ensuremath{\symup{\Omega}}} % U+2126 absent de la police
\DeclareSIUnit\division{div} % réglages d'oscilloscope (ms/div, V/div)
\DeclareSIUnit\tr{tr} % tours (vitesse de rotation, ex. \SI{1500}{\tr\per\minute})
\DeclareSIUnit\molar{M} % concentration molaire (ex. \SI{3}{\molar}, \SI{1}{\micro\molar})
\DeclareSIUnit\an{an} % année (le modèle confond parfois avec \year, un compteur TeX) — ex. \SI{5}{\kilo\meter\per\an}
\DeclareSIUnit\FCFA{FCFA} % franc CFA (ex. \SI{500}{\FCFA\per\kilo\gram})
\DeclareSIUnit\degreeBrix{\textdegree Brix} % teneur en sucre d'un sirop/jus (réfractométrie)
\DeclareSIUnit\month{mois} % le modèle utilise parfois \month, un compteur TeX, comme unité de durée
\DeclareSIUnit\microliter{\micro\liter} % le modèle écrit parfois \microliter en un seul token
\DeclareSIUnit\million{millions} % dénombrement (ex. \SI{15}{\million\per\mL} spermatozoïdes)
\usepackage{qrcode}
% \degree : commande fréquemment utilisée par les modèles pour un angle ou une
% température en dehors de \SI{}{} (non fournie par les paquets chargés).
\providecommand{\degree}{^{\circ}}
% \qed (fin de démonstration), utilisé par les modèles sans amsthm
\providecommand{\qed}{\hfill$\square$}
% \barre : double barre des valeurs interdites dans un tableau de variations (tkz-tab)
\providecommand{\barre}{\ensuremath{\|}}
% environnement proof (amsthm non chargé) utilisé par les modèles
\ifdefined\proof\else\newenvironment{proof}[1][Démonstration]{\par\noindent\textit{#1.}\ }{\qed\par}\fi
\usepackage{lastpage}
\usepackage{hyperref}
\hypersetup{hidelinks}

% ============================================================
% Palette « Pop » OnBuch+ — mêmes hex que la classe P (plus_ui.dart)
% ============================================================
\definecolor{poporange}{HTML}{F59321}
\definecolor{popdark}{HTML}{E07A0C}
\definecolor{popcream}{HTML}{FFF6E8}
\definecolor{popink}{HTML}{1C1714}
\definecolor{poppurple}{HTML}{7A5AE0}
\definecolor{popgreen}{HTML}{2E9E6A}
\definecolor{poppink}{HTML}{E0457B}
\definecolor{popblue}{HTML}{2F7DE1}
\definecolor{popgold}{HTML}{C98A12}
\definecolor{popmuted}{HTML}{8A7F76}
\definecolor{popline}{HTML}{EADFD2}
\definecolor{popgray}{HTML}{8A7F76}
\colorlet{popgrey}{popgray}
% Alias tolérants (le modèle nomme parfois les couleurs différemment)
\colorlet{popwhite}{white}
\colorlet{popblack}{popink}
\colorlet{popred}{poppink}
\colorlet{popyellow}{popgold}
% Teintes claires (fonds de tableaux/encadrés) — le modèle nomme parfois
% les couleurs avec un suffixe L (ex. popblueL) sans les avoir définies.
\colorlet{poporangeL}{poporange!20}
\colorlet{popdarkL}{popdark!20}
\colorlet{poppurpleL}{poppurple!20}
\colorlet{popgreenL}{popgreen!20}
\colorlet{poppinkL}{poppink!20}
\colorlet{popblueL}{popblue!20}
\colorlet{popgoldL}{popgold!20}
\colorlet{popredL}{popred!20}
\colorlet{popgrayL}{popgray!20}
% Teintes claires pour fonds de tableaux / zones
\colorlet{popblueL}{popblue!10!white}
\colorlet{poporangeL}{poporange!14!white}
\colorlet{popgreenL}{popgreen!12!white}
\colorlet{poppurpleL}{poppurple!12!white}
\colorlet{poppinkL}{poppink!12!white}
\colorlet{popgoldL}{popgold!14!white}

\pagecolor{popcream}

% ============================================================
% Typographie
% ============================================================
\defaultfontfeatures{Ligatures=TeX}
\setmainfont{PlusJakartaSans-Medium.ttf}[Path=../../../fonts/,BoldFont=PlusJakartaSans-Bold.ttf]
% Maths en sans-serif (Fira Math) avec chiffres et lettres droites de Plus
% Jakarta, pour que les formules se fondent dans le texte.
\usepackage[math-style=ISO,bold-style=ISO]{unicode-math}
\setmathfont{FiraMath-Regular.otf}[Path=../../../fonts/]
\setmathfont{PlusJakartaSans-Medium.ttf}[Path=../../../fonts/,range={up/num,up/{Latin,latin},\mathup}]
\usepackage{icomma} % virgule décimale sans espace en mode math
% Caractères mathématiques Unicode absents des polices de texte (Plus Jakarta,
% Archivo Black) : rendus par leur équivalent mathématique, en texte comme en
% formule — sinon ils s'affichent en carré vide (ex. « Arithmétique dans ℤ »).
\usepackage{newunicodechar}
\newunicodechar{ℝ}{\ensuremath{\mathbb{R}}}
\newunicodechar{ℤ}{\ensuremath{\mathbb{Z}}}
\newunicodechar{ℂ}{\ensuremath{\mathbb{C}}}
\newunicodechar{ℕ}{\ensuremath{\mathbb{N}}}
\newunicodechar{ℚ}{\ensuremath{\mathbb{Q}}}
\newunicodechar{𝔻}{\ensuremath{\mathbb{D}}}
\newunicodechar{α}{\ensuremath{\alpha}}
\newunicodechar{β}{\ensuremath{\beta}}
\newunicodechar{γ}{\ensuremath{\gamma}}
\newunicodechar{δ}{\ensuremath{\delta}}
\newunicodechar{ε}{\ensuremath{\varepsilon}}
\newunicodechar{θ}{\ensuremath{\theta}}
\newunicodechar{λ}{\ensuremath{\lambda}}
\newunicodechar{μ}{\ensuremath{\mu}}
\newunicodechar{σ}{\ensuremath{\sigma}}
\newunicodechar{τ}{\ensuremath{\tau}}
\newunicodechar{φ}{\ensuremath{\varphi}}
\newunicodechar{ψ}{\ensuremath{\psi}}
\newunicodechar{ω}{\ensuremath{\omega}}
\newunicodechar{Σ}{\ensuremath{\Sigma}}
% majuscules produites par \MakeUppercase dans les titres (α-aminés -> Α)
\newunicodechar{Α}{\ensuremath{\alpha}}
\newunicodechar{Β}{\ensuremath{\beta}}
\newunicodechar{Γ}{\ensuremath{\Gamma}}
\newunicodechar{Δ}{\ensuremath{\Delta}}
\newunicodechar{Ε}{\ensuremath{\varepsilon}}
\newunicodechar{Θ}{\ensuremath{\Theta}}
\newunicodechar{Λ}{\ensuremath{\Lambda}}
\newunicodechar{Μ}{\ensuremath{\mu}}
\newunicodechar{Τ}{\ensuremath{\tau}}
\newunicodechar{Φ}{\ensuremath{\Phi}}
\newunicodechar{Ψ}{\ensuremath{\Psi}}
\newunicodechar{Ω}{\ensuremath{\Omega}}
\newunicodechar{↦}{\ensuremath{\mapsto}}
\newunicodechar{∈}{\ensuremath{\in}}
\newunicodechar{∉}{\ensuremath{\notin}}
\newunicodechar{∧}{\ensuremath{\wedge}}
\newunicodechar{⇔}{\ensuremath{\Leftrightarrow}}
\newunicodechar{⟺}{\ensuremath{\Longleftrightarrow}}
\newunicodechar{⇒}{\ensuremath{\Rightarrow}}
\newunicodechar{⟹}{\ensuremath{\Longrightarrow}}
\newunicodechar{∘}{\ensuremath{\circ}}
\newunicodechar{∪}{\ensuremath{\cup}}
\newunicodechar{∩}{\ensuremath{\cap}}
\newunicodechar{≡}{\ensuremath{\equiv}}
\newunicodechar{⊂}{\ensuremath{\subset}}
\newunicodechar{⊥}{\ensuremath{\perp}}
\newunicodechar{∃}{\ensuremath{\exists}}
\newunicodechar{∀}{\ensuremath{\forall}}
\newunicodechar{∛}{\ensuremath{\sqrt[3]{\,}}}
\newunicodechar{∜}{\ensuremath{\sqrt[4]{\,}}}
\newunicodechar{⃗}{\ensuremath{{}^{\to}}}
\newunicodechar{ⁿ}{\ifmmode^{n}\else\textsuperscript{\ensuremath{n}}\fi}
\newunicodechar{ˣ}{\ifmmode^{x}\else\textsuperscript{\ensuremath{x}}\fi}
\newunicodechar{ᵇ}{\ifmmode^{b}\else\textsuperscript{\ensuremath{b}}\fi}
\newunicodechar{ʳ}{\ifmmode^{r}\else\textsuperscript{\ensuremath{r}}\fi}
\newunicodechar{ᵘ}{\ifmmode^{u}\else\textsuperscript{\ensuremath{u}}\fi}
\newunicodechar{ʸ}{\ifmmode^{y}\else\textsuperscript{\ensuremath{y}}\fi}
\newunicodechar{ᵃ}{\ifmmode^{a}\else\textsuperscript{\ensuremath{a}}\fi}
\newunicodechar{ᵉ}{\ifmmode^{e}\else\textsuperscript{\ensuremath{e}}\fi}
\newunicodechar{ᵏ}{\ifmmode^{k}\else\textsuperscript{\ensuremath{k}}\fi}
\newunicodechar{ᵖ}{\ifmmode^{p}\else\textsuperscript{\ensuremath{p}}\fi}
\newunicodechar{ᴺ}{\ifmmode^{N}\else\textsuperscript{\ensuremath{N}}\fi}
\newunicodechar{ᶜ}{\ifmmode^{c}\else\textsuperscript{\ensuremath{c}}\fi}
\newunicodechar{ʰ}{\ifmmode^{h}\else\textsuperscript{\ensuremath{h}}\fi}
\newunicodechar{ᵠ}{\ifmmode^{q}\else\textsuperscript{\ensuremath{q}}\fi}
\newunicodechar{ₙ}{\ifmmode_{n}\else\textsubscript{\ensuremath{n}}\fi}
\newunicodechar{ₐ}{\ifmmode_{a}\else\textsubscript{\ensuremath{a}}\fi}
\newunicodechar{ᵢ}{\ifmmode_{i}\else\textsubscript{\ensuremath{i}}\fi}
\newunicodechar{ᵣ}{\ifmmode_{r}\else\textsubscript{\ensuremath{r}}\fi}
\newunicodechar{ₖ}{\ifmmode_{k}\else\textsubscript{\ensuremath{k}}\fi}
\newunicodechar{ᵦ}{\ifmmode_{\beta}\else\textsubscript{\ensuremath{\beta}}\fi}
\newunicodechar{ₓ}{\ifmmode_{x}\else\textsubscript{\ensuremath{x}}\fi}
\newfontfamily\popdisplay{ArchivoBlack-Regular.ttf}[Path=../../../fonts/]
\newfontfamily\poplogo{SpaceGrotesk-Bold.ttf}[Path=../../../fonts/]
\newfontfamily\popsemi{PlusJakartaSans-SemiBold.ttf}[Path=../../../fonts/]
\newfontfamily\popxbold{PlusJakartaSans-ExtraBold.ttf}[Path=../../../fonts/]
\newfontfamily\popxboldls{PlusJakartaSans-ExtraBold.ttf}[Path=../../../fonts/,LetterSpace=3.0]

\setlist[enumerate]{leftmargin=1.7em,itemsep=5pt,topsep=4pt}
\setlist[itemize]{leftmargin=1.7em,itemsep=4pt,topsep=4pt,label={\color{poporange}\textbullet}}
\setlength{\parindent}{0pt}
\setlength{\parskip}{5pt}
\renewcommand{\arraystretch}{1.25}

% pgfplots : style Pop par défaut
\pgfplotsset{
  every axis/.append style={axis line style={popink,line width=1pt},tick style={popink},
    grid style={popline,line width=0.5pt},label style={font=\small},tick label style={font=\footnotesize}},
  popcurve/.style={poporange,line width=1.8pt,no markers,smooth},
}
% Même style disponible dans un \draw TikZ simple (hors pgfplots)
\tikzset{popcurve/.style={poporange,line width=1.8pt}}
\tikzset{popgrid/.style={popline,very thin}}
\colorlet{popcurve}{poporange} % le nom du style est parfois utilisé comme couleur

% ============================================================
% Pill / squiggle
% ============================================================
\newcommand{\poppill}[3]{%
  \tikz[baseline=-0.55ex]\node[draw=popink,line width=1.2pt,rounded corners=2.6pt,%
    fill=#2,inner xsep=6.5pt,inner ysep=3pt]{%
    \popxboldls\fontsize{7}{7}\selectfont\color{#3}\MakeUppercase{#1}};%
}
\newcommand{\popsquiggle}[1]{%
  \tikz[baseline=-0.4ex]{
    \draw[#1,line width=2.6pt,line cap=round]
      plot [smooth] coordinates {(0,0.09) (0.34,0.01) (0.68,0.21) (1.02,0.01) (1.36,0.10)};
  }%
}
% Mot-clé surligné (marqueur orange sous le texte)
\newcommand{\cle}[1]{{\popxbold\color{popdark}#1}}

% ============================================================
% En-tête / pied
% ============================================================
\pagestyle{fancy}
\fancyhf{}
\renewcommand{\headrulewidth}{0pt}
\renewcommand{\footrulewidth}{0pt}
\lhead{\raisebox{-0.15ex}{\poplogo\fontsize{13}{13}\selectfont\color{popink}OnBuch\color{poporange}+}}
\rhead{\poppill{\DOCMATIERE\ \textbullet\ \DOCNIVEAU\ \textbullet\ Leçon \DOCLECON}{white}{popink}}
\lfoot{\popxbold\fontsize{7.6}{7.6}\selectfont\color{popmuted}\MakeUppercase{Cours signé OnBuch+}\ \textbullet\ {\popsemi\DOCTITRE}}
\rfoot{\tikz[baseline=-0.6ex]\node[rounded corners=8.5pt,fill=popink,minimum height=17pt,minimum width=17pt,inner xsep=5pt,inner sep=0]{\popxbold\fontsize{8}{8}\selectfont\color{popcream}\thepage\,/\,\pageref*{LastPage}};}

% ============================================================
% Titres
% ============================================================
\newcommand{\chaptitre}[2]{%
  \begin{tcolorbox}[enhanced,colback=poporange,colframe=popink,boxrule=2.6pt,arc=11pt,
      drop shadow=popink,left=15pt,right=15pt,top=12pt,bottom=11pt,before skip=3pt,after skip=18pt]
    \poppill{#2}{popink}{popcream}\par\vspace{9pt}
    {\raggedright\hyphenpenalty=10000\exhyphenpenalty=10000\popdisplay\fontsize{22}{24}\selectfont\color{popcream}\MakeUppercase{#1}\par}
  \end{tcolorbox}
}
% Section numérotée : pastille orange avec le numéro + titre Archivo
\newcounter{popsec}
\newcommand{\coursec}[1]{%
  \stepcounter{popsec}\setcounter{popsub}{0}%
  \par\vspace{16pt}\noindent
  \begin{minipage}{\linewidth}
  \tikz[baseline=-0.7ex]\node[draw=popink,line width=1.6pt,fill=poporange,rounded corners=4pt,minimum size=22pt,inner sep=0]{\popdisplay\fontsize{12}{12}\selectfont\color{popcream}\thepopsec};%
  \hspace{8pt}{\popdisplay\fontsize{15}{17}\selectfont\color{popink}\MakeUppercase{#1}}\par
  \vspace{4pt}\noindent{\color{poporange}\rule{36pt}{3.6pt}}
  \end{minipage}\par\nopagebreak\vspace{8pt}
}
\newcounter{popsub}
\newcommand{\courssub}[1]{%
  \stepcounter{popsub}%
  \par\vspace{9pt}\noindent{\popxbold\fontsize{11.5}{13}\selectfont\color{popink}{\color{poporange}\thepopsec.\thepopsub}\hspace{6pt}#1}\par\nopagebreak\vspace{4pt}
}

% ============================================================
% Boîtes pédagogiques — même recette : fond blanc, bordure épaisse colorée,
% ombre dure encre, badge plein. Titre optionnel [#1].
% ============================================================
\tcbset{popbase/.style={breakable,enhanced,colback=white,boxrule=2.3pt,arc=9pt,
  shadow={3pt}{-3pt}{0pt}{popink},left=14pt,right=14pt,top=11pt,bottom=11pt,before skip=11pt,after skip=15pt,
  fontupper=\popsemi\fontsize{10.6}{14.5}\selectfont\color{popink},
  fontlower=\popsemi\fontsize{10.6}{14.5}\selectfont\color{popink}}}
% Coche dessinée (le glyphe ✓ n'existe pas dans les polices du document)
\newcommand{\popcheck}[1]{\tikz[baseline=-0.5ex]\draw[#1,line width=1.6pt,line cap=round,line join=round](-0.13,0.0)--(-0.04,-0.09)--(0.13,0.1);}
\providecommand{\coche}{\popcheck{popgreen}}
\providecommand{\resultat}[1]{\par\smallskip\noindent\fcolorbox{popgreen}{popgreenL}{\parbox{\dimexpr\linewidth-2\fboxsep-2\fboxrule\relax}{\textbf{Résultat.} #1}}\par\smallskip}
\newcommand{\popboxhead}[3]{\poppill{#1}{#2}{white}\ifx\relax#3\relax\else\hspace{6pt}{\popxbold\fontsize{10.6}{12}\selectfont\color{popink}#3}\fi\par\vspace{7pt}}

\newtcolorbox{aretenir}[1][]{popbase,colframe=popgreen,before upper={\popboxhead{À retenir}{popgreen}{#1}}}
\newtcolorbox{exemplebox}[1][]{popbase,colframe=poporange,before upper={\popboxhead{Exemple}{poporange}{#1}}}
\newtcolorbox{definition}[1][]{popbase,colframe=popblue,before upper={\popboxhead{Définition}{popblue}{#1}}}
\newtcolorbox{propriete}[1][]{popbase,colframe=poppurple,before upper={\popboxhead{Propriété}{poppurple}{#1}}}
\newtcolorbox{methode}[1][]{popbase,colframe=popgold,before upper={\popboxhead{Méthode}{popgold}{#1}}}
\newtcolorbox{attention}[1][]{popbase,colframe=poppink,before upper={\popboxhead{Attention}{poppink}{#1}}}
\newtcolorbox{experience}[1][]{popbase,colframe=popblue,colback=popblueL,before upper={\popboxhead{Expérience}{popblue}{#1}}}
% « Le savais-tu ? » : carte violette pleine (contexte, culture, Cameroun)
\newtcolorbox{savaistu}[1][]{popbase,colback=poppurple,colframe=popink,
  fontupper=\popsemi\fontsize{10.4}{14.2}\selectfont\color{white},
  before upper={\poppill{Le savais-tu\,?}{white}{poppurple}\ifx\relax#1\relax\else\hspace{6pt}{\popxbold\color{white}#1}\fi\par\vspace{7pt}}}
% Exercice résolu : énoncé en haut, \tcblower puis la solution sur fond crème
\newcounter{popexo}\newcounter{popexr}
\newtcolorbox{exoresolu}[1][]{popbase,colframe=poporange,colbacklower=popcream,
  segmentation style={popink,line width=1.2pt,dashed},
  before upper={\stepcounter{popexr}\popboxhead{Exercice résolu \thepopexr}{poporange}{#1}},
  before lower={\poppill{Solution}{popink}{popcream}\par\vspace{6pt}}}
% Exercice (non résolu en place) — la correction est dans la partie Corrigés
\newtcolorbox{exercice}[1][]{popbase,colframe=popink,
  before upper={\stepcounter{popexo}\popboxhead{Exercice \thepopexo}{popink}{#1}}}
% Corrigé
\newtcolorbox{corrige}[1][]{popbase,colframe=popgreen,colback=popgreenL,
  before upper={\popboxhead{Corrigé}{popgreen}{#1}}}
% Objectifs / prérequis / bilan
\newtcolorbox{objectifs}{popbase,colframe=popink,colback=poporangeL,
  before upper={\popboxhead{Objectifs de la leçon}{popink}{}}}
\newtcolorbox{prerequis}{popbase,colframe=popink,colback=popblueL,
  before upper={\popboxhead{Prérequis}{popblue}{}}}
\newtcolorbox{bilan}{popbase,colframe=popink,colback=white,boxrule=2.8pt,
  before upper={\popboxhead{Fiche bilan}{poporange}{L'essentiel à connaître pour l'examen}}}
% Figure encadrée avec légende
\newtcolorbox{popfigure}[1][]{enhanced,breakable=false,colback=white,colframe=popline,boxrule=1.4pt,arc=8pt,
  left=10pt,right=10pt,top=10pt,bottom=8pt,before skip=10pt,after skip=12pt,halign=center,
  lower separated=false,halign lower=center,
  fontlower=\popsemi\fontsize{9}{11.5}\selectfont\color{popmuted},
  #1}
\newcommand{\legende}[1]{\par\vspace{6pt}{\popsemi\fontsize{9}{11.5}\selectfont\color{popmuted}#1}}

% ============================================================
% Signature OnBuch+
% ============================================================
\newcommand{\poplogoinline}[1]{{\poplogo\fontsize{#1}{#1}\selectfont\color{popink}OnBuch\color{poporange}+}}
\newcommand{\signatureonbuch}{%
  \begin{tcolorbox}[enhanced,colback=popink,colframe=popink,boxrule=2.2pt,arc=10pt,
      drop shadow=poporange,left=14pt,right=14pt,top=10pt,bottom=10pt,before skip=0pt,after skip=0pt]
    \tikz[baseline=-0.7ex]\node[circle,fill=popgreen,draw=popcream,line width=1.3pt,minimum size=22pt,inner sep=0]{\popcheck{white}};%
    \hspace{9pt}\begin{minipage}[c]{0.62\linewidth}
      {\popxbold\fontsize{12}{13}\selectfont\color{popcream}Signé\ \textbullet\ {\poplogo OnBuch\color{poporange}+}}\par\vspace{2pt}
      {\raggedright\popsemi\fontsize{8}{10.5}\selectfont\color{popline}\MakeUppercase{Équipe pédagogique OnBuch+}\\\MakeUppercase{Conforme au programme officiel MINESEC}\par}
    \end{minipage}\hfill
    {\poplogo\fontsize{22}{22}\selectfont\color{popcream}OnBuch\color{poporange}+}
  \end{tcolorbox}%
}

% ============================================================
% Page de garde
% #1 titre · #2 matière · #3 niveau · #4 module · #5 n° leçon · #6 durée ·
% #7 description / objectifs (texte ou itemize)
% ============================================================
\newcommand{\pagegarde}[7]{%
  \thispagestyle{empty}
  \newgeometry{a4paper,margin=1.7cm,top=1.6cm,bottom=1.4cm}%
  \begin{tikzpicture}[remember picture,overlay]
    \fill[poporange!18] ([xshift=-3.2cm,yshift=-3.6cm]current page.north east) circle (4.6cm);
    \fill[poppurple!14] ([xshift=2.4cm,yshift=7.6cm]current page.south west) circle (3.8cm);
  \end{tikzpicture}%
  {\poplogo\fontsize{30}{30}\selectfont\color{popink}OnBuch\color{poporange}+}\hfill
  \raisebox{0.6ex}{\poppill{Cours signé \textbullet\ Premium}{popink}{popcream}}\par
  \vspace{1pt}\popsquiggle{poppurple}\par
  \vspace{8pt}
  \begin{tcolorbox}[enhanced,colback=poporange,colframe=popink,boxrule=2.8pt,arc=13pt,
      drop shadow=popink,left=18pt,right=18pt,top=12pt,bottom=13pt,after skip=10pt]
    \poppill{#2 \textbullet\ #3}{popink}{popcream}\hspace{5pt}\poppill{#4}{white}{popink}\par\vspace{8pt}
    {\popdisplay\fontsize{14}{14}\selectfont\color{popink}LEÇON #5}\par\vspace{4pt}
    {\raggedright\hyphenpenalty=10000\exhyphenpenalty=10000\popdisplay\fontsize{25}{27}\selectfont\color{popcream}\MakeUppercase{#1}\par}
    \vspace{8pt}
    {\popxbold\fontsize{9.5}{9.5}\selectfont\color{popink}\MakeUppercase{Durée conseillée : #6}}
  \end{tcolorbox}
  \begin{tcolorbox}[enhanced,colback=white,colframe=popink,boxrule=2.2pt,arc=10pt,
      drop shadow=popink,left=16pt,right=16pt,top=10pt,bottom=10pt,after skip=8pt]
    \poppill{Dans cette leçon}{poppurple}{white}\par\vspace{5pt}
    \popsemi\fontsize{9.3}{12.6}\selectfont\color{popink}#7
  \end{tcolorbox}
  \noindent\begin{minipage}[t]{0.487\linewidth}
    \begin{tcolorbox}[enhanced,colback=popgreen,colframe=popink,boxrule=2.2pt,arc=10pt,
        drop shadow=popink,left=11pt,right=11pt,top=10pt,bottom=10pt,height=3.1cm,valign=center]
      \tcbox[colback=white,colframe=popink,boxrule=1.3pt,arc=5pt,boxsep=0pt,nobeforeafter,
        left=3pt,right=3pt,top=3pt,bottom=3pt,valign=center]{\qrcode[height=1.9cm]{https://play.google.com/store/apps/details?id=com.luvvix.onbuch&hl=fr}}%
      \hspace{8pt}\begin{minipage}[c]{0.5\linewidth}
        {\popdisplay\fontsize{11}{12.5}\selectfont\color{white}\MakeUppercase{Télécharge OnBuch}}\par\vspace{4pt}
        {\raggedright\popsemi\fontsize{8.4}{11}\selectfont\color{white}Gratuit \textbullet\ Google Play\\Tuteur IA, annales, concours, résultats\par}
      \end{minipage}
    \end{tcolorbox}
  \end{minipage}\hfill
  \begin{minipage}[t]{0.487\linewidth}
    \begin{tcolorbox}[enhanced,colback=poppurple,colframe=popink,boxrule=2.2pt,arc=10pt,
        drop shadow=popink,left=11pt,right=11pt,top=10pt,bottom=10pt,height=3.1cm,valign=center]
      \tikz[baseline=-0.7ex]\node[circle,fill=white,draw=popink,line width=1.3pt,minimum size=1.6cm,inner sep=0]{\popdisplay\fontsize{24}{24}\selectfont\color{poppurple}+};%
      \hspace{8pt}\begin{minipage}[c]{0.6\linewidth}
        {\popdisplay\fontsize{11}{12.5}\selectfont\color{white}\MakeUppercase{Passe à OnBuch+}}\par\vspace{4pt}
        {\raggedright\popsemi\fontsize{8.4}{11}\selectfont\color{white}Tous les cours signés, Tuteur IA illimité, fiches PDF — onglet OnBuch+ de l'app\par}
      \end{minipage}
    \end{tcolorbox}
  \end{minipage}
  \vspace{8pt}
  \popcontenu
  \vfill
  \signatureonbuch
  \restoregeometry
  \newpage
}

% Rangée « Ce cours contient » (page de garde)
\newcommand{\poptile}[3]{% #1 couleur #2 gros texte #3 légende
  \begin{tcolorbox}[enhanced,colback=#1,colframe=popink,boxrule=2pt,arc=9pt,shadow={2.5pt}{-2.5pt}{0pt}{popink},
      left=6pt,right=6pt,top=9pt,bottom=9pt,halign=center,valign=center,height=1.75cm,width=\linewidth,nobeforeafter]
    {\popdisplay\fontsize{12.5}{13}\selectfont\color{white}\MakeUppercase{#2}}\par\vspace{3pt}
    {\centering\popsemi\fontsize{7.8}{9.5}\selectfont\color{white}#3\par}
  \end{tcolorbox}}
\newcommand{\popcontenu}{%
  \par\noindent{\popxboldls\fontsize{7.5}{7.5}\selectfont\color{popmuted}CE COURS SIGNÉ CONTIENT}\par\vspace{7pt}
  \noindent\begin{minipage}[t]{0.235\linewidth}\poptile{popblue}{Cours}{complet, illustré, pas à pas}\end{minipage}\hfill
  \begin{minipage}[t]{0.235\linewidth}\poptile{poporange}{Méthodes}{et exercices résolus}\end{minipage}\hfill
  \begin{minipage}[t]{0.235\linewidth}\poptile{poppink}{Exercices}{type examen + corrigés}\end{minipage}\hfill
  \begin{minipage}[t]{0.235\linewidth}\poptile{popgreen}{Fiche bilan}{l'essentiel à retenir}\end{minipage}\par}

% Sommaire
\newtcolorbox{sommaire}{enhanced,colback=white,colframe=popink,boxrule=2.2pt,arc=10pt,
  drop shadow=popink,left=18pt,right=18pt,top=13pt,bottom=13pt,before skip=6pt,after skip=20pt,
  before upper={\poppill{Sommaire}{poppurple}{white}\par\vspace{9pt}\popsemi\fontsize{10.6}{14.5}\selectfont\color{popink}}}

% Signature de fin de cours — poussée tout en bas de la dernière page
\newcommand{\findecours}{%
  \par\vfill\nopagebreak
  \begin{center}\popsquiggle{poporange}\end{center}\vspace{4pt}
  \signatureonbuch
}
\AtBeginDocument{\def\llbracket{\mathopen{[\mkern-3.5mu[}}\def\rrbracket{\mathclose{]\mkern-3.5mu]}}} % intervalles d'entiers (glyphe ⟦ absent de la police)
% Glyphes absents de Fira Math : redéfinis à partir de glyphes disponibles
\AtBeginDocument{%
  \def\setminus{\mathbin{\backslash}}\let\smallsetminus\setminus
  \def\vdots{\mathord{\vbox{\baselineskip3.2pt\lineskiplimit0pt\kern4pt\hbox{.}\hbox{.}\hbox{.}}}}%
  \def\ddots{\mathinner{\mkern1mu\raise6.5pt\hbox{.}\mkern2mu\raise3.6pt\hbox{.}\mkern2mu\raise0.7pt\hbox{.}\mkern1mu}}%
  \def\triangle{\mathord{\scalebox{0.9}{$\symup{\Delta}$}}}%
  \def\nabla{\mathord{\raisebox{\depth}{\scalebox{1}[-1]{$\symup{\Delta}$}}}}%
  \def\checkmark{\popcheck{popdark}}%
  \def\star{\mathbin{\ast}}\def\bigstar{\mathord{\ast}}%
  \def\diamond{\mathbin{\scalebox{0.6}{$\Diamond$}}}%
  \def\bigcup{\mathop{\mathchoice{\raisebox{-0.3em}{\scalebox{1.7}{$\cup$}}}{\scalebox{1.25}{$\cup$}}{\cup}{\cup}}\displaylimits}%
  \def\bigcap{\mathop{\mathchoice{\raisebox{-0.3em}{\scalebox{1.7}{$\cap$}}}{\scalebox{1.25}{$\cap$}}{\cap}{\cap}}\displaylimits}%
  \def\lesseqqgtr{\mathrel{\lessgtr}}\def\bigoplus{\mathop{\oplus}\displaylimits}%
}
\providecommand{\ssi}{si et seulement si} % abréviation fréquente
\AtBeginDocument{\providecommand{\wideparen}{\overparen}} % arc de cercle

%% FIGFIX-BEGIN v3 — correctifs globaux des figures (docs/onbuch-generation/tools/figfix)
\usepackage{adjustbox}\usepackage{etoolbox}
% toute figure TikZ/pgfplots plus large que la ligne est réduite à la largeur disponible
\BeforeBeginEnvironment{tikzpicture}{\begin{adjustbox}{max width=\linewidth}}
\AfterEndEnvironment{tikzpicture}{\end{adjustbox}}
% légendes pgfplots lisibles par-dessus les courbes
\pgfplotsset{every axis/.append style={legend style={fill=white,fill opacity=0.92,text opacity=1,draw=black!15,inner sep=3pt}}}
% étiquettes d'arêtes (syntaxe edge["..."]) avec fond blanc
\tikzset{every edge quotes/.append style={fill=white,inner sep=1.2pt}}
% bulles de cartes mentales : texte à la ligne dans la bulle, bulle un peu plus grande
\tikzset{every concept/.append style={text width=2.0cm,align=center,minimum size=2.3cm,inner sep=2pt}}
%% FIGFIX-END
\begin{document}

\pagegarde{Formuler les idées essentielles d’un texte à résumer}{Français}{Tle Tech}{Français – classe de Terminale technique}{N25}{10 séances (fiche de progression harmonisée)}{Cette leçon apprend à identifier, hiérarchiser et reformuler les idées essentielles d'un texte argumentatif en maîtrisant les techniques de réduction linguistique et les outils de cohésion, à travers l'étude de "Le Vieux Nègre et la Médaille" et la préparation d'exposés oraux.
\vspace{4pt}
\begin{multicols}{2}\small
\begin{itemize}
\item Analyser la structure argumentative d'un texte (thèse, arguments, exemples, conclusion).
\item Différencier les contenus manifestes (explicites) et latents (implicites) d'un discours.
\item Appliquer les techniques de réduction (suppression, généralisation, substitution) sans trahir le sens.
\item Maîtriser l'usage des conjonctions de coordination et de subordination pour assurer la cohérence du résumé.
\item Reformuler les idées essentielles en changeant la formulation (syntaxe, vocabulaire) tout en conservant l'intention de l'auteur.
\item Préparer et présenter un exposé oral structuré à partir d'un dossier documentaire.
\end{itemize}\end{multicols}}

\begin{sommaire}
\begin{multicols}{2}
\begin{enumerate}
\item 1. Décryptage du texte argumentatif \& Lecture méthodique de l'œuvre
\item 2. Techniques de contraction : Structure argumentative et réduction
\item 3. Outils de liaison : Conjonctions de coordination et de subordination
\item 4. Atelier pratique : Reformulation et synthèse des idées essentielles
\item 5. Expression orale : Préparation et présentation d'un exposé
\item 6. Évaluation intégrée : Compte-rendu, Remédiation et Projection Bac
\item Activité d'intégration
\item Méthodes et exercices résolus
\item Exercices
\item Corrigés des exercices
\item Fiche bilan
\end{enumerate}
\end{multicols}
\end{sommaire}

\begin{objectifs}
\begin{itemize}
\item Analyser la structure argumentative d'un texte (thèse, arguments, exemples, conclusion).
\item Différencier les contenus manifestes (explicites) et latents (implicites) d'un discours.
\item Appliquer les techniques de réduction (suppression, généralisation, substitution) sans trahir le sens.
\item Maîtriser l'usage des conjonctions de coordination et de subordination pour assurer la cohérence du résumé.
\item Reformuler les idées essentielles en changeant la formulation (syntaxe, vocabulaire) tout en conservant l'intention de l'auteur.
\item Préparer et présenter un exposé oral structuré à partir d'un dossier documentaire.
\item Rédiger une contraction de texte respectant le taux de réduction imposé (1/3 ou 1/4) pour le Baccalauréat.
\end{itemize}
\end{objectifs}

\begin{prerequis}
\begin{itemize}
\item Reconnaissance des types et formes de phrases (déclarative, interrogative, etc.) et des connecteurs logiques simples (mais, donc, parce que).
\item Notion de paragraphe et d'idée force / idée secondaire (programme de Première).
\item Lecture cursive de "Le Vieux Nègre et la Médaille" de Ferdinand Oyono (contexte colonial, personnages, intrigue principale).
\item Techniques de base du résumé appris en classe de Seconde (surlignage, plan détaillé).
\item Conjugaison des temps de l'indicatif et du subjonctif (nécessaire pour la subordination).
\end{itemize}
\end{prerequis}

\newpage

\coursec{Décryptage du texte argumentatif \& Lecture méthodique de l'œuvre}

\textbf{Situation d'accroche.} Dans un bureau de la Mairie de Douala, un jeune fonctionnaire doit présenter en 5 minutes une synthèse d'un rapport de 20 pages sur l'assainissement du quartier New-Bell pour le Maire. Il panique : comment distinguer l'essentiel de l'accessoire, garder la logique de l'auteur et utiliser les bons connecteurs pour que ce soit fluide ? C'est exactement le défi de la contraction de texte argumentatif : avant de réduire un texte, il faut d'abord \emph{comprendre} comment il est construit, quelle idée il défend et par quels moyens il convainc. Telle est la problématique de cette première séquence : \textbf{comment décrypter la structure argumentative d'un texte, en particulier d'un roman comme \emph{Le Vieux Nègre et la Médaille} de Ferdinand Oyono, pour en dégager les idées essentielles ?}

\courssub{Qu'est-ce qu'un texte argumentatif ?}

\textbf{Intuition d'abord.} Quand un ami essaie de te convaincre d'aller voir un match au stade Ahmadou Ahidjo plutôt que de rester dormir, il ne se contente pas de dire « viens ». Il donne des raisons (« c'est la finale », « les places sont à 500 FCFA »), il illustre (« tu te souviens de l'ambiance l'année dernière ? »), il réfute tes objections (« tu pourras dormir demain »). Sans le savoir, il produit un texte argumentatif oral. L'écrivain fait la même chose, avec plus d'art.

\begin{definition}[Le texte argumentatif]
Un \cle{texte argumentatif} est un texte dont la \cle{visée} est de persuader (toucher l'affectif) ou de convaincre (s'adresser à la raison) un destinataire d'adhérer à une idée, un point de vue ou une conduite. Sa structure type comporte quatre éléments :
\begin{itemize}
\item la \cle{thèse} : l'idée défendue (explicite quand elle est énoncée clairement, implicite quand elle est suggérée) ;
\item les \cle{arguments} : les raisons logiques, morales ou émotionnelles qui soutiennent la thèse ;
\item les \cle{exemples} : les faits, anecdotes, statistiques ou illustrations qui rendent les arguments concrets et vérifiables ;
\item la \cle{conclusion} : le point d'arrivée qui referme le raisonnement, rappelle la thèse et ouvre parfois sur une perspective.
\end{itemize}
\end{definition}

\begin{attention}[Ne pas confondre sujet et thèse]
Erreur fréquente au Bac : confondre le \textbf{sujet} d'un texte (ce dont il parle, le thème) et sa \textbf{thèse} (ce qu'il affirme sur ce sujet, le point de vue). Dans \emph{Le Vieux Nègre et la Médaille}, le sujet est « l'attente et la remise d'une médaille à un vieil Africain » ; la thèse est « la critique du système colonial et de l'illusion de l'assimilation ». Le sujet se formule en un groupe nominal ; la thèse exige une phrase complète avec un jugement.
\end{attention}

\courssub{Structure globale du \emph{Vieux Nègre et la Médaille} : une satire coloniale}

\textbf{Observation préalable.} Lisons le début et la fin du roman. Au début, Meka, le vieux nègre, apprend qu'il va recevoir une médaille pour avoir « donné » ses terres à la mission et ses deux fils à la guerre ; il attend la cérémonie avec fierté. À la fin, ivre, humilié, battu par les forces de l'ordre après une nuit d'errance, il rejette la médaille et rit de sa propre illusion. Le point de départ et le point d'arrivée se répondent : c'est une \cle{structure circulaire}, où le retour au point initial révèle que tout a changé dans le regard du personnage.

\begin{definition}[Ironie et satire]
L'\cle{ironie} est un procédé qui consiste à dire le contraire de ce qu'on pense, ou à laisser voir l'écart entre les apparences et la réalité, pour faire réfléchir ou rire (souvent jaune). La \cle{satire} est une dénonciation des vices, des abus ou des ridicules d'une société par le rire, l'exagération et la caricature. \emph{Le Vieux Nègre et la Médaille} est une \textbf{satire coloniale} : Oyono ridiculise l'administration coloniale et la mission religieuse pour dénoncer l'injustice du système.
\end{definition}

La structure circulaire du roman est elle-même un argument : en ramenant Meka à son point de départ dépouillé de ses illusions, Oyono montre que la médaille — symbole de la « reconnaissance » coloniale — n'a rien changé à la condition du vieil homme. L'honneur promis était un piège.

\courssub{Lecture méthodique n°4 : repérage de la thèse implicite et des arguments narratifs}

\textbf{Problème.} Un roman ne dit jamais « voici ma thèse ». Comment la dégager ? Il faut observer ce que le récit \emph{montre} : les faits subis par les personnages, les contrastes, les répétitions.

\begin{methode}[Dégager la thèse implicite d'un récit]
\begin{enumerate}
\item Relever les faits marquants subis par le personnage principal (ce qu'il donne, ce qu'il perd, ce qu'il reçoit en retour).
\item Repérer les contrastes entre les discours officiels (promesses, idéologie) et la réalité vécue (souffrance, spoliation).
\item Observer le dénouement : que devient le personnage ? Quelle leçon le lecteur tire-t-il de la chute ?
\item Formuler la thèse en une phrase complète : « L'auteur montre que... » ou « L'auteur dénonce... ».
\end{enumerate}
\end{methode}

\textbf{Application.} Meka a donné ses terres à la mission catholique, ses deux fils sont morts « pour la France », et il reçoit en échange... une médaille en métal et une gifle de la réalité. Le contraste est flagrant : le discours colonial parle de « reconnaissance » et de « civilisation », la réalité montre la spoliation et le mépris. Le dénouement — Meka battu, enfermé, puis lucide — confirme la leçon.

\begin{aretenir}[Thèse implicite du roman]
Dans \emph{Le Vieux Nègre et la Médaille} (1956), Ferdinand Oyono défend la thèse suivante : \textbf{l'assimilation coloniale est une illusion humiliante} ; les « bienfaits » de la colonisation (médailles, discours, religion) masquent la spoliation et le mépris dont sont victimes les Africains.
\end{aretenir}

Les \textbf{arguments narratifs} sont les épisodes du roman qui prouvent cette thèse sans la formuler explicitement :
\begin{enumerate}
\item L'attente naïve et fière de Meka, qui croit à la consécration (l'illusion entretenue).
\item La cérémonie du 14 juillet, où les discours officiels sonnent creux et où l'interprète trahit la pensée du chef de poste (la mascarade institutionnelle).
\item La nuit d'errance, l'arrestation et les coups, qui révèlent le vrai visage du système répressif (la violence réelle).
\item Le rejet final de la médaille, signe de la lucidité retrouvée (la prise de conscience).
\end{enumerate}

\courssub{Lecture méthodique n°5 : personnages et symboles}

\textbf{Approfondissons} en interrogeant les figures du roman, car ce sont elles qui portent l'argumentation narrative.

\textbf{Meka.} Vieil homme du village de Doum, il incarne la génération africaine qui a cru aux promesses coloniales. Sa naïveté initiale (il se vante auprès des villageois, se fait confectionner des chaussures neuves qui le blesseront) puis sa lucidité finale (« plus jamais ça ») tracent un parcours d'apprentissage douloureux. Il est à la fois ridicule (par sa soumission) et pathétique (par sa souffrance) : l'ironie d'Oyono le moque, mais la tendresse de l'auteur le sauve et force le respect du lecteur.

\textbf{La médaille.} Elle symbolise la fausse reconnaissance coloniale : un objet brillant, sans valeur réelle, offert en échange de sacrifices immenses (terres, fils, vie). C'est le symbole de l'\emph{illusion} de l'assimilation : on décore l'indigène pour mieux l'asservir.

\textbf{Le prêtre (le père Vandermayer) et l'administrateur (le commandant).} Le prêtre représente la mission religieuse, qui a recueilli les terres de Meka ; l'administrateur représente le pouvoir colonial, qui décore puis fait brutaliser le même homme. Leur complicité lors de la cérémonie montre l'alliance de la croix et de l'épée au service du même système. Leurs discours pompeux, traduits maladroitement par l'interprète, tournent à la farce : l'ironie d'Oyono les démasque.

\courssub{Contenu manifeste et contenu latent}

\begin{definition}[Contenu manifeste et contenu latent]
Le \cle{contenu manifeste} d'un texte est ce qui est dit explicitement : l'intrigue, les dialogues, les actions, les descriptions, le décor. Le \cle{contenu latent} est ce qui est suggéré, sous-entendu, dénoncé indirectement : les valeurs, les critiques, les émotions, la thèse que le lecteur doit reconstruire par inférence. Un bon lecteur (et un bon contracteur de texte) lit le manifeste mais comprend le latent pour ne garder que l'essentiel.
\end{definition}

\begin{center}
\begin{tabularx}{\linewidth}{|l|X|X|}
\hline
\rowcolor{popblueL} \textbf{Passage} & \textbf{Contenu manifeste (dit explicitement)} & \textbf{Contenu latent (suggéré, à inférer)} \\
\hline
La remise de la médaille & Le commandant épingle la médaille, prononce un discours sur la fidélité, la foule applaudit. & La cérémonie est une mascarade : les discours sont vides, la « reconnaissance » est un théâtre pour masquer l'exploitation. \\
\hline
Les chaussures neuves de Meka & Meka porte des chaussures qui le blessent ; il marche avec peine mais ne les quitte pas. & L'accoutrement à l'européenne est une souffrance physique et symbolique : l'assimilation blesse celui qui s'y soumet. \\
\hline
La nuit d'arrestation & Meka, ivre, est battu et enfermé par les gardes après avoir troublé l'ordre. & Le colonisé décoré le matin redevient un « indigène » frappé le soir : l'honneur colonial est un mensonge, la violence est la vraie loi. \\
\hline
\end{tabularx}
\end{center}

\courssub{Les opérateurs argumentatifs dans le texte}

Pour suivre la logique d'un texte, il faut repérer ses \cle{opérateurs argumentatifs} : les \textbf{connecteurs logiques} (mais, donc, cependant, enfin, ainsi, par conséquent), les \textbf{modalisateurs} (mots qui marquent le degré de certitude ou l'attitude de l'énonciateur : certes, sans doute, il est clair que, évidemment) et les \textbf{temps verbaux} (le présent de vérité générale pour les affirmations intemporelles, le passé composé et l'imparfait pour le récit des faits arguments).

\begin{exemplebox}[Exemple : la cérémonie de remise]
Dans l'extrait de la cérémonie (édition de référence CLE, p. 80-90 environ), le discours du chef de poste utilise \textbf{5 connecteurs logiques} qui structurent sa rhétorique : \emph{donc} (conséquence), \emph{mais} (opposition), \emph{cependant} (concession), \emph{ainsi} (conséquence/synthèse), \emph{enfin} (clôture). Ces connecteurs donnent au discours une apparence de rigueur logique — que le contenu vide et hypocrite dément : c'est précisément cet écart entre la forme (rigoureuse) et le fond (mensonger) qui crée l'ironie critique d'Oyono.
\end{exemplebox}

\courssub{Le schéma argumentatif de l'œuvre}

\begin{methode}[Construire le schéma argumentatif d'un texte littéraire]
\begin{enumerate}
\item Formuler la thèse (souvent latente) en une phrase complète et autonome.
\item Dégager 2 à 4 arguments narratifs majeurs (épisodes ou étapes du récit qui prouvent la thèse).
\item Associer à chaque argument un exemple concret précis tiré du texte (scène, réplique, détail signifiant).
\item Identifier la conclusion ou la chute (souvent ironique chez Oyono) et la relier à la thèse.
\item Dessiner le schéma : thèse en haut, arguments au centre, conclusion en bas, reliés par des flèches logiques.
\end{enumerate}
\end{methode}

\begin{popfigure}
\begin{tikzpicture}[node distance=1.2cm and 0.5cm, every node/.style={font=\small, align=center, draw, rectangle, rounded corners, inner sep=4pt}]
\node[fill=popblueL, text width=8cm] (thesis) {Thèse latente : l'assimilation est une illusion humiliante};
\node[fill=popgreenL, text width=5cm, below=of thesis] (arg1) {Arg 1 : Meka attend la médaille comme une consécration};
\node[fill=popgreenL, text width=4.5cm, left=of arg1] (arg2) {Arg 2 : la cérémonie tourne à la farce (discours vides)};
\node[fill=popgreenL, text width=4.5cm, right=of arg1] (arg3) {Arg 3 : le rejet final de la médaille par Meka};
\node[fill=poppinkL, text width=7cm, below=of arg1] (concl) {Conclusion : la lucidité douloureuse du « vieux nègre »};
\draw[->, thick] (thesis) -- (arg1);
\draw[->, thick] (thesis) -- (arg2);
\draw[->, thick] (thesis) -- (arg3);
\draw[->, thick] (arg1) -- (concl);
\draw[->, thick] (arg2) -- (concl);
\draw[->, thick] (arg3) -- (concl);
\end{tikzpicture}
\legende{Schéma argumentatif du \emph{Vieux Nègre et la Médaille} : la thèse latente (en haut) est prouvée par trois arguments narratifs (au centre), qui convergent vers la chute ironique (en bas).}
\end{popfigure}

\begin{exoresolu}[Dégager thèse et arguments d'un extrait]
\textbf{Énoncé.} Voici un extrait adapté du roman : « Meka regardait sa médaille. Elle brillait au soleil. Le commandant avait parlé de reconnaissance, de fidélité, de la France éternelle. Mais le soir même, les gardes le frappaient comme un chien. » (a) Formulez la thèse implicite. (b) Relevez deux arguments narratifs. (c) Identifiez un connecteur et précisez sa fonction argumentative.
\tcblower
\textbf{Solution.} (a) Thèse : les honneurs coloniaux sont un mensonge qui masque le mépris réel et la violence du système envers les Africains. (b) Argument 1 : le discours officiel exalte la « reconnaissance » et la « fidélité » (paroles officielles). Argument 2 : les gardes battent Meka le soir même (actes réels) ; le contraste paroles/actes prouve l'hypocrisie du système. (c) Le connecteur \emph{mais} introduit une opposition forte (adversative) : il fait basculer le passage de l'apparence (cérémonie, lumière) à la réalité (violence, nuit), et c'est ce basculement brutal qui porte toute la charge critique du texte.
\end{exoresolu}

\begin{savaistu}[Oyono, romancier et diplomate]
Ferdinand Oyono (1929-2010), né à Ngoulémakong au Cameroun, n'a pas été seulement écrivain : il a mené une longue carrière de diplomate et d'homme d'État (ambassadeur en plusieurs pays, ministre des Affaires étrangères). Publié en 1956, \emph{Le Vieux Nègre et la Médaille} a préfiguré, par sa critique lucide de la colonisation, l'indépendance du Cameroun proclamée le 1er janvier 1960. La littérature africaine de cette époque fut une véritable arme de combat politique et de reconquête de la dignité.
\end{savaistu}

\begin{exercice}[À toi de jouer]
Relis la scène où Meka essaie ses chaussures neuves avant la cérémonie (début du roman). (1) Décris le contenu manifeste de la scène en deux phrases. (2) Dégage son contenu latent en une phrase. (3) Intègre cette scène dans le schéma argumentatif de l'œuvre : à quel argument sert-elle d'exemple et pourquoi ?
\end{exercice}

\begin{corrige}[Exercice — éléments de réponse]
(1) Contenu manifeste : Meka enfile des chaussures neuves, trop étroites, qui lui font mal aux pieds ; il marche avec difficulté mais persiste à les garder par fierté. (2) Contenu latent : l'adoption des habitudes et de l'apparence européennes est une souffrance imposée ; l'assimilation blesse littéralement celui qui s'y plie sans en avoir les codes culturels. (3) Cette scène illustre l'argument 1 (l'attente de la consécration / l'illusion entretenue) : Meka accepte la douleur physique parce qu'il croit encore à l'honneur promis par la médaille ; elle prépare ironiquement la désillusion finale où il comprendra qu'il a souffert pour rien.
\end{corrige}

\begin{aretenir}[L'essentiel de la séquence]
Décrypter un texte argumentatif, c'est : distinguer le \cle{sujet} de la \cle{thèse} ; repérer les \cle{arguments} et les \cle{exemples} ; lire au-delà du \cle{contenu manifeste} pour atteindre le \cle{contenu latent} ; suivre les \cle{opérateurs argumentatifs} (connecteurs, modalisateurs, temps) ; et construire le \cle{schéma argumentatif} du texte. Chez Oyono, la structure circulaire et l'ironie font du roman entier une démonstration implicite : l'assimilation coloniale est une illusion humiliante. Cette compétence de lecture analytique est la condition \emph{sine qua non} de la contraction de texte, que nous aborderons dans la prochaine séquence.
\end{aretenir}

\coursec{Techniques de contraction : Structure argumentative et réduction}

\courssub{Transition depuis la section précédente}
Dans la section précédente, nous avons appris à \emph{décrypter} un texte argumentatif : identifier la thèse, repérer les arguments, analyser les connecteurs logiques et distinguer les contenus manifestes des contenus latents lors de la lecture méthodique du \emph{Vieux nègre et la médaille}. Cette analyse fine est le \textbf{prérequis indispensable} à toute contraction réussie. On ne peut réduire ce que l'on n'a pas compris en profondeur. Nous passons maintenant de la \emph{réception} (lecture/analyse) à la \emph{production} (écriture/contraction). L'objectif : produire un texte court (le tiers de l'original pour les séries techniques) qui soit le \emph{miroir fidèle} de la pensée de l'auteur, sans en être la paraphrase servile.

\courssub{Le cadre officiel : Consignes du Baccalauréat Technique}

\begin{definition}[Contraction de texte (Résumé contraint)]
La \cle{contraction de texte} est un exercice de réécriture consistant à réduire un texte source à une fraction de sa longueur (généralement \textbf{un tiers} pour les séries techniques au Cameroun) tout en conservant \textbf{intégralement les idées essentielles}, l'\textbf{enchaînement logique} et la \textbf{tonalité} (modalités d'assertion, nuances) de l'auteur. Elle interdit tout ajout personnel, toute interprétation ou tout commentaire.
\end{definition}

\begin{definition}[Taux de réduction]
Le \cle{taux de réduction} est le rapport entre le nombre de mots du texte cible (la contraction) et le nombre de mots du texte source. Pour le Probatoire/Baccalauréat technique : 
\[ \text{Taux cible} = \frac{1}{3} \quad (\text{soit environ } 33,3\% ) \]
Une tolérance de \textbf{$\pm 10\%$} sur le nombre de mots final est généralement admise (ex: pour un source de 300 mots, cible = 100 mots ; intervalle acceptable : 90 à 110 mots).
\end{definition}

\begin{aretenir}[Les 4 piliers de la contraction au Bac]
\begin{enumerate}
    \item \textbf{Fidélité au sens} : Aucune idée ajoutée, aucune idée majeure omise, respect des nuances (probabilité, certitude, ironie).
    \item \textbf{Neutralité} : Disparition du \emph{je} de l'élève. On utilise l'impersonnel (\emph{Il est question de...}, \emph{L'auteur montre que...}) ou le discours indirect libre.
    \item \textbf{Cohérence} : Le texte réduit doit être lisible comme un texte autonome, fluide, correct syntaxiquement.
    \item \textbf{Respect du taux} : Comptage rigoureux (mots outils + mots lexicaux).
\end{enumerate}
\end{aretenir}

\courssub{Les trois opérations chirurgicales de la réduction}
Réduire un texte, ce n'est pas « couper au hasard », c'est opérer trois gestes techniques précis, souvent combinés.

\begin{definition}[Suppression (Élagage)]
La \cle{suppression} consiste à éliminer les segments redondants ou secondaires par rapport à l'idée force : 
\begin{itemize}
    \item Les \textbf{exemples} et illustrations (sauf s'ils \emph{sont} l'argument unique).
    \item Les \textbf{détails} descriptifs, anecdotiques, temporels/spatiaux non essentiels à la démonstration.
    \item Les \textbf{répétitions}, reformulations, amplifications rhétoriques.
    \item Les \textbf{incises} et parenthèses commentées.
\end{itemize}
\end{definition}

\begin{definition}[Généralisation (Montée en hyperonymie)]
La \cle{généralisation} remplace une énumération, une série de termes précis ou une explication développée par un \textbf{hyperonyme} (terme plus général) ou une proposition synthétique.
\begin{center}
    \textbf{Exemple :} \textit{« Le chômage, la pauvreté, la corruption, l'insécurité minent la société »} $\rightarrow$ \textit{« Les fléaux sociaux minent la société »}.
\end{center}
\end{definition}

\begin{definition}[Substitution (Condensation syntaxique)]
La \cle{substitution} remplace une structure lourde (phrase, proposition) par une structure plus compacte (nominalisation, participe, infinitif, adjectif, pronom relatif) sans perte de sens.
\begin{center}
    \begin{tabular}{|l|l|l|}
    \hline
    \textbf{Source (lourd)} & \textbf{Opération} & \textbf{Cible (compact)} \\
    \hline
    \textit{« Meka attend l'administrateur qui doit lui remettre la médaille »} & Nominalisation / Relatif & \textit{« Meka attend l'administrateur pour la remise de la médaille »} \\
    \textit{« Bien qu'il soit vieux, il court »} & Concession (participe) & \textit{« Malgré son âge, il court »} \\
    \textit{« Il a dit : "Je viendrai" »} & Discours indirect & \textit{« Il a annoncé sa venue »} \\
    \hline
    \end{tabular}
\end{center}
\end{definition}

\begin{popfigure}
\begin{tikzpicture}[scale=0.85, every node/.style={font=\small}]
    % Axe principal
    \draw[<->, thick, popdark] (0,0) -- (13,0);
    % Graduations et labels
    \foreach \x/\label in {
        0/Texte Source (300 mots),
        3/Opération 1: Suppression (détails{,} exemples{,} répétitions),
        6/Opération 2: Généralisation (énumérations $\to$ hyperonymes),
        9/Opération 3: Substitution (phrases $\to$ noms/participes),
        13/Texte Cible (100 mots)} {
        \draw[thick, popdark] (\x,0.15) -- (\x,-0.15) node[below=2mm, align=center, text width=2.5cm] {\label};
    }
    % Boîte de contrôle flottante
    \node[draw=popblue, thick, rectangle, rounded corners, fill=popblueL, align=center, above=1.2cm, text width=8cm] at (6.5,0) {
        \textbf{Contrôle qualité continu} \\
        Taux = 1/3 ? \quad Sens conservé ? \quad Connecteurs logiques préservés ? \quad Neutralité respectée ?
    };
    % Flèches de progression
    \draw[poporange, thick, -{Stealth[length=3mm]}] (1.5,0.5) -- (2.5,0.5);
    \draw[poporange, thick, -{Stealth[length=3mm]}] (4.5,0.5) -- (5.5,0.5);
    \draw[poporange, thick, -{Stealth[length=3mm]}] (7.5,0.5) -- (8.5,0.5);
    \draw[popgreen, thick, -{Stealth[length=3mm]}] (10.5,0.5) -- (11.5,0.5);
\end{tikzpicture}
\legende{Le processus de contraction vu comme une chaîne opératoire séquentielle : chaque opération réduit le volume textuel tout devant préserver le noyau sémantique.}
\end{popfigure}

\courssub{Application sur la structure argumentative : Noyau et Satellites}
Un texte argumentatif suit une architecture : \textbf{Thèse} $\rightarrow$ \textbf{Arguments} (souvent organisés : fort / moyen / fort) $\rightarrow$ \textbf{Conclusion}. La contraction doit respecter cette \textbf{architecture}. On ne contracte pas « au fil de l'eau » mot à mot, mais \emph{idée par idée}.

\begin{methode}[Identifier l'idée force et ses satellites (Étape cruciale avant rédaction)]
\begin{enumerate}
    \item \textbf{Surligner la phrase nucléaire (Thèse) :} Souvent en début ou fin de paragraphe, parfois implicite.
    \item \textbf{Numéroter les arguments :} Repérer les connecteurs (\emph{Premièrement, En outre, Paradoxalement, En définitive}).
    \item \textbf{Qualifier les satellites :} Pour chaque argument, distinguer : 
    \begin{itemize}
        \item L'\textbf{idée directrice} (à garder impérativement).
        \item La \textbf{preuve/exemple/développement} (à généraliser ou supprimer).
        \item Le \textbf{connecteur logique} (à garder ou adapter pour la fluidité).
    \end{itemize}
    \item \textbf{Construire le « squelette » :} Une liste numérotée des idées directrices dans l'ordre de l'auteur. C'est le plan de la contraction.
\end{enumerate}
\end{methode}

\begin{propriete}[Principe de fidélité structurelle]
La contraction ne doit contenir \textbf{aucune idée absente du source}, ni \textbf{modifier la portée des assertions} (modalité). Si l'auteur écrit \emph{« Il est probable que... »}, la contraction ne doit pas écrire \emph{« Il est certain que... »}. L'ordre des arguments de l'auteur est \textbf{inviolable} : ne jamais réorganiser la démonstration pour la « rendre plus logique ».
\end{propriete}

\courssub{Gestion des citations et du discours rapporté}
En contraction, la \textbf{citation directe} (guillemets, deux-points) est \textbf{proscrite} (trop gourmande en mots, rompt l'homogénéité). On pratique l'\textbf{intégration indirecte}.

\begin{definition}[Intégration indirecte]
Transformation du discours direct/indirect en proposition subordonnée (complétive ou relative) ou en groupe nominal.
\begin{itemize}
    \item \textbf{Que + Indicatif} : pour les faits, certitudes (\emph{L'auteur affirme \textbf{que} le chômage augmente}).
    \item \textbf{Que + Subjonctif} : pour la volonté, le doute, la nécessité, l'émotion (\emph{L'auteur exige \textbf{que} des mesures soient prises}).
    \item \textbf{De + Infinitif / Nominalisation} : plus compact (\emph{L'auteur préconise \textbf{la création d'emplois}}).
\end{itemize}
\end{definition}

\begin{attention}[Pièges fréquents sur les citations]
\begin{itemize}
    \item \textbf{Conserver les guillemets :} \textit{« L'auteur dit : "Il faut agir" »} (9 mots) $\rightarrow$ \textit{L'auteur préconise l'action} (4 mots).
    \item \textbf{Changer la modalité :} Source : \textit{« Peut-être la solution viendra-t-elle des jeunes »}. Contraction fautive : \textit{« La solution viendra des jeunes »} (perte de l'incertitude). Contraction correcte : \textit{« L'auteur envisage une solution venue des jeunes »}.
    \item \textbf{Attribuer à tort :} Ne pas écrire \emph{« L'auteur pense que... »} si le texte rapporte l'avis d'un tiers (\emph{« Selon X, ... »} $\rightarrow$ \emph{« X estime que... »}).
\end{itemize}
\end{attention}

\courssub{Méthodologie complète : Les 4 étapes de la contraction}

\begin{methode}[Opérer une contraction en 4 étapes (Méthode Bac)]
\begin{enumerate}
    \item \textbf{Lecture active \& Surlignage (15-20\% du temps) :} 
    \begin{itemize}
        \item 1ère lecture : compréhension globale (sujet, thèse, ton).
        \item 2ème lecture : surligner \emph{uniquement} les mots-clés des idées forces (thèse, arguments principaux, conclusion). Ignorer les exemples.
        \item Compter les mots du source ($N_{source}$). Calculer la cible ($N_{cible} = N_{source} / 3$) et l'intervalle $\pm 10\%$.
    \end{itemize}
    \item \textbf{Plan détaillé des idées essentielles (20\% du temps) :} 
    \begin{itemize}
        \item Rédiger au brouillon la liste séquentielle des idées retenues (une ligne par idée force).
        \item Vérifier l'exhaustivité : ai-je la thèse ? Tous les arguments majeurs ? La conclusion ?
        \item Vérifier l'ordre : est-il identique au source ?
    \end{itemize}
    \item \textbf{Rédaction du brouillon : Application S/G/S (40-50\% du temps) :} 
    \begin{itemize}
        \item Écrire des phrases complètes en enchaînant les idées du plan.
        \item Appliquer systématiquement : \textbf{S}uppression (coupure), \textbf{G}énéralisation (hyperonyme), \textbf{S}ubstitution (nominalisation, discours indirect, participes).
        \item Soigner les \textbf{connecteurs logiques} (donc, cependant, par conséquent, en somme) pour assurer la cohérence du texte réduit.
    \end{itemize}
    \item \textbf{Comptage \& Peaufinage linguistique (15-20\% du temps) :} 
    \begin{itemize}
        \item Compter \textbf{chaque mot} (voir méthode officielle ci-dessous).
        \item Si $> N_{max}$ : resserrer (plus de nominalisation, couper adjectifs qualificatifs restants, fusionner phrases).
        \item Si $< N_{min}$ : développer légèrement (rétablir un connecteur, préciser un agent).
        \item Relire à voix haute : fluidité, orthographe, accord, ponctuation.
        \item Recopier proprement sur la copie d'examen.
    \end{itemize}
\end{enumerate}
\end{methode}

\courssub{Comptage officiel des mots : La méthode rigoureuse}
Au Baccalauréat, le comptage est strict. La règle officielle (MINESEC) :
\begin{itemize}
    \item \textbf{Mots lexicaux} (noms, verbes, adjectifs, adverbes) = 1 mot.
    \item \textbf{Mots outils} (déterminants, prépositions, conjonctions, pronoms, auxiliaires) = 1 mot.
    \item \textbf{Mots composés} (trait d'union) = 1 mot (\emph{aujourd'hui, porte-parole, c'est-à-dire}).
    \item \textbf{Nombres écrits en chiffres} (\emph{2025}) = 1 mot. Écrits en lettres (\emph{deux mille vingt-cinq}) = 4 mots.
    \item \textbf{Signes de ponctuation} = 0 mot.
    \item \textbf{Apostrophe} : \emph{L'homme} = 2 mots (\emph{le} + \emph{homme}) ; \emph{Qu'il} = 2 mots (\emph{que} + \emph{il}).
\end{itemize}

\begin{exemplebox}[Comptage pratique]
\textit{« Meka, le vieux chef, attend avec une impatience fébrile l'arrivée de l'administrateur qui doit lui remettre la médaille de la Légion d'honneur. »}
\begin{enumerate}
    \item Meka (1)
    \item le (2)
    \item vieux (3)
    \item chef (4)
    \item attend (5)
    \item avec (6)
    \item une (7)
    \item impatience (8)
    \item fébrile (9)
    \item l'arrivée (10) [l' + arrivée]
    \item de (11)
    \item l'administrateur (12) [l' + administrateur]
    \item qui (13)
    \item doit (14)
    \item lui (15)
    \item remettre (16)
    \item la (17)
    \item médaille (18)
    \item de (19)
    \item la (20)
    \item Légion (21)
    \item d'honneur (22) [d' + honneur]
\end{enumerate}
\textbf{Total : 22 mots.}
\end{exemplebox}

\courssub{Exercice résolu progressif : Éditorial « Cameroon Tribune »}

\begin{exoresolu}[Contraction d'un éditorial sur l'emploi des jeunes]
\textbf{Énoncé :} Contractez le texte suivant (152 mots) en 50 mots $\pm 10\%$ (soit 45 à 55 mots). Respectez la neutralité et l'ordre des idées.

\textbf{Texte source (152 mots) :}
\emph{« L'emploi des jeunes constitue aujourd'hui le défi majeur de notre nation. En effet, chaque année, des milliers de diplômés sortent des universités et des grandes écoles sans trouver de débouchés professionnels. Cette situation engendre une frustration sociale palpable, source d'instabilité politique potentielle. Le gouvernement a bien lancé plusieurs programmes d'insertion, comme le PIAASI ou le Programme Spécial Jeunes, mais leur impact reste limité par la lourdeur administrative et le manque de financement structurel. Par ailleurs, le secteur privé, moteur naturel de la croissance, peine à absorber cette main-d'œuvre qualifiée à cause d'un environnement des affaires encore perfectible. Il est donc impératif de repenser en profondeur notre modèle éducatif pour l'aligner sur les besoins réels du marché du travail, tout en assainissant le climat des affaires pour encourager l'investissement privé. L'avenir du pays en dépend. »}

\textbf{Étape 1 : Analyse \& Comptage Source}
$N_{source} = 152$ mots. Cible $= 50,6 \approx 51$ mots. Intervalle : \textbf{46 à 56 mots}.
Idées forces (surlignage mental) :
1. Emploi jeunes = défi majeur.
2. Diplômés nombreux sans débouchés.
3. Frustration $\rightarrow$ instabilité.
4. Programmes gov (PIAASI, etc.) existent mais limités (admin, financement).
5. Secteur privé peine à absorber (environnement affaires).
6. Solutions : Réformer éducation (adéquation marché) + Assainir climat affaires.
7. Enjeu : Avenir pays.

\textbf{Étape 2 : Plan détaillé (Squelette)}
1. L'emploi des jeunes est le défi majeur national.
2. L'afflux de diplômés sans emploi crée frustration et risque d'instabilité.
3. Les programmes publics d'insertion sont inefficaces (lourdeurs, sous-financement).
4. Le secteur privé n'absorbe pas cette main-d'œuvre (climat des affaires défavorable).
5. Il faut réformer l'éducation (adéquation marché) et améliorer le climat des affaires.
6. L'avenir national est en jeu.

\textbf{Étape 3 : Rédaction brouillon (Application S/G/S)}
\textit{Brouillon 1 :} L'emploi des jeunes est le défi majeur. Des milliers de diplômés sans emploi créent une frustration risquant l'instabilité. Les programmes gouvernementaux d'insertion échouent par lourdeur administrative et manque de fonds. Le secteur privé, freiné par un mauvais climat des affaires, n'absorbe pas ces qualifications. Il faut urgemment réformer l'éducation vers les besoins du marché et assainir le climat des affaires pour attirer l'investissement. L'avenir du pays en dépend.
\textit{Comptage brouillon 1 :} 58 mots (Hors tolérance haute 56). Il faut couper encore.

\textit{Optimisation (Substitution/Généralisation) :}
- « Des milliers de diplômés sans emploi » $\rightarrow$ « Le chômage des diplômés » (G).
- « créent une frustration risquant l'instabilité » $\rightarrow$ « génère frustration et instabilité » (S).
- « échouent par lourdeur administrative et manque de fonds » $\rightarrow$ « échouent (lourdeurs, sous-financement) » (S).
- « freiné par un mauvais climat des affaires » $\rightarrow$ « freiné par le climat des affaires » (S).
- « vers les besoins du marché » $\rightarrow$ « aux besoins du marché » (S).
- Supprimer « L'avenir du pays en dépend » (implicite dans « impératif » ou garder si place). Gardons-le pour la conclusion.

\textit{Brouillon 2 (Final) :}
\emph{L'emploi des jeunes est le défi majeur national. Le chômage des diplômés génère frustration et instabilité. Les programmes publics d'insertion échouent par lourdeurs et sous-financement. Le secteur privé, freiné par le climat des affaires, n'absorbe pas cette main-d'œuvre. Il faut réformer l'éducation aux besoins du marché et assainir le climat des affaires pour attirer l'investissement. L'avenir national en dépend.}

\textbf{Étape 4 : Comptage final \& Vérification}
1. L'emploi (1) 2. des (2) 3. jeunes (3) 4. est (4) 5. le (5) 6. défi (6) 7. majeur (7) 8. national (8). 9. Le (9) 10. chômage (10) 11. des (11) 12. diplômés (12) 13. génère (13) 14. frustration (14) 15. et (15) 16. instabilité (16). 17. Les (17) 18. programmes (18) 19. publics (19) 20. d'insertion (20) 21. échouent (21) 22. par (22) 23. lourdeurs (23) 24. et (24) 25. sous-financement (25). 26. Le (26) 27. secteur (27) 28. privé (28) 29. freiné (29) 30. par (30) 31. le (31) 32. climat (32) 33. des (33) 34. affaires (34) 35. n'absorbe (35) 36. pas (36) 37. cette (37) 38. main-d'œuvre (38). 39. Il (39) 40. faut (40) 41. réformer (41) 42. l'éducation (42) 43. aux (43) 44. besoins (44) 45. du (45) 46. marché (46) 47. et (47) 48. assainir (48) 49. le (49) 50. climat (50) 51. des (51) 52. affaires (52) 53. pour (53) 54. attirer (54) 55. l'investissement (55). 56. L'avenir (56) 57. national (57) 58. en (58) 59. dépend (59).

\textbf{Résultat : 59 mots.} Toujours légèrement au-dessus (56 max). Dernière coupe : supprimer « nationale » (mot 8), « cette » (mot 37), « national » (mot 57).
\textbf{Version finale (56 mots) :}
\emph{L'emploi des jeunes est le défi majeur. Le chômage des diplômés génère frustration et instabilité. Les programmes publics d'insertion échouent par lourdeurs et sous-financement. Le secteur privé, freiné par le climat des affaires, n'absorbe pas la main-d'œuvre. Il faut réformer l'éducation aux besoins du marché et assainir le climat des affaires pour attirer l'investissement. L'avenir en dépend.}
\textbf{Contrôle :} 56 mots. Dans l'intervalle [46-56]. Sens fidèle. Ordre respecté. Neutralité totale. Connecteurs logiques (par, et, pour) efficaces.
\tcblower
\textbf{Analyse pédagogique de la résolution :} 
Notez l'usage massif de la \textbf{nominalisation} (\emph{chômage, insertion, sous-financement, absorption, investissement}) et de la \textbf{généralisation} (\emph{programmes publics} au lieu de citer PIAASI). La \textbf{suppression} a éliminé les chiffres (« milliers »), les noms propres de programmes, et la phrase finale légèrement redondante a été condensée (« L'avenir en dépend »). Le taux est respecté de justesse, prouvant que le peaufinage final (étape 4) est critique.
\end{exoresolu}

\courssub{Exemple chiffré de référence (Modèle minimal)}

\begin{exemplebox}[Modèle de gain : Extrait « Le Vieux nègre et la médaille »]
\textbf{Source (22 mots) :} 
\emph{« Meka, le vieux chef, attend avec une impatience fébrile l'arrivée de l'administrateur qui doit lui remettre la médaille de la Légion d'honneur. »}

\textbf{Cible (13 mots) :} 
\emph{« Meka attend fébrilement l'administrateur pour la remise de sa médaille. »}

\textbf{Analyse des opérations :}
\begin{itemize}
    \item \textbf{Suppression :} « le vieux chef », « avec une impatience », « l'arrivée de », « qui doit lui », « de la Légion d'honneur » (détails honorifiques non essentiels à l'action centrale).
    \item \textbf{Généralisation :} « impatience fébrile » $\rightarrow$ « fébrilement » (adverbe) ; « remise de la médaille » (nominalisation).
    \item \textbf{Substitution :} « attend ... l'administrateur qui doit lui remettre » $\rightarrow$ « attend l'administrateur pour la remise » (infinitif de but).
\end{itemize}
\textbf{Gain :} 40,9\% (réduction à 59\% du texte). \textbf{Attention :} Ce taux dépasse le tiers requis (33\%). Cet exemple illustre uniquement les \emph{opérations linguistiques} ; au Bac, il faudrait viser 7-8 mots (ex: \emph{« Meka attend fébrilement l'administrateur pour sa médaille »}, 10 mots, encore haut ; \emph{« Meka attend l'administrateur pour sa médaille »}, 8 mots, dans la norme).
\end{exemplebox}

\courssub{Erreurs fatales à éviter (Check-list de relecture)}

\begin{attention}[Les 5 erreurs éliminatoires au Bac]
\begin{enumerate}
    \item \textbf{Hors taux :} $< 45$ ou $> 55$ mots (pour une cible de 50). \emph{Solution : Comptage au brouillon, ajustement lexical.}
    \item \textbf{Subjectivité :} « Je pense que l'auteur a raison », « Heureusement, le gouvernement agit ». \emph{Solution : Impersonnel / Discours indirect.}
    \item \textbf{Désordre logique :} Inverser l'argument 2 et 3. \emph{Solution : Suivre le plan du source à la lettre.}
    \item \textbf{Perte de modalité :} Transformer « L'auteur suggère » en « L'auteur prouve ». \emph{Solution : Reporter les verbes introducteurs (suggérer, affirmer, nier, ironiser).}
    \item \textbf{Télégraphique / Agrammatical :} « Emploi jeunes défi majeur. Diplômés chômage. Instabilité. Programmes échec. Privé freiné. Réformer éducation. » \emph{Solution : Phrases complètes, connecteurs, syntaxe correcte.}
\end{enumerate}
\end{attention}

\courssub{Entraînement guidé : Application immédiate}

\begin{exercice}[Contraction d'un paragraphe argumentatif court]
Contractez le texte ci-dessous en \textbf{40 mots $\pm 10\%$ (36 à 44 mots)}. Appliquez la méthode en 4 étapes. Présentez votre plan détaillé et le comptage final.

\textbf{Texte source (124 mots) :}
\emph{« La corruption gangrène l'administration camerounaise à tous les niveaux. Du bas de l'échelle au sommet de la hiérarchie, les fonctionnaires exigent des pots-de-vin pour accomplir les tâches les plus élémentaires : délivrer un acte de naissance, signer un permis de construire, valider un dossier de retraite. Cette pratique, devenue quasi systémique, décourage l'investissement étranger et pénalise les citoyens les plus démunis qui ne peuvent pas payer. Le gouvernement a mis en place la CONAC et le Tribunal Criminel Spécial, mais la répression seule ne suffit pas tant les mentalités sont ancrées. Il faut donc coupler la sanction à une éducation civique renouvelée et à la numérisation des procédures administratives pour supprimer le contact direct entre l'usager et l'agent corrompu. Seule une approche globale éradiquera ce fléau. »}
\end{exercice}

\begin{corrige}[Exercice n°1 - Proposition de correction]
\textbf{Étape 1 : Comptage source = 124 mots. Cible = 41 mots. Intervalle [37 - 45 mots].}
\textbf{Idées forces :}
1. Corruption généralisée dans l'administration (tous niveaux).
2. Pot-de-vin pour actes élémentaires (exemples : naissance, construire, retraite).
3. Conséquences : décourage investissement, pénalise pauvres.
4. Réponse gov : CONAC, TCS (répression) mais insuffisante (mentalités).
5. Solutions nécessaires : Éducation civique + Numérisation (supprimer contact).
6. Conclusion : Approche globale seule efficace.

\textbf{Étape 2 : Plan détaillé (Squelette).}
\textbf{Étape 3 : Rédaction optimisée.}
\emph{La corruption gangrène toute l'administration : pots-de-vin exigés pour actes élémentaires. Cela décourage l'investissement et pénalise les démunis. La répression (CONAC, TCS) échoue face aux mentalités ancrées. Il faut y ajouter éducation civique et numérisation des procédures pour supprimer le contact usager-agent. Seule une approche globale éradiquera ce fléau.}

\textbf{Étape 4 : Comptage final.}
1.La 2.corruption 3.gangrène 4.toute 5.l'administration 6.: 7.pots-de-vin 8.exigés 9.pour 10.actes 11.élémentaires 12.Cela 13.décourage 14.l'investissement 15.et 16.pénalise 17.les 18.démunis 19.La 20.répression 21.(CONAC, 22.TCS) 23.échoue 24.face 25.aux 26.mentalité 27.ancrées 28.Il 29.faut 30.y 31.ajouter 32.éducation 33.civique 34.et 35.numérisation 36.des 37.procédures 38.pour 39.supprimer 40.le 41.contact 42.usager-agent 43.Seule 44.une 45.approche 46.globale 47.éradiquera 48.ce 49.fléau.
\textbf{Total : 49 mots.} \textbf{Hors tolérance haute (45).}

\textbf{Ajustement final (Coupe 4 mots) :} Supprimer « toute » (4), « élémentaires » (11), « y » (30), « ce » (48).
\emph{La corruption gangrène l'administration : pots-de-vin exigés pour actes. Cela décourage l'investissement et pénalise les démunis. La répression (CONAC, TCS) échoue face aux mentalités ancrées. Il faut ajouter éducation civique et numérisation des procédures pour supprimer contact usager-agent. Seule approche globale éradiquera fléau.}
\textbf{Comptage ajusté : 45 mots. \checkmark Dans l'intervalle.}
\end{corrige}

\courssub{Synthèse et projection vers la Section 3}
Maîtriser la contraction, c'est maîtriser la \textbf{pensée synthétique} : distinguer l'essentiel de l'accessoire, restructurer l'information sans la trahir, manier la syntaxe pour la densité. Ces compétences reposent sur une connaissance précise des \textbf{outils de liaison} (conjonctions, prépositions, adverbes de liaison) qui assurent la cohérence du texte réduit. C'est l'objet de la prochaine section : \emph{« Outils de liaison : Conjonctions de coordination et de subordination »}.

\coursec{Outils de liaison : Conjonctions de coordination et de subordination}

Dans la section précédente, nous avons appris à repérer la structure argumentative d'un texte et à le réduire en supprimant les éléments secondaires. Mais un résumé n'est pas une simple suite de phrases amputées : c'est un texte \cle{fluide}, où les idées s'enchaînent logiquement. Pour relier ces idées entre elles sans alourdir, le rédacteur dispose d'outils précis : les \cle{conjonctions}. C'est l'objet de cette troisième séquence. Bien les maîtriser, c'est gagner en cohérence, en économie de mots et en points au Bac.

\courssub{Révision : les conjonctions de coordination}

\begin{definition}[Conjonction de coordination]
Une \cle{conjonction de coordination} est un mot invariable qui relie deux unités de \textbf{même nature} et de \textbf{même fonction} : deux mots, deux groupes de mots ou deux propositions. Il en existe sept, mémorisables grâce au moyen mnémotechnique : \emph{Mais où est donc Ornicar ?} — \textbf{mais, ou, et, donc, or, ni, car}.
\end{definition}

L'intuition est simple : la coordination met deux éléments sur un pied d'\textbf{égalité}. Aucune des deux propositions ne dépend de l'autre ; on peut les inverser ou les séparer sans détruire le sens grammatical.

\begin{exemplebox}[La coordination à l'œuvre]
\begin{itemize}
\item \emph{Meka est vieux, \textbf{mais} il reste courageux.} (opposition entre deux propositions indépendantes)
\item \emph{Il reçoit la médaille \textbf{et} le village l'acclame.} (addition)
\item \emph{Il accepte l'honneur, \textbf{car} il ne comprend pas le piège.} (explication)
\end{itemize}
Dans chaque cas, chaque proposition pourrait vivre seule : \emph{Meka est vieux. Il reste courageux.}
\end{exemplebox}

\begin{attention}[Le piège de « car »]
\emph{Car} exprime l'explication, mais c'est une conjonction de \textbf{coordination} : il ne peut \textbf{jamais} commencer une phrase. On écrit : \emph{Il accepta, car il était flatté.} et non \emph{Car il était flatté, il accepta.} En revanche, \emph{parce que} (subordination) peut tout à fait ouvrir la phrase : \emph{Parce qu'il était flatté, il accepta.} Confondre les deux est une faute classique qui coûte cher dans un résumé.
\end{attention}

\courssub{Les conjonctions de subordination}

\begin{definition}[Conjonction de subordination]
Une \cle{conjonction de subordination} est un mot ou une locution invariable (\emph{que, si, quand, comme, parce que, bien que, afin que...}) qui introduit une \textbf{proposition subordonnée}, c'est-à-dire une proposition qui \textbf{dépend} d'une proposition principale et ne peut pas être employée seule.
\end{definition}

Ici, les deux idées ne sont plus égales : il y a une \textbf{idée principale} et une \textbf{idée secondaire} qui la complète. La subordination est donc l'outil privilégié du résumé, car elle permet de hiérarchiser les idées — exactement ce qu'exige la contraction de texte.

On distingue trois grandes familles de subordonnées :
\begin{itemize}
\item les \textbf{complétives} : \emph{Il dit \textbf{que} la médaille est un honneur.} (introduites par \emph{que}, \emph{si}) ;
\item les \textbf{relatives} : \emph{La médaille \textbf{qu'}il reçut le perdit.} (introduites par \emph{qui, que, dont, où}) ;
\item les \textbf{circonstancielles}, qui sont notre objet principal.
\end{itemize}

\begin{definition}[Proposition subordonnée circonstancielle]
Une \cle{proposition subordonnée circonstancielle} précise les \textbf{circonstances} de l'action exprimée par la principale : sa \textbf{cause}, sa \textbf{conséquence}, son \textbf{but}, une \textbf{concession}, une \textbf{condition} ou le \textbf{temps}. Elle est introduite par une conjonction ou une locution conjonctive et peut souvent être déplacée dans la phrase : \emph{Bien qu'il fût fatigué, il dansa} = \emph{Il dansa, bien qu'il fût fatigué.}
\end{definition}

\begin{popfigure}
\begin{tikzpicture}[node distance=1.4cm, every node/.style={font=\small, align=center}]
\node[draw, ellipse, fill=popblueL, thick, popblue, align=center] (coord){Coordination\\ \emph{mais, ou, et, donc, or, ni, car}\\ Proposition 1 = Proposition 2};
\node[draw, ellipse, fill=poporangeL, thick, poporange, right=3.2cm of coord, align=center] (subord){Subordination\\ \emph{que, si, quand, comme, bien que...}\\ Principale $>$ Subordonnée};
\node[draw, rectangle, rounded corners, fill=popgreenL, thick, popgreen, below=1.6cm of coord, text width=5.2cm] (circ) {Circonstancielles :\\ cause, conséquence, but,\\ concession, condition, temps};
\node[draw, rectangle, rounded corners, fill=poppinkL, thick, poppink, below=1.6cm of subord, text width=5.2cm] (comp) {Complétives : \emph{que, si}\\ Relatives : \emph{qui, que, dont, où}};
\draw[->, very thick, popblue] (coord.south) -- (circ.north) node[midway, right, font=\footnotesize] {indépendance};
\draw[->, very thick, poporange] (subord.south) -- (comp.north) node[midway, right, font=\footnotesize] {dépendance};
\draw[<->, thick, popmuted] (coord.east) -- (subord.west) node[midway, above, font=\footnotesize] {deux logiques};
\end{tikzpicture}
\legende{Schéma fonctionnel des deux grandes liaisons : la coordination associe des propositions égales ; la subordination hiérarchise une idée principale et une idée dépendante.}
\end{popfigure}

\courssub{Tableau de correspondance : relation logique et connecteurs}

Pour choisir le bon outil, il faut d'abord nommer la relation logique qui unit deux idées. Voici le tableau de correspondance à connaître pour le Bac.

\begin{center}
\begin{tabularx}{\linewidth}{|l|X|X|}
\hline
\rowcolor{popblueL} \textbf{Relation logique} & \textbf{Conjonctions / locutions} & \textbf{Mode du verbe} \\
\hline
Cause & parce que, comme, puisque, vu que & indicatif \\
\hline
Conséquence & si bien que, de sorte que, donc, alors & indicatif \\
\hline
But & pour que, afin que, de peur que & \textbf{subjonctif} \\
\hline
Concession & bien que, quoique, encore que, même si & \textbf{subjonctif} (sauf \emph{même si} : indicatif) \\
\hline
Condition & si, à condition que, pourvu que & indicatif après \emph{si} ; subjonctif après \emph{à condition que} \\
\hline
Temps & quand, lorsque, dès que, pendant que, avant que & indicatif ; subjonctif après \emph{avant que} \\
\hline
Opposition & mais, alors que, tandis que & indicatif \\
\hline
\end{tabularx}
\end{center}

\courssub{Le piège du subjonctif}

\begin{definition}[Le mode subjonctif]
Le \cle{subjonctif} est le mode de l'\textbf{éventualité}, du souhait, du doute, du sentiment ou du but non encore réalisé. Il s'oppose à l'\textbf{indicatif}, mode du fait réel et certain. Il est imposé par certaines conjonctions et locutions conjonctives.
\end{definition}

\begin{propriete}[Règle du subjonctif obligatoire]
Le \textbf{subjonctif} est obligatoire après :
\begin{itemize}
\item les locutions de \textbf{concession} : \emph{bien que, quoique, encore que} ;
\item les locutions de \textbf{but} : \emph{pour que, afin que} ;
\item les locutions de \textbf{crainte} : \emph{de peur que, de crainte que} ;
\item les verbes de \textbf{doute} et de volonté : \emph{douter que, vouloir que, il faut que}.
\end{itemize}
Exemple correct : \emph{\textbf{Bien qu'il soit} vieux, il court.} — Faute grave : \emph{Bien qu'il \textbf{est} vieux...}
\end{propriete}

\begin{attention}[Erreurs fréquentes au Bac]
\begin{enumerate}
\item \emph{Bien qu'il est fatigué} : faute de mode. Écrire \emph{bien qu'il \textbf{soit} fatigué}.
\item \emph{Après qu'il soit venu} : faute. \emph{Après que} exige l'\textbf{indicatif} : \emph{après qu'il \textbf{est} venu} (le fait est accompli, donc réel).
\item \emph{Malgré qu'il soit} : faute d'usage. \emph{Malgré} est une préposition suivie d'un nom : \emph{malgré sa fatigue}.
\end{enumerate}
\end{attention}

\courssub{Méthode : choisir le bon connecteur}

\begin{methode}[Trois étapes pour relier deux idées]
\begin{enumerate}
\item \textbf{Identifier la relation logique} entre les deux idées : est-ce une opposition, une cause, une conséquence, un but, un ajout ?
\item \textbf{Choisir la structure} : coordination si les deux idées sont d'égale importance ; subordination s'il y a une idée principale et une idée secondaire (cas le plus fréquent dans un résumé).
\item \textbf{Vérifier la grammaire} : le connecteur choisi exige-t-il l'indicatif ou le subjonctif ? Contrôler la conjugaison du verbe introduit.
\end{enumerate}
\end{methode}

\begin{exemplebox}[Le gain de mots, un exemple chiffré]
Comparer trois formulations :
\begin{itemize}
\item \emph{Il est vieux. Il court.} : 6 mots, deux phrases sèches, aucun lien explicite.
\item \emph{Bien qu'il soit vieux, il court.} : 8 mots, mais \textbf{une seule phrase} fluide qui exprime la concession.
\item \emph{Malgré son âge, il court.} : 5 mots, grâce à la préposition \emph{malgré} + nom.
\end{itemize}
La subordination ne raccourcit pas toujours le décompte brut, mais elle \textbf{fusionne} les idées et rend le résumé plus dense et plus élégant ; la préposition, elle, offre la contraction maximale.
\end{exemplebox}

\courssub{Application : l'articulation des idées chez Ferdinand Oyono}

Dans \emph{Le Vieux Nègre et la médaille}, Oyono articule constamment les idées par la subordination, ce qui donne à son récit son ironie et sa fluidité. Observons quelques mécanismes représentatifs :
\begin{itemize}
\item \emph{Bien qu'il fût fatigué, Meka dansa toute la nuit.} : la subordonnée de concession (imparfait du subjonctif, \emph{fût}) souligne l'héroïsme dérisoire du vieil homme épuisé par les réjouissances.
\item \emph{Parce qu'il avait donné ses terres à la mission, il recevait une médaille.} : la subordonnée de cause expose le marché inégal que la coordination aurait masqué.
\item \emph{Il attendit que la pluie cessât.} : la complétive au subjonctif marque l'attente et l'incertitude du dénouement.
\end{itemize}
En résumant un passage du roman, l'élève doit conserver ces rapports logiques : supprimer la subordonnée de concession, c'est trahir la pensée de l'auteur.

\begin{exoresolu}[Transformer des phrases juxtaposées en phrase complexe]
Énoncé : transformez la suite suivante en une phrase complexe comportant au moins une subordonnée circonstancielle, sans changer le sens : \emph{Meka était épuisé. La pluie tombait. Il attendit le matin sous le hangar.}
\tcblower
Solution détaillée :
\begin{enumerate}
\item Relations logiques : \emph{épuisé} + \emph{attendit} = concession ; \emph{pluie} + \emph{attendit sous le hangar} = cause.
\item Choix des connecteurs : \emph{bien que} (concession, subjonctif) et \emph{parce que} (cause, indicatif).
\item Vérification des modes : \emph{bien qu'il fût} (subjonctif imparfait) ; \emph{parce que la pluie tombait} (indicatif).
\end{enumerate}
Phrase obtenue : \emph{Bien qu'il fût épuisé, Meka attendit le matin sous le hangar, parce que la pluie tombait.} — Trois phrases juxtaposées (13 mots, sans liens) deviennent une phrase unique, hiérarchisée et fluide, prête à intégrer un résumé.
\end{exoresolu}

\courssub{Boîte à outils personnelle : les 15 connecteurs indispensables}

\begin{aretenir}[Fiche récapitulative à recopier et à mémoriser]
\begin{itemize}
\item \textbf{Addition} : et, de plus, également.
\item \textbf{Opposition} : mais, cependant, alors que.
\item \textbf{Concession} : bien que (+ subj.), quoique (+ subj.), malgré (+ nom).
\item \textbf{Cause} : parce que, puisque, comme, car.
\item \textbf{Conséquence} : donc, si bien que, c'est pourquoi.
\item \textbf{But} : pour que (+ subj.), afin de (+ inf.).
\item \textbf{Condition} : si, à condition que (+ subj.).
\item \textbf{Temps} : quand, dès que, avant que (+ subj.).
\end{itemize}
Règle d'or : après \emph{bien que, quoique, pour que, afin que, avant que, à condition que}, toujours vérifier le \textbf{subjonctif}.
\end{aretenir}

\begin{savaistu}[Les connecteurs, un critère de notation au Bac]
À l'épreuve de contraction de texte du Baccalauréat technique, la \emph{qualité de la langue} représente environ 5 points sur 20. Les correcteurs distinguent explicitement les copies qui utilisent des connecteurs variés et correctement accordés (subordination, subjonctif maîtrisé) de celles qui se contentent de \emph{et} et de \emph{mais}. Une boîte à outils bien apprise, c'est donc une note directement améliorée.
\end{savaistu}

\begin{exercice}[Entraînement]
\begin{enumerate}
\item Complétez avec le connecteur qui convient et accordez le verbe : \emph{Meka accepta la médaille, \dots\ il (se douter) du mépris des administrateurs.} (concession)
\item Transformez en une phrase complexe : \emph{Le village était pauvre. Les villageois organisèrent une fête. Ils voulaient honorer Meka.}
\item Relevez, dans un chapitre du \emph{Vieux Nègre et la médaille}, trois conjonctions de subordination et précisez la relation logique exprimée.
\end{enumerate}
\end{exercice}

\begin{corrige}[Exercice : entraînement]
\begin{enumerate}
\item \emph{Meka accepta la médaille, \textbf{bien qu'il ne se doutât pas} du mépris des administrateurs.} Concession : \emph{bien que} + subjonctif imparfait (\emph{se doutât}).
\item \emph{Bien que le village fût pauvre, les villageois organisèrent une fête pour honorer Meka.} Concession (\emph{bien que} + subjonctif) puis but ramené à \emph{pour} + infinitif, ce qui contracte davantage.
\item Exemple attendu : \emph{parce que} (cause), \emph{bien que} (concession), \emph{quand} (temps) — l'élève justifie chaque relation à partir du contexte de la phrase relevée.
\end{enumerate}
\end{corrige}

En résumé, la coordination juxtapose des idées égales, la subordination les hiérarchise ; le choix du connecteur dépend de la relation logique, et sa grammaire (indicatif ou subjonctif) doit être vérifiée systématiquement. Ces outils de liaison seront mis en pratique dans l'atelier de reformulation de la section suivante, où il s'agira de produire un résumé complet, fluide et fidèle.

\coursec{Atelier pratique : Reformulation et synthèse des idées essentielles}

Dans la section précédente, nous avons appris à manier les conjonctions de coordination et de subordination, ces outils de liaison qui permettent d'articuler les idées entre elles. Mais à quoi servent ces connecteurs si l'on ne sait pas quelles idées relier ? C'est précisément l'objet de cet atelier : une fois les idées essentielles repérées dans le texte (savoir acquis en section 2), il faut les \emph{reformuler} avec ses propres mots, les condenser, puis les enchaîner grâce aux conjonctions étudiées. Nous passerons ici de la théorie à la pratique, en travaillant sur des extraits du \emph{Vieux nègre et la médaille} de Ferdinand Oyono.

\courssub{Paraphrase ou contraction : deux exercices à ne pas confondre}

Commençons par une intuition simple. Imaginez qu'un camarade absent vous demande de lui raconter un film. Vous pouvez soit lui raconter \emph{toutes} les scènes dans l'ordre, avec vos propres mots (cela prendrait autant de temps que le film), soit lui donner l'essentiel en cinq minutes. La première démarche est la \cle{paraphrase} ; la seconde est la \cle{contraction}.

\begin{definition}[Reformulation : paraphrase et contraction]
La \cle{reformulation} consiste à exprimer avec d'autres mots le contenu d'un texte. Elle prend deux formes distinctes :
\begin{itemize}
\item La \cle{paraphrase} : reprise intégrale du texte avec d'autres termes, \textbf{sans réduction de longueur ni suppression d'informations}. Elle sert à vérifier la compréhension ou à expliquer un passage difficile.
\item La \cle{contraction} (ou résumé) : reprise \textbf{sélective et réduite} du texte. On ne conserve que les idées essentielles, en respectant une consigne de longueur (au Probatoire et au Baccalauréat technique, généralement le \textbf{tiers} du texte source).
\end{itemize}
\end{definition}

\begin{attention}[Ne pas paraphraser quand on demande une contraction]
L'erreur la plus fréquente à l'examen est de produire une paraphrase déguisée : le candidat recopie toutes les idées en changeant quelques mots, et dépasse largement la limite imposée. La contraction exige des \textbf{coupes} : supprimer les exemples, les répétitions, les descriptions d'ambiance, les détails secondaires. Dire la même chose en moins de mots suppose d'avoir \emph{choisi} ce qu'on garde.
\end{attention}

\courssub{La réécriture par changement de point de vue}

Pour reformuler sans plagier le texte, une technique efficace consiste à \textbf{changer de point de vue énonciatif}. Cela oblige le rédacteur à s'approprier le contenu au lieu de le recopier.

\begin{methode}[Transformer le point de vue énonciatif]
\begin{enumerate}
\item \textbf{Discours direct $\rightarrow$ discours indirect} : « Meka s'écria : "Je suis un homme important désormais !" » devient « Meka affirma qu'il était désormais un homme important ». Les verbes de parole (dire, affirmer, déclarer) introduisent la subordonnée avec \emph{que}.
\item \textbf{Passé simple $\rightarrow$ passé composé ou imparfait} : « il rentra » devient « il est rentré » ou « il rentrait ». Cela rapproche le récit du lecteur et allège le style.
\item \textbf{3\textsuperscript{e} personne $\rightarrow$ tournure impersonnelle} : « Meka comprit qu'il avait été dupé » peut devenir « il apparaît que Meka a été dupé » ou « il faut reconnaître que la médaille n'a rien changé ». La tournure impersonnelle (\emph{il faut que}, \emph{il semble que}) donne un ton plus analytique, adapté au résumé.
\end{enumerate}
\end{methode}

\begin{definition}[Discours indirect libre et discours indirect intégré]
Le \cle{discours indirect} rapporte les paroles ou pensées d'un personnage dans une subordonnée introduite par \emph{que} : « Il dit qu'il était fier. » Le \cle{discours indirect libre} fusionne la pensée du personnage avec la voix du narrateur, sans verbe introducteur ni guillemets : « Meka regarda la médaille. Quelle dérision ! Tant de souffrances pour ce bout de métal. » Dans une contraction, on \textbf{intègre} systématiquement le discours : on transforme le direct et l'indirect libre en indirect classique, plus compact et plus neutre.
\end{definition}

\begin{definition}[Nominalisation]
La \cle{nominalisation} consiste à transformer un verbe (ou une proposition entière) en nom : « il décide de partir » $\rightarrow$ « sa décision de partir » ; « les Blancs l'ont décoré » $\rightarrow$ « sa décoration par les Blancs ». C'est l'outil de compression par excellence : une subordonnée de six mots devient un groupe nominal de trois mots.
\end{definition}

\courssub{La technique du « Noyau dur »}

Voici maintenant la méthode opérationnelle de réduction, à appliquer phrase par phrase lors du premier passage.

\begin{methode}[La technique du Noyau dur]
Pour chaque phrase du texte source :
\begin{enumerate}
\item \textbf{Isoler le verbe principal et son sujet} : c'est le noyau irremplaçable. Exemple : dans « Meka, épuisé par une longue journée d'attente sous le soleil implacable, rentra enfin chez lui », le noyau est « Meka rentra ».
\item \textbf{Garder les compléments indispensables} : les compléments d'objet et les compléments circonstanciels de \emph{cause} ou de \emph{but}, car ils portent la logique du texte (« chez lui » est indispensable ici ; « sous le soleil implacable » ne l'est pas).
\item \textbf{Supprimer les adjectifs qualificatifs non significatifs} : « implacable », « longue », « épuisé » sont des effets de style, pas des informations.
\item \textbf{Nominaliser les actions secondaires} : « épuisé par une longue journée d'attente » devient « après l'attente » ou « fatigué d'attendre ».
\end{enumerate}
Résultat : « Après une longue attente, Meka rentre chez lui, épuisé » (9 mots au lieu de 17).
\end{methode}

\begin{attention}[Conserver la nuance modale]
En réduisant, ne perdez jamais la \textbf{modalité} de l'énoncé. « Il \emph{devrait} partir » n'est pas « Il part » : la suppression du modal fait disparaître l'idée de conseil ou d'obligation. Conservez les verbes modaux (\emph{devoir, pouvoir, falloir}) et les adverbes de nuance (\emph{probablement, sans doute, peut-être}). De même, « Meka semble déçu » n'est pas « Meka est déçu » : le narrateur d'Oyono laisse souvent planer le doute, et le résumé doit le respecter.
\end{attention}

\courssub{Que faire des métaphores et des images d'Oyono ?}

Le style d'Oyono est riche en images. Faut-il les conserver dans une contraction ? La règle est celle de la \textbf{hiérarchie} :

\begin{itemize}
\item \textbf{Garder} l'image si elle est \emph{essentielle} au sens de l'œuvre. « La médaille » elle-même est le symbole central du roman : elle représente la fausse reconnaissance coloniale. On ne peut pas la supprimer ni la remplacer.
\item \textbf{Expliquer} l'image si elle est \emph{opaque} pour le lecteur. Si Oyono écrit que la médaille est « un hochet », la contraction peut expliciter : « un jouet dérisoire, une récompense sans valeur ».
\item \textbf{Généraliser} l'image si elle est \emph{secondaire}. Les détails de la chaleur, de la poussière, des uniformes peuvent se condenser en « dans une atmosphère solennelle et artificielle ».
\end{itemize}

\begin{popfigure}
\begin{tikzpicture}[scale=0.9, every node/.style={font=\small}]
\matrix (grid) [matrix of nodes, nodes={draw, minimum width=3cm, minimum height=1.2cm, align=center}, row sep=5mm, column sep=5mm] {
|[fill=popblueL]| Texte original (extrait de la fin) & |[fill=popgoldL]| Idées essentielles repérées (mots-clés) & |[fill=popgreenL]| Reformulation contractée \\
"Meka rentra chez lui, la médaille au cou. Il ne la regarda pas. Il la posa sur la table. Elle ne valait pas la peine qu'on la regardât." & Meka / rentre / médaille / refuse de la regarder / déception / inutilité & De retour chez lui, Meka dépose sa médaille sans la regarder, conscient de sa vanité. \\
};
\end{tikzpicture}
\legende{Les trois étapes de la contraction : du texte source aux mots-clés, puis à la phrase reformulée. La métaphore centrale (« la médaille ») est conservée ; les répétitions et le subjonctif imparfait sont éliminés.}
\end{popfigure}

\courssub{Travail sur corpus : trois micro-contractions}

L'atelier se déroule sur trois extraits représentatifs du roman, correspondant aux trois moments de la structure narrative. Consigne : produire pour chacun une \textbf{micro-contraction d'environ 50 mots}.

\begin{exemplebox}[Corpus de travail]
\textbf{Extrait 1 — Le début : l'attente.} Meka et les villageois d'Ekoumdoum attendent la visite du commandant. L'atmosphère est fiévreuse : on nettoie le village, on prépare les fêtes, Meka se vante de la médaille qu'il va recevoir.\par
\textbf{Extrait 2 — Le milieu : la cérémonie.} Le commandant décore Meka. Discours officiels, fanfare, embrassades. Mais Meka, perdu dans la foule des invités blancs, se sent déjà étranger à lui-même ; la nuit suivante, errant et ivre, il est arrêté par les gendarmes.\par
\textbf{Extrait 3 — La fin : le retour.} Meka rentre chez lui, humilié. La médaille, posée sur la table, n'est plus qu'un objet dérisoire. La communauté partage sa désillusion.
\end{exemplebox}

\begin{exoresolu}[Micro-contraction de l'extrait 1 (l'attente)]
Énoncé : réduisez en 50 mots environ le passage décrivant l'attente de la visite du commandant à Ekoumdoum.
\tcblower
Solution détaillée : on repère les mots-clés (Meka, village, attente, commandant, médaille, fierté, préparatifs). On supprime les descriptions pittoresques (nettoyage des cases, coiffures des femmes), on nominalise (« on prépara les fêtes » $\rightarrow$ « les préparatifs »), on intègre les propos de Meka au discours indirect.\\
\textbf{Production possible (48 mots)} : « À Ekoumdoum, l'annonce de la visite du commandant provoque une effervescence générale. Tandis que le village se prépare à fêter l'événement, Meka, que les Blancs doivent décorer, éprouve une fierté immense : il croit que cette médaille fera de lui l'égal des colons. »
\end{exoresolu}

\begin{exercice}[Micro-contractions des extraits 2 et 3]
Rédigez une contraction d'environ 50 mots pour l'extrait 2 (la cérémonie) et pour l'extrait 3 (le retour). Appliquez la technique du Noyau dur, intégrez le discours direct au discours indirect, et conservez la métaphore de la médaille en l'expliquant si nécessaire. Vous utiliserez au moins deux conjonctions de subordination par contraction.
\end{exercice}

\begin{corrige}[Exercice : micro-contractions des extraits 2 et 3]
\textbf{Extrait 2 (52 mots)} : « Lors de la cérémonie, le commandant remet sa médaille à Meka au milieu des discours officiels. Pourtant, au lieu de la joie attendue, le vieillard découvre l'isolement : égaré dans la nuit, il est arrêté par les gendarmes, ce qui révèle que la décoration ne lui confère aucune protection réelle. »\\
\textbf{Extrait 3 (47 mots)} : « De retour au village, Meka, humilié, comprend que sa médaille n'est qu'un objet dérisoire : elle n'a changé ni sa condition ni le regard des Blancs. Parce qu'elle partage sa désillusion, la communauté d'Ekoumdoum mesure le mensonge de la prétendue reconnaissance coloniale. »\\
Justification des choix : nominalisations (« la remise » $\rightarrow$ « remet » pour garder un verbe ; « l'arrestation » intégrée en proposition), modaux et nuances conservés (« prétendue »), image de la médaille explicitée (« objet dérisoire », « mensonge »).
\end{corrige}

\courssub{Mutualisation : l'évaluation par les pairs}

Les productions sont échangées en classe. Chaque élève corrige la copie d'un camarade à l'aide d'une grille critériée, ce qui développe l'esprit critique et prépare à l'auto-correction.

\begin{center}
\begin{tabularx}{\linewidth}{|l|X|c|}
\hline
\rowcolor{popblueL} \textbf{Critère} & \textbf{Attendu} & \textbf{Note /5} \\
\hline
Fidélité & Aucune idée ajoutée, aucune idée essentielle omise, nuance modale respectée & \\
\hline
Réduction & Nombre de mots conforme à la consigne (marge de 10 \% admise) & \\
\hline
Langue & Orthographe, grammaire, accords, ponctuation corrects & \\
\hline
Connecteurs & Présence de conjonctions de coordination et de subordination variées et pertinentes & \\
\hline
\end{tabularx}
\end{center}

\begin{savaistu}[Le correcteur du Bac déteste le « style télégramme »]
Les correcteurs du Probatoire et du Baccalauréat sanctionnent le résumé rédigé en « télégramme » : une suite de groupes nominaux sans verbes ni connecteurs (« Meka. Médaille. Retour village. Déception. »). Une phrase complète avec subordination — « Meka rentre au village, où il comprend que sa médaille est dérisoire » — vaut toujours mieux que trois groupes nominaux juxtaposés. Le résumé est un \emph{texte}, pas une liste de mots-clés.
\end{savaistu}

\courssub{Examen blanc : contraction intégrale d'un chapitre}

L'atelier se conclut par une épreuve blanche chronométrée (1 h 30) : contraction intégrale d'un chapitre du roman, avec une cible de \textbf{150 mots}.

\begin{exemplebox}[Objectif chiffré du Bac]
Au Baccalauréat technique, la consigne impose généralement le tiers du texte : un passage source de \textbf{200 mots} doit être réduit à environ \textbf{66 mots}, avec une marge tolérée de \textbf{60 à 72 mots}. Au-delà, le barème applique une pénalité de \textbf{0,5 point par tranche de 5 mots} en sus ou en moins. Pour l'examen blanc de cet atelier : chapitre source d'environ 450 mots, cible 150 mots, marge admise de 135 à 165 mots.
\end{exemplebox}

\begin{methode}[Gérer les 1 h 30 de l'épreuve blanche]
\begin{enumerate}
\item \textbf{Lecture analytique (20 min)} : lire deux fois le chapitre, souligner les mots-clés, numéroter les idées essentielles dans la marge.
\item \textbf{Plan de réduction (10 min)} : regrouper les idées en 4 ou 5 blocs, décider des coupes (exemples, répétitions, descriptions).
\item \textbf{Rédaction (40 min)} : rédiger en appliquant le Noyau dur, la nominalisation et le discours indirect ; compter les mots par lignes au brouillon.
\item \textbf{Relecture (20 min)} : vérifier la fidélité, les connecteurs, les accords ; recompter les mots et ajuster.
\end{enumerate}
\end{methode}

\courssub{Correction type au tableau}

La séance suivante est consacrée à la correction collective. Le professeur projette ou écrit au tableau une \textbf{contraction type}, puis la confronte aux copies des élèves. Trois questions guident l'analyse des écarts :

\begin{itemize}
\item \textbf{Pourquoi cette coupe ?} Chaque suppression doit être justifiée : « J'ai supprimé la description de l'uniforme du commandant parce qu'elle ne fait pas avancer l'action ni l'argumentation ; j'ai gardé l'arrestation de Meka parce qu'elle prouve l'inutilité de la médaille. »
\item \textbf{Pourquoi cette reformulation ?} On compare les nominalisations choisies, les transformations du discours, et l'on vérifie qu'aucune nuance modale n'a été perdue.
\item \textbf{Où sont les connecteurs ?} On souligne les conjonctions (\emph{tandis que, parce que, ce qui, pourtant}) et l'on montre comment elles reconstruisent la logique argumentative du texte en version réduite.
\end{itemize}

\begin{aretenir}[L'essentiel de l'atelier]
Reformuler, ce n'est pas recopier autrement : c'est \textbf{sélectionner} (les idées essentielles), \textbf{compresser} (Noyau dur, nominalisation, discours indirect) et \textbf{relier} (conjonctions de coordination et de subordination). La contraction réussie est fidèle au sens, conforme à la longueur imposée, rédigée en phrases complètes, et respecte les nuances du texte de départ. Ces quatre exigences — fidélité, réduction, langue, connecteurs — sont exactement celles de la grille de correction du Probatoire et du Baccalauréat technique.
\end{aretenir}

\coursec{Expression orale : Préparation et présentation d'un exposé}

Après avoir appris, dans l'atelier précédent, à repérer, reformuler et condenser les idées essentielles d'un texte à l'écrit, nous allons maintenant transposer ces compétences à l'\cle{oral}. Car résumer, c'est déjà sélectionner l'essentiel ; exposer, c'est le dire clairement devant un public. La contraction de texte et l'exposé oral partagent donc la même exigence : aller droit à l'idée, avec ordre et précision. Cette compétence est directement évaluée au Probatoire et au Baccalauréat technique, où l'épreuve orale de français mesure la \cle{compétence communicationnelle} du candidat.

\courssub{Qu'est-ce qu'un exposé oral ?}

\begin{definition}[L'exposé oral]
L'\cle{exposé oral} est une communication orale \textbf{structurée}, \textbf{documentée} et \textbf{limitée dans le temps} (généralement 5 à 10 minutes), portant sur un sujet précis, et suivie d'un temps d'\textbf{échanges} (questions du jury ou du public). Il repose sur un plan annoncé, des supports (fiches, diaporama, affiches) et une expression orale soignée en français standard.
\end{definition}

Il faut bien distinguer l'exposé des autres productions que tu connais :

\begin{center}
\begin{tabularx}{\linewidth}{|l|X|X|X|}
\hline
\rowcolor{popblueL} \textbf{Critère} & \textbf{Exposé} & \textbf{Dissertation} & \textbf{Compte-rendu} \\
\hline
Mode & Oral & Écrit & Écrit (ou oral bref) \\
\hline
But & Informer, expliquer, convaincre un public & Discuter un problème par arguments & Rapporter fidèlement des faits ou un texte \\
\hline
Durée / volume & 5 à 10 min & 2 à 4 pages & Court, objectif \\
\hline
Support & Fiches, diaporama, tableau & Copie d'examen & Document écrit \\
\hline
Échanges & Questions-réponses après & Aucun & Parfois \\
\hline
\end{tabularx}
\end{center}

\begin{attention}[Piège fréquent]
Un exposé n'est \textbf{pas} une dissertation lue à voix haute. Celui qui rédige intégralement son texte puis le lit tête baissée perd l'essentiel : le contact avec le public. L'exposé se \emph{prépare} par écrit, mais se \emph{réalise} à l'oral, avec des mots-clés pour seuls repères.
\end{attention}

\courssub{Les étapes de la préparation}

\begin{methode}[Préparer son exposé en 4 étapes]
\begin{enumerate}
\item \textbf{Analyser le sujet (la commande).} Dégager le thème, les limites du sujet, le public visé, le temps imparti. Formuler une \cle{problématique} : quelle question mon exposé doit-il résoudre ?
\item \textbf{Mener la recherche documentaire.} Consulter des sources fiables : manuels scolaires, encyclopédies, articles de presse sérieux, sites institutionnels. Noter les références au fur et à mesure pour la bibliographie.
\item \textbf{Construire un plan en 3 parties} (Introduction, Développement, Conclusion), rédiger les \cle{transitions} entre les parties et sélectionner les idées essentielles — c'est ici que la technique de la contraction de texte devient utile.
\item \textbf{Fabriquer des fiches bristol} et préparer le support visuel, puis \textbf{répéter} en se chronométrant.
\end{enumerate}
\end{methode}

\begin{definition}[La fiche bristol]
La \cle{fiche bristol} (ou aide-mémoire) est une petite carte cartonnée sur laquelle on note uniquement des \textbf{mots-clés}, des \textbf{chiffres}, des \textbf{citations courtes} et des \textbf{repères de plan}. Elle ne contient jamais de phrases complètes rédigées : son rôle est de relancer la mémoire, pas d'être lue.
\end{definition}

\begin{definition}[La transition orale]
La \cle{transition orale} est une phrase qui relie deux parties de l'exposé et guide l'auditeur. Exemples : \og Passons maintenant à l'analyse littéraire proprement dite \fg{}, \og Après avoir situé le contexte, nous verrons ensuite que les romanciers répondent par l'ironie \fg{}, \og Cela nous amène à notre dernière partie \fg{}.
\end{definition}

\begin{definition}[La bibliographie]
La \cle{bibliographie} est la liste ordonnée des sources utilisées. En norme simplifiée (inspirée des normes APA/ISO), on indique : \textbf{NOM de l'auteur, Prénom. \emph{Titre de l'ouvrage}. Lieu : Éditeur, année.} Pour un site : auteur ou organisme, titre de la page, adresse, date de consultation.
\end{definition}

\begin{exemplebox}[Une fiche bristol réussie]
Pour un exposé sur Ferdinand Oyono, une fiche peut contenir : \og Oyono, 1929-2010 / \emph{Le Vieux nègre et la médaille}, 1956 / Meka : médaille = récompense illusoire / ironie + satire coloniale / citation courte : \og la médaille \fg{} = symbole / transition : $\rightarrow$ Beti \fg. Huit lignes de mots-clés suffisent pour parler deux minutes.
\end{exemplebox}

\courssub{La structure de l'exposé : un parcours fléché}

Un exposé suit toujours la même architecture, que le public doit pouvoir suivre sans effort :

\begin{popfigure}
\begin{center}
\begin{tikzpicture}[node distance=1cm, every node/.style={font=\small, align=center}]
\node[draw, rectangle, rounded corners, fill=popblueL, align=center] (intro){Introduction\\(Accroche, Sujet, Problématique, Plan)};
\node[draw, rectangle, rounded corners, fill=popgreenL, below of=intro, align=center] (dev1){Partie I : Contexte historique\\(Colonisation, Assimilation)};
\node[draw, rectangle, rounded corners, fill=popgreenL, below of=dev1, align=center] (dev2){Partie II : Analyse littéraire\\(Oyono, Beti, Ironie)};
\node[draw, rectangle, rounded corners, fill=popgreenL, below of=dev2, align=center] (dev3){Partie III : Postérité\\(Littérature actuelle, Mémoire)};
\node[draw, rectangle, rounded corners, fill=poppinkL, below of=dev3, align=center] (ccl){Conclusion\\(Bilan, Ouverture, Merci)};
\draw[->, thick] (intro) -- (dev1);
\draw[->, thick] (dev1) -- node[right] {Transition 1} (dev2);
\draw[->, thick] (dev2) -- node[right] {Transition 2} (dev3);
\draw[->, thick] (dev3) -- (ccl);
\end{tikzpicture}
\end{center}
\legende{Schéma de la structure d'un exposé en trois parties, appliqué au thème de simulation : \og L'héritage colonial dans la littérature camerounaise d'expression française \fg{}. Les transitions, annoncées oralement, relient les parties.}
\end{popfigure}

\courssub{Les techniques de l'oral : le corps qui parle}

Un bon contenu mal dit est un exposé raté. Quatre paramètres font la différence :

\begin{itemize}
\item \textbf{La voix} : débit ni trop rapide (stress) ni trop lent (ennui), articulation nette, volume adapté au fond de la salle. Marquer de courtes pauses après chaque idée essentielle.
\item \textbf{Le regard} : balayer lentement l'ensemble du public, sans fixer le sol, le plafond ou ses fiches. Le regard crée le contact et capte l'attention.
\item \textbf{Les gestes} : posture droite, gestes naturels pour appuyer une idée, utilisation ponctuelle du tableau ou du pointeur. Éviter les gestes parasites (stylo, cheveux, poches).
\item \textbf{La gestion du stress} : respirer profondément avant de commencer (inspirer 4 secondes, expirer 6 secondes), sourire, accepter une petite hésitation sans paniquer : reprendre calmement sa fiche.
\end{itemize}

\begin{attention}[Les trois erreurs qui coûtent cher]
\begin{enumerate}
\item \textbf{Lire son texte tête baissée} : perte du contact visuel, débit monotone, note d'expression orale effondrée.
\item \textbf{Dépasser le temps imparti} : c'est le signe d'un manque de répétition chronométrée. Le jury peut interrompre le candidat, qui ne pourra pas conclure.
\item \textbf{Surcharger les slides} : police inférieure à 24 pt, plus de 6 lignes par diapositive, paragraphes entiers recopiés. Le support doit illustrer, pas remplacer la parole.
\end{enumerate}
\end{attention}

\courssub{Les supports visuels}

Le support visuel est \textbf{obligatoire} dans un exposé scolaire. Deux options :

\begin{itemize}
\item \textbf{Le diaporama} : maximum 8 diapositives — 1. Titre ; 2. Plan ; 3 à 6. Contenu (une idée par slide, mots-clés, une image ou un chiffre) ; 7. Conclusion ; 8. Bibliographie.
\item \textbf{L'affiche ou le tableau} : titre visible, plan apparent, schémas simples, couleurs lisibles de loin.
\end{itemize}

\begin{exemplebox}[Exemple chiffré : la répartition du temps]
Temps standard au Bac et aux concours : \textbf{10 min d'exposé + 5 min de questions}. Répartition conseillée des 10 minutes : Introduction : 1 min ; Développement : 6 min (2 min par partie) ; Conclusion : 1 min ; Transitions : 1 min au total ; Marge de sécurité : 1 min. Répéter avec un chronomètre jusqu'à tenir ce découpage.
\end{exemplebox}

\courssub{Méthode complète : réussir son exposé en trois temps}

\begin{methode}[AVANT -- PENDANT -- APRÈS]
\textbf{AVANT :} construire le plan détaillé ; fabriquer les fiches bristol ; répéter à voix haute en se chronométrant (au moins 3 fois) ; tester le matériel (vidéoprojecteur, fichier, craies).\par
\textbf{PENDANT :} commencer par une \cle{accroche} (question, citation, chiffre) ; annoncer le plan ; regarder le public ; articuler ; expliciter chaque transition (\og Premièrement\ldots Ensuite\ldots En définitive\ldots \fg) ; terminer par une conclusion percutante et remercier.\par
\textbf{APRÈS :} répondre calmement aux questions (écouter jusqu'au bout, reformuler si besoin, répondre brièvement) ; remplir sa grille d'auto-évaluation ; noter les conseils des pairs.
\end{methode}

\begin{savaistu}[L'oral au Probatoire et au Bac technique]
L'épreuve orale de français au Probatoire et au Baccalauréat technique évalue la \og compétence communicationnelle \fg{} du candidat. Un exposé bien structuré, porté par des connecteurs oraux clairs (\og Premièrement \fg{}, \og En outre \fg{}, \og En définitive \fg{}), peut rapporter jusqu'à 4 points sur la prestation globale. Autant dire que la forme orale n'est pas un luxe : c'est une stratégie de réussite.
\end{savaistu}

\courssub{Simulation en classe : l'exposé par groupes}

\begin{experience}[Simulation : exposé par groupes de 3]
\textbf{Thème transversal :} \og L'héritage colonial dans la littérature camerounaise d'expression française \fg{} (en lien avec Ferdinand Oyono, Mongo Beti et Patrice Nganang).\par
\textbf{Matériel :} fiches bristol, diaporama ou affiches, chronomètre, grilles d'évaluation.\par
\textbf{Protocole :} chaque groupe répartit les rôles : un élève pour l'introduction et la Partie I (contexte colonial), un élève pour la Partie II (analyse d'\emph{Une vie de boy} ou du \emph{Vieux nègre et la médaille} et de l'ironie chez Oyono et Beti), un élève pour la Partie III (postérité : Patrice Nganang, littérature actuelle) et la conclusion. Répétition chronométrée, puis passage devant la classe.\par
\textbf{Observations attendues :} respect du temps, transitions audibles entre les parties, regard du public, support lisible.\par
\textbf{Interprétation :} l'exposé réussi est celui où le public peut restituer le plan et les idées essentielles sans avoir pris de notes.
\end{experience}

\courssub{L'évaluation : grille officielle et retour des pairs}

\begin{center}
\begin{tabularx}{\linewidth}{|l|X|c|}
\hline
\rowcolor{popblueL} \textbf{Critère MINESEC} & \textbf{Ce qui est observé} & \textbf{Points} \\
\hline
Contenu & Justesse et richesse des informations, sources citées & /5 \\
\hline
Structure & Plan en 3 parties, transitions explicites & /4 \\
\hline
Expression orale & Français standard, aisance, voix, regard & /5 \\
\hline
Support & Qualité, lisibilité, pertinence des visuels & /3 \\
\hline
Respect du temps & 10 min $\pm$ 1 min & /3 \\
\hline
\textbf{Total} & & \textbf{/20} \\
\hline
\end{tabularx}
\end{center}

Après chaque exposé, les camarades formulent un retour selon la méthode du \cle{feedback sandwich} : un point positif, un point à améliorer, un point positif. Exemple : \og Ton introduction était accrocheuse ; pense à regarder davantage le fond de la salle ; ta conclusion était très claire. \fg{} L'orateur complète ensuite sa propre grille d'auto-évaluation et la compare à celle du jury.

\begin{exoresolu}[Construire l'introduction d'un exposé]
\textbf{Énoncé.} Rédige l'introduction orale (environ 1 minute) de l'exposé : \og L'héritage colonial dans la littérature camerounaise d'expression française \fg{}. Elle doit contenir une accroche, la présentation du sujet, une problématique et l'annonce du plan.
\tcblower
\textbf{Solution détaillée.}\par
\emph{Accroche :} \og En 1956, un écrivain camerounais raconte l'histoire d'un vieil homme à qui l'administration coloniale remet une médaille. Cette médaille, censée être un honneur, devient dans le roman un symbole d'humiliation. \fg{}\par
\emph{Sujet :} \og C'est toute la question de notre exposé : comment la littérature camerounaise d'expression française, de Ferdinand Oyono à Patrice Nganang, s'est-elle construite sur l'héritage de la colonisation ? \fg{}\par
\emph{Problématique :} \og Comment les écrivains camerounais dénoncent-ils, dépassent-ils ou assument-ils cet héritage colonial ? \fg{}\par
\emph{Annonce du plan :} \og Pour répondre à cette question, nous verrons d'abord le contexte historique de la colonisation et de l'assimilation ; nous analyserons ensuite l'ironie chez Oyono et Mongo Beti ; enfin, nous montrerons la postérité de cette littérature dans l'œuvre de Patrice Nganang. \fg{}\par
Cette introduction tient en une minute, respecte les quatre éléments exigés et annonce clairement le fil conducteur de l'exposé.
\end{exoresolu}

\begin{exercice}[À toi de jouer]
Par groupes de trois, préparez en 30 minutes un mini-exposé de 5 minutes sur le thème : \og La médaille de Meka : honneur ou piège ? \fg{} à partir du \emph{Vieux nègre et la médaille}. Fabriquez trois fiches bristol (une par élève), répétez en vous chronométrant, puis présentez devant la classe. Les autres groupes remplissent la grille d'évaluation et formulent un feedback sandwich.
\end{exercice}

\begin{corrige}[Exercice : mini-exposé sur la médaille de Meka]
Éléments attendus : \textbf{Introduction} avec accroche (la scène de la remise de la médaille), problématique (\og la médaille est-elle une récompense ou un instrument de domination ? \fg) et plan annoncé. \textbf{Partie I} : le contexte — Meka, vieux nègre ayant donné ses terres et ses fils à la France, reçoit une médaille ; l'administration coloniale instrumentalise cette cérémonie. \textbf{Partie II} : l'analyse — l'ironie d'Oyono : la médaille n'apporte ni terre ni dignité ; Meka, ivre et arrêté pendant la tornade, découvre que la \og reconnaissance \fg{} coloniale est illusoire. \textbf{Partie III} : l'ouverture — la médaille comme symbole de l'assimilation trompeuse, thème repris par toute la littérature africaine francophone. \textbf{Conclusion} : bilan (la médaille est un piège symbolique) et ouverture (vers Mongo Beti). Critères oraux observés : voix audible, regard balayant le public, transitions explicites, temps respecté (5 min $\pm$ 30 s).
\end{corrige}

\begin{aretenir}[L'essentiel de la section]
Un exposé est une communication orale structurée, documentée et chronométrée, suivie d'échanges. Il se prépare en quatre étapes (analyse du sujet, recherche, plan en trois parties, fiches bristol) et se réussit en trois temps : AVANT (répétition chronométrée), PENDANT (accroche, plan annoncé, regard, transitions explicites, conclusion percutante), APRÈS (questions, auto-évaluation). Le support visuel illustre sans se substituer à la parole, et le respect du temps est un critère d'évaluation à part entière.
\end{aretenir}

\coursec{Évaluation intégrée : Compte-rendu, Remédiation et Projection Bac}

Après l'atelier d'expression orale de la section précédente, où vous avez appris à préparer et à présenter un exposé en mobilisant les connecteurs et les idées essentielles, il est temps de faire le point. Toute séquence d'apprentissage sérieuse se termine par un moment de vérité : l'évaluation. Mais attention, évaluer ne signifie pas seulement « noter ». Évaluer, c'est d'abord \cle{mesurer} où l'on en est, \cle{comprendre} ses erreurs, puis \cle{progresser} grâce à des activités de remédiation. Cette dernière section vous prépare aussi concrètement à l'épreuve de Français du Baccalauréat technique, où la contraction de texte est un exercice incontournable.

\courssub{Bilan des acquis : l'auto-positionnement}

Avant toute épreuve, le bon élève se pose une question simple : « Qu'est-ce que je sais vraiment faire ? ». C'est le principe de l'\cle{auto-positionnement} : on s'évalue soi-même, honnêtement, compétence par compétence, pour savoir où concentrer ses efforts.

\begin{definition}[Évaluation formative et évaluation sommative]
L'\cle{évaluation formative} est une évaluation qui a lieu \emph{pendant} l'apprentissage : elle ne vise pas à donner une note définitive, mais à repérer les difficultés pour les corriger (exemple : les exercices des sections précédentes). L'\cle{évaluation sommative} est une évaluation \emph{bilan}, qui certifie les acquis à la fin d'une séquence ou d'une année : elle donne lieu à une note qui compte (exemple : la composition, le Probatoire, le Baccalauréat).
\end{definition}

Voici la grille d'auto-positionnement sur les quatre grandes compétences de l'unité. Pour chaque item, cochez votre niveau : 1 = non acquis, 2 = en cours, 3 = acquis.

\begin{center}
\begin{tabularx}{\linewidth}{|l|X|c|c|c|}
\hline
\rowcolor{popblueL} \textbf{Compétence} & \textbf{Je sais...} & \textbf{1} & \textbf{2} & \textbf{3} \\
\hline
Lire / Comprendre & identifier le thème et la thèse d'un texte argumentatif & & & \\
\hline
Lire / Comprendre & repérer les arguments et les exemples qui soutiennent la thèse & & & \\
\hline
Écrire / Contracter & respecter le taux de réduction imposé (ex. 1/4 du texte) & & & \\
\hline
Écrire / Contracter & reformuler sans recopier les phrases du texte & & & \\
\hline
Parler / Exposer & présenter oralement une synthèse en un temps limité & & & \\
\hline
Langue / Connecteurs & employer correctement conjonctions de coordination et de subordination & & & \\
\hline
Langue / Connecteurs & utiliser le subjonctif après les bonnes conjonctions & & & \\
\hline
\end{tabularx}
\end{center}

\begin{methode}[Interpréter sa grille d'auto-positionnement]
\begin{enumerate}
\item Totalisez vos points (maximum 21).
\item De 18 à 21 : vous êtes prêt pour l'épreuve sommative ; visez l'excellence.
\item De 12 à 17 : des acquis solides existent ; identifiez les items notés 1 ou 2 et rejoignez l'atelier de remédiation correspondant.
\item Moins de 12 : reprenez les sections 1 à 4 de la leçon avant l'épreuve, avec l'aide du professeur.
\end{enumerate}
\end{methode}

\courssub{Épreuve sommative type Bac (2 heures)}

L'épreuve ci-dessous reproduit fidèlement les conditions du Baccalauréat technique. Elle est notée sur 20 points.

\begin{exercice}[Sujet type Bac — Contraction et étude de la langue (2 h, 20 pts)]
\textbf{Texte support (inédit, 280 mots)} : extrait argumentatif sur « Le téléphone portable à l'école : outil ou obstacle ? » (texte distribué en classe).

\textbf{I. Contraction (12 pts)} : Résumez le texte en \textbf{80 à 100 mots}. Vous indiquerez le nombre de mots à la fin de votre résumé. Tout écart supérieur à 10 \% du taux imposé sera sanctionné.

\textbf{II. Questions de linguistique (6 pts)} :
\begin{enumerate}
\item Relevez dans le texte deux conjonctions de subordination et précisez la relation logique qu'elles expriment (cause, concession, but, condition). (2 pts)
\item « Bien que les élèves adorent leur téléphone, il nuit à leur concentration. » Justifiez le mode du verbe « adorent » et proposez une reformulation avec « malgré ». (2 pts)
\item Transformez à la voix passive : « Les enseignants confisquent les téléphones en classe. » (2 pts)
\end{enumerate}

\textbf{III. Reformulation (2 pts)} : Reformulez l'idée essentielle du dernier paragraphe en une seule phrase, sans reprendre aucune expression de plus de trois mots du texte.
\end{exercice}

\begin{attention}[Les deux erreurs qui coûtent cher]
\begin{itemize}
\item \textbf{Négliger les questions de linguistique.} Beaucoup de candidats se jettent sur le résumé et bâclent les questions de langue. Or celles-ci rapportent souvent 5 à 6 points \emph{faciles} : une conjonction relevée, une transformation correcte, c'est un point gagné en deux minutes. Traitez-les en premier si vous manquez de temps.
\item \textbf{Ne pas relire sa copie.} Coquilles, accords oubliés, mots manquants : trois minutes de relecture finale peuvent sauver un à deux points. Prévoyez-les dans votre gestion du temps.
\end{itemize}
\end{attention}

\courssub{Correction collective commentée : analyse de copies types}

La correction collective est un moment d'apprentissage à part entière. On y compare trois copies anonymisées représentatives des trois niveaux constatés.

\begin{definition}[Grille critériée]
Une \cle{grille critériée} est un outil de correction qui décompose la tâche en critères précis, chacun associé à des \cle{indicateurs observables} (ce qu'on peut voir concrètement dans la copie) et à un barème. Elle rend la notation objective et identique pour tous les correcteurs.
\end{definition}

\begin{center}
\begin{tabularx}{\linewidth}{|l|X|c|}
\hline
\rowcolor{popblueL} \textbf{Critère} & \textbf{Indicateur observable} & \textbf{Barème} \\
\hline
Fidélité au sens & Aucune idée ajoutée, aucune idée essentielle omise, aucun contresens & /4 \\
\hline
Respect du taux & Nombre de mots compris entre 80 et 100, indiqué sur la copie & /2 \\
\hline
Reformulation & Phrases personnelles, pas de recopie de plus de 3 mots & /3 \\
\hline
Qualité de la langue & Orthographe, accords, ponctuation, connecteurs variés et corrects & /3 \\
\hline
\end{tabularx}
\end{center}

\begin{exemplebox}[Trois copies types commentées]
\textbf{Copie A — Excellente (16--18/20).} Le résumé fait 92 mots (taux respecté), s'ouvre sur la thèse de l'auteur, enchaîne les arguments avec \emph{« d'abord... ensuite... cependant... c'est pourquoi »}. Aucun contresens, deux petites fautes d'accord. Points forts : fidélité, connecteurs de concession bien employés.

\textbf{Copie B — Moyenne (10--13/20).} Le résumé fait 130 mots (taux dépassé : $-1$ pt). Les idées sont fidèles mais deux phrases sont recopiées presque mot pour mot. Les connecteurs se limitent à \emph{« et »} et \emph{« mais »}. Faiblesse récurrente : la réduction n'a pas été assez travaillée.

\textbf{Copie C — Insuffisante (moins de 10/20).} Le résumé contient un \cle{contresens} : l'élève fait dire à l'auteur le contraire de sa thèse (il écrit que l'auteur est favorable au téléphone en classe alors que l'auteur conclut à une utilisation encadrée). Le texte est truffé de fautes et dépasse 150 mots. Cause principale : la lecture méthodique (étape 1) a été négligée.
\end{exemplebox}

Le graphique ci-dessous compare les notes moyennes des trois profils de copies sur les deux grands critères de la grille.

\begin{popfigure}
\begin{tikzpicture}[scale=0.8, every node/.style={font=\small}]
\begin{axis}[
    ybar, 
    bar width=10pt, 
    width=12cm, 
    height=6cm, 
    ylabel={Note /20}, 
    symbolic x coords={Excellente, Moyenne, Insuffisante}, 
    xtick=data, 
    nodes near coords, 
    ymin=0, 
    ymax=20,
    xticklabel style={align=center, text width=2.5cm},
    legend style={at={(0.5,-0.25)}, anchor=north, legend columns=-1}
]
\addplot[fill=popgreen!60] coordinates { (Excellente, 18) (Moyenne, 12) (Insuffisante, 6) };
\addplot[fill=popblue!60] coordinates { (Excellente, 17) (Moyenne, 11) (Insuffisante, 5) };
\legend{Fidélité Sens, Qualité Langue/Connecteurs}
\end{axis}
\end{tikzpicture}
\legende{Notes moyennes par profil de copie sur les critères « fidélité au sens » et « qualité de la langue / connecteurs » (classe de Terminale technique, session d'entraînement). On observe que les deux critères évoluent ensemble : qui lit bien écrit bien.}
\end{popfigure}

\courssub{Remédiation ciblée : des ateliers différenciés}

\begin{definition}[Remédiation pédagogique]
La \cle{remédiation} est un ensemble d'activités proposées après une évaluation pour aider chaque élève à surmonter ses difficultés précises. Elle est \cle{différenciée} lorsque les activités ne sont pas les mêmes pour tous : chacun travaille sur \emph{son} point faible, identifié grâce à la grille critériée.
\end{definition}

Concrètement, la classe est répartie en trois groupes de besoins, constitués d'après la correction des copies.

\begin{center}
\begin{tabularx}{\linewidth}{|l|X|X|}
\hline
\rowcolor{popblueL} \textbf{Groupe} & \textbf{Difficulté diagnostiquée} & \textbf{Atelier de remédiation} \\
\hline
Groupe A & Lecture : l'idée force n'est pas repérée, d'où contresens & Entraînement au surlignage (thème en jaune, thèse en vert, arguments en bleu) puis construction du plan du texte en 3 ou 4 étapes \\
\hline
Groupe B & Réduction : taux non respecté, recopie & Exercices de suppression (exemples, répétitions, détails) et de généralisation (« mangues, oranges, bananes » $\rightarrow$ « fruits ») \\
\hline
Groupe C & Langue : connecteurs pauvres ou fautifs, modes mal employés & Exercices de transformation de phrases ; subjonctif après \emph{bien que, pour que, il faut que} ; réemploi des connecteurs dans des phrases personnelles \\
\hline
\end{tabularx}
\end{center}

\begin{exoresolu}[Exercice de remédiation — Groupe C (connecteurs et modes)]
Énoncé : Complétez avec la conjonction qui convient et mettez le verbe au bon mode : « Le gouvernement encourage la lecture, \dots\ les jeunes (préférer) les écrans. » Puis : « Il faut que chaque élève (lire) au moins un livre par trimestre. »
\tcblower
Solution détaillée : 1) La relation logique est une opposition : on choisit \emph{alors que} ou \emph{tandis que}, qui se construisent avec l'indicatif (constat réel) : « Le gouvernement encourage la lecture, \textbf{alors que} les jeunes \textbf{préfèrent} les écrans. » 2) Après \emph{il faut que}, le subjonctif est obligatoire : « Il faut que chaque élève \textbf{lise} au moins un livre par trimestre. » Piège : ne pas confondre \emph{alors que} (indicatif, opposition constatée) et \emph{bien que} (subjonctif, concession).
\end{exoresolu}

\courssub{Stratégies de révision pour le Bac}

\begin{definition}[Annales]
Les \cle{annales} sont les sujets d'examen des années antérieures, souvent accompagnés de corrigés. S'entraîner sur les annales 2020--2024 permet de se familiariser avec le format réel de l'épreuve, ses exigences et son barème.
\end{definition}

Votre fiche \textbf{« Mémo Contraction »} doit tenir sur une page et contenir trois blocs : les \textbf{étapes} (lire deux fois $\rightarrow$ dégager thème et thèse $\rightarrow$ repérer les arguments $\rightarrow$ rédiger avec ses propres mots $\rightarrow$ compter et vérifier le taux $\rightarrow$ relire) ; les \textbf{connecteurs stars} (\emph{d'abord, ensuite, en effet, cependant, bien que + subjonctif, par conséquent, c'est pourquoi, enfin}) ; les \textbf{pièges} (recopie, contresens, taux dépassé, avis personnel ajouté).

\begin{methode}[Réviser efficacement la contraction en 1 semaine]
\begin{enumerate}
\item \textbf{J1} : Relire le mémo de méthode + apprendre 3 connecteurs par jour avec une phrase d'exemple.
\item \textbf{J2} : Faire une contraction chronométrée (2 h) sur un sujet d'annales.
\item \textbf{J3} : Corriger sa copie avec la grille officielle, en s'auto-notant sévèrement.
\item \textbf{J4} : Atelier connecteurs : 10 phrases à transformer avec le subjonctif.
\item \textbf{J5} : Refaire la contraction du J2 en corrigeant toutes les erreurs repérées.
\item \textbf{J6} : Repos ou lecture plaisir (le cerveau consolide).
\item \textbf{J7} : Simulation complète de Bac blanc dans les conditions réelles.
\end{enumerate}
\end{methode}

\begin{exemplebox}[Statistiques Bac technique 2023--2024]
La moyenne nationale à l'exercice de contraction tourne autour de \num{9}/20, avec un écart-type fort : les copies sont très inégales. Observation décisive des jurys : les candidats qui maîtrisent les connecteurs de subordination (concession : \emph{bien que} ; cause : \emph{parce que, puisque}) gagnent en moyenne \textbf{3 points de plus} sur la note de langue. Trois points, c'est souvent la différence entre l'échec et la mention.
\end{exemplebox}

\courssub{Lien avec l'orientation : une compétence professionnelle}

La contraction n'est pas un exercice scolaire artificiel : c'est la \cle{note de synthèse} du monde professionnel. En BTS comme à l'Université, vous devrez rédiger des \textbf{rapports de stage} condensant des semaines d'observation ; dans la fonction publique et les entreprises, on rédige des \textbf{notes de service} et des comptes rendus qui doivent aller à l'essentiel. Aux concours administratifs camerounais (ENAM, ENSET, concours des Impôts, des Douanes), la note de synthèse est une épreuve reine : on vous donne un dossier de 20 pages et on attend une synthèse d'une page, fidèle, claire, bien articulée. Exactement ce que vous venez d'apprendre.

\begin{savaistu}[Une compétence transversale]
La compétence « synthétiser une information complexe » est transversale : elle sert en Géographie (croquis, commentaire de carte), en Philosophie (analyse de texte), en Économie (note de synthèse) et dans tout métier de cadre. Selon les enquêtes d'insertion professionnelle, la capacité à rédiger une synthèse claire figure parmi les compétences les plus recherchées par les employeurs camerounais, dans la fonction publique comme dans les entreprises privées.
\end{savaistu}

\courssub{Clôture et motivation finale}

Vous voilà au terme de cette leçon. Retenez l'essentiel : résumer, ce n'est pas « raccourcir » un texte, c'est en \cle{reformuler la pensée} avec fidélité, économie et une langue soignée. Cette compétence vous servira au Baccalauréat technique dans quelques mois, puis pendant toute votre vie d'étudiant et de professionnel. Chaque fois que vous lirez un document long et que vous devrez en dire l'essentiel en quelques lignes — à un examinateur, à un chef de service, à un jury —, vous appliquerez la méthode apprise ici. Travaillez vos connecteurs, entraînez-vous sur les annales, relisez toujours vos copies : la réussite au Bac se prépare ainsi, méthodiquement, une contraction après l'autre.

\begin{aretenir}[L'essentiel de la section]
\begin{itemize}
\item L'évaluation formative prépare, l'évaluation sommative certifie ; la grille critériée rend la correction objective.
\item Les questions de linguistique rapportent des points faciles : ne les négligez jamais.
\item La remédiation différenciée cible trois difficultés : lecture (groupe A), réduction (groupe B), langue et connecteurs (groupe C).
\item Une semaine de révision méthodique (mémo, annales, auto-correction, Bac blanc) suffit à progresser nettement.
\item La synthèse est la compétence n°1 des concours administratifs et du monde du travail au Cameroun.
\end{itemize}
\end{aretenir}

\newpage

\coursec{Activité d'intégration}

\coursec{Leçon 25 : Formuler les idées essentielles d'un texte à résumer}

\courssub{Objectifs de l'activité}

Cette activité d'intégration mobilise l'ensemble des savoirs de la leçon dans une situation professionnelle simulée. À l'issue de la mission, l'élève doit être capable de :

\begin{itemize}
\item repérer la \cle{structure argumentative} d'un texte administratif (thèse, arguments, exemples, conclusion) ;
\item appliquer les \cle{techniques de contraction} (suppression, généralisation, reformulation) pour produire une note de synthèse respectant la contrainte du tiers du volume initial ;
\item \cle{reformuler} les idées essentielles sans copier les formulations du texte de départ ;
\item identifier et justifier l'emploi des \cle{conjonctions de subordination} (cause, concession, but, condition) dans un texte ;
\item distinguer les \cle{contenus manifestes} des \cle{contenus latents} d'un discours administratif ;
\item préparer et présenter un \cle{exposé oral} bref à l'aide de fiches bristol.
\end{itemize}

\courssub{Situation}

\textbf{Mission « Synthèse Administrative » : le rapport du Délégué du Gouvernement.}

Vous êtes stagiaire au Secrétariat Général de la Région du Littoral, à Douala. Ce matin, à 8 h 30, le Délégué du Gouvernement vous remet un rapport de trois pages (environ 450 mots) rédigé par le service de l'hygiène et du salubrité. Son titre : \emph{« La gestion des déchets plastiques à Douala : enjeux et solutions »}. Le Délégué doit partir en réunion de coordination à 10 h précises. Il vous confie la mission suivante : « Je n'aurai pas le temps de lire ce document. Résumez-le-moi de manière fiable et préparez-moi de quoi le présenter oralement. »

\textbf{Extrait du rapport (début du texte à contracter) :}

\begin{quote}
« La ville de Douala produit chaque jour près de 1 500 tonnes de déchets, dont environ 12 \% sont des plastiques. Bien que des arrêtés municipaux interdisent la distribution gratuite des sachets plastiques depuis 2014, leur usage demeure massif, parce qu'ils restent bon marché et pratiques pour les commerçants comme pour les consommateurs. Ces déchets bouchent les caniveaux, provoquant des inondations récurrentes dans les quartiers de Ndogpassi, Bépanda et Deïdo, et contaminent le fleuve Wouri. Afin d'enrayer ce phénomène, trois pistes sont proposées : le renforcement du contrôle des arrêtés, le développement de filières de recyclage porteuses d'emplois pour les jeunes, et une campagne de sensibilisation durable auprès des ménages. Si ces mesures sont appliquées avec rigueur, Douala pourrait réduire de moitié ses déchets plastiques d'ici cinq ans, bien que le coût de mise en œuvre, estimé à 2 milliards de FCFA, constitue un frein réel. »
\end{quote}

Le rapport complet comporte trois parties : I. Un constat alarmant ; II. Des causes multiples ; III. Des solutions réalistes mais coûteuses.

\courssub{Consignes}

\begin{enumerate}
\item \textbf{Lecture et décryptage.} Lisez le rapport en entier. Dégagez la thèse défendue, les trois arguments principaux et les exemples qui les illustrent. Présentez le résultat dans un tableau (argument / exemple / fonction).
\item \textbf{Contenus manifestes et latents.} Relevez deux informations explicitement données (contenu manifeste) et une information implicite (contenu latent) que le Délégué doit comprendre sans qu'elle soit écrite.
\item \textbf{Note de synthèse écrite.} Rédigez une note de synthèse de \textbf{150 mots maximum} (soit environ le tiers du rapport complet de 450 mots), en respectant le plan du texte (I, II, III), sans phrase copiée, sans avis personnel, et en employant au moins trois connecteurs logiques variés.
\item \textbf{Étude grammaticale.} Identifiez dans l'extrait ci-dessus \textbf{cinq conjonctions ou locutions conjonctives de subordination} et justifiez leur emploi : une exprimant la cause, deux la concession, une le but, une la condition. Précisez pour chacune le mode du verbe de la subordonnée.
\item \textbf{Exposé oral.} Préparez sur fiches bristol un exposé de \textbf{3 minutes} destiné au chef de cabinet : introduction (sujet et plan), développement (constat, causes, solutions), conclusion. Les fiches ne doivent contenir que des mots-clés, jamais de phrases rédigées.
\item \textbf{Travail en binôme et relecture croisée.} L'un rédige la note, l'autre prépare l'oral ; échangez, relisez, corrigez, puis remettez la production finale. L'évaluation porte sur la grille officielle du Bac : Contenu, Structure, Langue, Oral.
\end{enumerate}

\courssub{Ressources à mobiliser}

\begin{aretenir}[Boîte à outils de la contraction]
\begin{itemize}
\item \textbf{Structure argumentative :} thèse $\rightarrow$ arguments $\rightarrow$ exemples $\rightarrow$ conclusion. Le résumé conserve ce squelette et élimine le reste.
\item \textbf{Techniques de réduction :} suppression (exemples, répétitions, détails chiffrés secondaires), généralisation (une série de termes remplacée par un terme générique), reformulation (paraphrase avec ses propres mots).
\item \textbf{Règle du tiers :} un texte de 450 mots se contracte en environ 150 mots.
\item \textbf{Conjonctions utiles :} coordination (\emph{mais, or, et, donc, ni, car}) ; subordination de cause (\emph{parce que, puisque, comme}), de concession (\emph{bien que, quoique} + subjonctif), de but (\emph{pour que, afin que} + subjonctif ; \emph{pour, afin de} + infinitif), de condition (\emph{si} + indicatif).
\item \textbf{Exposé oral :} introduction (accroche, sujet, plan), développement ordonné, conclusion (bilan, ouverture) ; regard vers l'auditoire, voix posée, fiches bristol en support.
\end{itemize}
\end{aretenir}

\begin{attention}[Pièges à éviter]
Ne copiez jamais une phrase entière du texte : c'est une faute de contraction. N'introduisez aucune opinion personnelle (« je pense que... »). Ne confondez pas \emph{car} (coordination, jamais en tête de phrase) et \emph{parce que} (subordination). Attention au subjonctif obligatoire après \emph{bien que} et \emph{afin que}. Respectez la contrainte de longueur (tiers du texte source).
\end{attention}

\begin{methode}[Démarche pas à pas pour la contraction]
\begin{enumerate}
\item \textbf{Lecture active :} identifier le thème, la thèse, le plan (arguments principaux), les connecteurs logiques.
\item \textbf{Analyse structurelle :} construire le tableau « Argument / Exemple / Fonction ».
\item \textbf{Sélection / Suppression :} barrer les illustrations détaillées (noms de quartiers, chiffres précis non essentiels), les répétitions, les reformulations internes.
\item \textbf{Généralisation :} remplacer les énumérations par des hyperonymes (ex. : « Ndogpassi, Bépanda, Deïdo » $\rightarrow$ « plusieurs quartiers »).
\item \textbf{Reformulation :} réécrire les idées retenues avec son propre vocabulaire, en changeant la syntaxe (nominalisation, verbalisation, changement de voix).
\item \textbf{Assemblage et liaison :} relier les phrases reformulées par des connecteurs logiques adaptés à la relation d'idées (cause, conséquence, concession, énumération).
\item \textbf{Contrôle :} compter les mots (cible $\pm$ 10 \%), vérifier la fidélité au sens, l'absence de copier-coller, la correction linguistique.
\end{enumerate}
\end{methode}

\courssub{Démarche et corrigé commenté}

\begin{exoresolu}[Note de synthèse du rapport complet (450 mots $\rightarrow$ $\sim$150 mots) et analyse grammaticale]
Énoncé : à partir de la structure du rapport (trois parties) et de l'extrait fourni, (a) produire une note de synthèse d'environ 150 mots, (b) relever cinq subordonnantes dans l'extrait et justifier l'emploi de leurs conjonctions.

\tcblower

\textbf{Étape 1 — Repérage de la structure du rapport complet.}
\begin{itemize}
\item \textbf{Thèse :} La gestion des déchets plastiques à Douala est une urgence environnementale et urbaine ; des solutions existent mais exigent une volonté politique et financière forte.
\item \textbf{Partie I (Constat) :} Production massive (1 500 t/j, 12 \% plastiques), pollution visible (caniveaux bouchés, inondations à Ndogpassi, Bépanda, Deïdo, contamination du Wouri).
\item \textbf{Partie II (Causes) :} Persistance des sachets interdits (bas prix, praticité), insuffisance du contrôle municipal, absence d'alternatives accessibles.
\item \textbf{Partie III (Solutions) :} Triple levier : contrôle rigoureux des arrêtés, filières de recyclage (emplois jeunes), sensibilisation durable. Objectif : -50 \% en 5 ans. Frein : coût (2 milliards FCFA).
\end{itemize}

\textbf{Étape 2 — Application des techniques de réduction.}
Suppression des noms de quartiers (détail localisable), généralisation des chiffres (1 500 t $\rightarrow$ « masse considérable », 12 \% $\rightarrow$ « part significative »), reformulation des causes et solutions.

\textbf{Étape 3 — Note de synthèse (environ 150 mots).}

\begin{quote}
\textbf{Note de synthèse : Gestion des déchets plastiques à Douala — Enjeux et solutions}

Douala génère quotidiennement une masse considérable de déchets, dont une part significative de plastiques. Malgré l'interdiction municipale de 2014, les sachets persistent car ils sont économiques et pratiques. Cette prolifération engendre des conséquences graves : obstruction des réseaux d'assainissement, inondations récurrentes dans plusieurs quartiers et pollution du fleuve Wouri.

Face à ce constat, le rapport identifie deux causes majeures : la faible application des textes réglementaires et l'absence d'alternatives viables pour les acteurs économiques.

Pour y remédier, trois axes d'action sont préconisés. Premièrement, le renforcement effectif du contrôle des arrêtés existants. Deuxièmement, la création de filières de recyclage, sources d'emplois pour la jeunesse. Troisièmement, une campagne de sensibilisation continue des ménages. Si cette stratégie est mise en œuvre avec rigueur, l'objectif d'une réduction de 50 \% des déchets plastiques à l'horizon 2029 est atteignable. Toutefois, l'enveloppe prévisionnelle de 2 milliards de FCFA nécessite un engagement budgétaire fort et une gouvernance transparente pour lever le frein financier.
\end{quote}
(Compte : $\sim$ 148 mots. Contrainte respectée. Plan I-II-III conservé. Connecteurs : \emph{Malgré, car, engendre, Face à, Premièrement/Deuxièmement/Troisièmement, Si, Toutefois}.)

\textbf{Étape 4 — Étude grammaticale des subordonnées (sur l'extrait fourni).}

\begin{center}
\begin{tabularx}{\linewidth}{|l|X|X|X|}\hline
\rowcolor{popblueL} \textbf{Conjonction / Locution} & \textbf{Relation logique} & \textbf{Extrait concerné} & \textbf{Mode du verbe} \\ \hline
\emph{parce que} & Cause & « \emph{parce qu'ils restent bon marché} » & Indicatif présent (\emph{restent}) \\ \hline
\emph{bien que} (1) & Concession & « \emph{Bien que des arrêtés municipaux interdisent} » & Subjonctif présent (\emph{interdisent} -- forme homographe à l'indicatif) \\ \hline
\emph{bien que} (2) & Concession & « \emph{bien que le coût \dots constitue un frein} » & Subjonctif présent (\emph{constitue} -- forme homographe à l'indicatif) \\ \hline
\emph{afin de} & But & « \emph{Afin d'enrayer ce phénomène} » & Infinitif (\emph{enrayer}) \\ \hline
\emph{si} & Condition & « \emph{Si ces mesures sont appliquées} » & Indicatif présent (\emph{sont appliquées}) \\ \hline
\end{tabularx}
\end{center}

\textbf{Étape 5 — Contenus manifestes et latents.}
\begin{itemize}
\item \textbf{Manifestes (explicites) :} 1) Volume quotidien de déchets (1 500 t) et part plastique (12 \%). 2) Coût estimé du plan d'action (2 milliards de FCFA).
\item \textbf{Latent (implicite) :} L'inefficacité actuelle de la réglementation (arrêtés de 2014 non appliqués) révèle un déficit de gouvernance locale ; la réussite du plan dépend moins de la technique que de la volonté politique de financer et de contrôler.
\end{itemize}
\end{exoresolu}

\begin{exemplebox}[Fiches bristol modèles pour l'exposé oral (3 min)]
\textbf{Fiche 1 (Intro) :} Titre : \emph{Déchets plastiques Douala}. Accroche : 1 500 t/j. Sujet : Enjeux / Solutions. Plan : 1. Constat 2. Causes 3. Solutions.
\textbf{Fiche 2 (Constat) :} 12 \% plastiques. Conséquences : caniveaux bouchés $\rightarrow$ inondations (Ndogpassi, etc.) + pollution Wouri.
\textbf{Fiche 3 (Causes) :} Sachets interdits mais utilisés (prix, usage). Contrôle faible. Pas d'alternatives.
\textbf{Fiche 4 (Solutions) :} 1. Contrôle rigoureux. 2. Recyclage = emplois jeunes. 3. Sensibilisation ménages. Objectif : -50 \% (5 ans).
\textbf{Fiche 5 (Conclusion) :} Frein : 2 Mds FCFA. Condition : volonté politique + gouvernance. Ouverture : modèle pour autres métropoles.
\end{exemplebox}

\courssub{Bilan de l'activité}

Cette mission a permis de mettre en évidence que la contraction n'est pas un simple raccourcissement : elle suppose d'abord une \cle{compréhension fine} de la structure argumentative, puis une \cle{reformulation fidèle} qui respecte la pensée de l'auteur sans la trahir ni la copier. L'étude des conjonctions de subordination a montré que les connecteurs ne sont pas décoratifs : ils portent la logique du texte (cause, concession, but, condition) et doivent être conservés ou remplacés avec soin dans le résumé. Enfin, la préparation de l'exposé oral a révélé qu'un même contenu peut prendre deux formes — note écrite et parole publique — selon le destinataire et la situation de communication. Ces compétences, directement évaluées au Probatoire et au Baccalauréat technique (épreuve de contraction de texte et exposé), sont aussi celles du monde professionnel : rédiger une note de synthèse claire et la présenter avec aisance est un savoir-faire attendu dans toute administration camerounaise.

\newpage

\coursec{Méthodes et exercices résolus}

\coursec{Formuler les idées essentielles d'un texte à résumer}

\begin{savaistu}[Contexte littéraire : Ferdinand Oyono et \emph{Le Vieux nègre et la médaille}]
\emph{Le Vieux nègre et la médaille} (1956) est le premier roman de \cle{Ferdinand Oyono} (1929--2010), figure majeure de la littérature camerounaise et africaine d'expression française. Diplômé de l'École normale de Dschang, administrateur civil, puis ministre et ambassadeur, Oyono y dépeint avec ironie et lucidité l'aliénation coloniale. L'œuvre, au programme de la \cle{Terminale technique} (Lecture méthodique n° 4 et 5), met en scène \cle{Meka}, vieillard naïf qui voit dans la médaille coloniale une reconnaissance, avant de découvrir l'humiliation masquée. Ce texte est un support privilégié pour travailler la \cle{contraction de texte}, l'\cle{argumentation} et l'analyse des \cle{contenus manifestes et latents}.
\end{savaistu}

\courssub{Méthode 1 : Appliquer les notions de l'unité — repérer la structure argumentative d'un texte}

\begin{methode}[Repérer la structure argumentative d'un texte]
Pour résumer un texte argumentatif, il faut d'abord en dégager l'ossature. Démarche en 5 étapes :
\begin{enumerate}
\item \textbf{Identifier le \cle{thème}} : de quoi parle le texte ? Formulez-le en un nom ou un groupe nominal (ex. : « la colonisation », « l'éducation des filles »).
\item \textbf{Dégager la \cle{thèse}} : quelle position l'auteur défend-il ? Cherchez les marques d'opinion (\emph{je pense que, il faut, il est évident que}) et la conclusion du texte.
\item \textbf{Repérer les \cle{arguments}} : chaque argument est une raison qui soutient la thèse. Ils sont souvent annoncés par des \cle{connecteurs logiques} (\emph{d'abord, ensuite, enfin, en effet, car}).
\item \textbf{Relever les \cle{exemples}} : ils illustrent les arguments mais ne sont pas essentiels ; ils seront supprimés ou réduits dans le résumé.
\item \textbf{Tracer le \cle{schéma}} : Thème $\rightarrow$ Thèse $\rightarrow$ Argument 1 $\rightarrow$ Argument 2 $\rightarrow$ Conclusion. Ce schéma est le squelette du résumé.
\end{enumerate}
\textbf{Réflexe} : soulignez au brouillon les connecteurs logiques ; ils balisent la progression de l'argumentation.
\end{methode}

\begin{exoresolu}[Analyse de la structure d'un texte argumentatif]
\textbf{Énoncé.} Lisez le texte suivant, puis : 1) dégagez le thème et la thèse ; 2) relevez les arguments en les numérotant.

\emph{« L'école doit former des citoyens et non des machines à réciter. En effet, l'élève qui apprend par cœur sans comprendre oublie tout dès l'examen passé. De plus, la vie professionnelle exige des personnes capables de résoudre des problèmes nouveaux, et non de répéter des leçons. C'est pourquoi il faut privilégier, dans nos classes, le raisonnement sur la mémoire. »}

\tcblower
\textbf{Solution détaillée.}

\textbf{1) Thème et thèse.}
\begin{itemize}
\item \textbf{Thème} : l'éducation, plus précisément la méthode d'enseignement (mémoire contre raisonnement).
\item \textbf{Thèse} : l'école doit privilégier le raisonnement plutôt que l'apprentissage par cœur. On la repère grâce à la formule conclusive « C'est pourquoi il faut privilégier... le raisonnement sur la mémoire » et à la phrase d'ouverture « former des citoyens et non des machines à réciter ».
\end{itemize}

\textbf{2) Arguments.}
\begin{itemize}
\item \textbf{Argument 1} (introduit par \emph{En effet}) : le savoir mémorisé sans compréhension est vite oublié.
\item \textbf{Argument 2} (introduit par \emph{De plus}) : la vie professionnelle demande des capacités d'adaptation et de résolution de problèmes, pas la répétition.
\end{itemize}

\textbf{Schéma de la structure} :
\begin{center}
\begin{tabularx}{\linewidth}{|l|X|}\hline
\rowcolor{popblueL} Élément & Contenu \\ \hline
Thème & La méthode d'enseignement \\ \hline
Thèse & Il faut privilégier le raisonnement sur la mémoire \\ \hline
Argument 1 & Le par cœur sans compréhension est vite oublié \\ \hline
Argument 2 & La vie professionnelle exige de résoudre des problèmes nouveaux \\ \hline
Conclusion & « C'est pourquoi il faut privilégier... » (rappel de la thèse) \\ \hline
\end{tabularx}
\end{center}

\textbf{Conclusion} : le texte défend une pédagogie fondée sur la compréhension, appuyée sur deux arguments (fragilité du par cœur, exigences de la vie active) et conclue par un connecteur de conséquence.
\end{exoresolu}

\courssub{Méthode 2 : Appliquer les notions de l'unité — contracter un texte (technique de réduction)}

\begin{methode}[Réduire un texte au quart de sa longueur]
Au Baccalauréat technique, le résumé demande généralement une \cle{contraction} à un quart ($1/4$) du texte, avec une marge de 10\,\%. Démarche :
\begin{enumerate}
\item \textbf{Compter les mots} du texte (ou estimer : nombre de mots d'une ligne $\times$ nombre de lignes). Diviser par 4 pour obtenir la longueur cible.
\item \textbf{Supprimer} : les exemples, les répétitions, les citations, les détails chiffrés secondaires, les énumérations (à remplacer par un terme générique), les figures de style (comparaisons, métaphores).
\item \textbf{Fusionner} : remplacer plusieurs phrases par une seule (subordination, \cle{nominalisation} : « il faut que les élèves comprennent » $\rightarrow$ « la compréhension des élèves »).
\item \textbf{Reformuler} avec ses propres mots, à la troisième personne, sans « je », sans citer le texte entre guillemets.
\item \textbf{Lier} les idées par des connecteurs logiques (\emph{d'abord, ensuite, enfin, car, c'est pourquoi}) et vérifier le nombre de mots final.
\end{enumerate}
\textbf{Formule à retenir} : $\text{longueur du résumé} = \dfrac{\text{nombre de mots du texte}}{4}$, tolérance de $\pm 10\,\%$.
\end{methode}

\begin{exoresolu}[Contraction d'un paragraphe]
\textbf{Énoncé.} Résumez le texte suivant (80 mots) en 20 mots environ.

\emph{« Le Vieux nègre et la médaille montre, à travers le personnage de Meka, un vieil homme naïf et crédule, les illusions et les humiliations que la colonisation a fait subir aux Africains, notamment lors de la cérémonie de remise de la médaille, véritable mascarade organisée par l'administration coloniale. »}

\tcblower
\textbf{Solution détaillée.}

\textbf{Étape 1 — Calcul.} Texte : environ 50 mots de contenu utile sur 80 ; la consigne fixe la cible : 20 mots (soit le quart de 80).

\textbf{Étape 2 — Idées essentielles à conserver.} (a) l'œuvre dénonce la colonisation ; (b) Meka incarne la naïveté de la victime ; (c) la médaille symbolise l'humiliation.

\textbf{Étape 3 — À supprimer.} Les adjectifs descriptifs (« naïf et crédule » $\rightarrow$ « naïf » suffit), la précision « véritable mascarade organisée par l'administration coloniale » (détail déjà contenu dans « humiliations »).

\textbf{Étape 4 — Reformulation.}

\emph{« À travers Meka, vieillard naïf, le roman dénonce les illusions et les humiliations infligées aux Africains par la colonisation. »}

\textbf{Étape 5 — Vérification.} Comptage : À(1) travers(2) Meka(3), vieillard(4) naïf(5), le(6) roman(7) dénonce(8) les(9) illusions(10) et(11) les(12) humiliations(13) infligées(14) aux(15) Africains(16) par(17) la(18) colonisation(19). = 19 mots, conforme à la consigne ($20 \pm 10\,\%$).

\textbf{Conclusion} : le résumé conserve la thèse (dénonciation de la colonisation), l'argument (le personnage de Meka) et supprime tous les détails, en 19 mots.
\end{exoresolu}

\courssub{Méthode 3 : Appliquer les notions de l'unité — reformuler les idées essentielles}

\begin{methode}[Reformuler sans trahir le texte]
Reformuler, c'est dire la même chose autrement. Quatre techniques :
\begin{enumerate}
\item \textbf{La \cle{nominalisation}} : transformer un verbe en nom. « Les colons exploitent les terres » $\rightarrow$ « l'exploitation des terres par les colons ».
\item \textbf{Le \cle{changement de voix}} : actif $\leftrightarrow$ passif. « L'administration humilie Meka » $\rightarrow$ « Meka est humilié par l'administration ».
\item \textbf{La \cle{substitution lexicale}} : remplacer par un synonyme ou un terme générique. « houe, machette, pirogue » $\rightarrow$ « outils » ; « tromper » $\rightarrow$ « duper ».
\item \textbf{La \cle{condensation syntaxique}} : deux phrases $\rightarrow$ une phrase avec subordonnée. « Meka reçoit une médaille. Cette médaille est une duperie. » $\rightarrow$ « La médaille que reçoit Meka est une duperie. »
\end{enumerate}
\textbf{Interdits} : copier des phrases entières du texte, ajouter son avis personnel, utiliser « je ».
\end{methode}

\begin{exoresolu}[Reformulation de phrases]
\textbf{Énoncé.} Reformulez chaque phrase en utilisant la technique indiquée, sans changer le sens.

1. « Les élèves doivent comprendre les leçons. » (nominalisation)

2. « La cérémonie de la médaille révèle la duplicité des colons. » (voix passive)

3. « Meka est joyeux, fier et plein d'espoir avant la cérémonie. » (condensation + terme générique)

\tcblower
\textbf{Solution détaillée.}

\textbf{1. Nominalisation.} Le verbe « comprendre » devient le nom « la compréhension » ; « doivent » devient « est nécessaire » :
\emph{« La compréhension des leçons est nécessaire aux élèves. »}
Le sens est identique ; seule la construction change.

\textbf{2. Voix passive.} Le complément d'objet direct (« la duplicité des colons ») devient sujet ; le sujet initial devient complément d'agent introduit par « par » :
\emph{« La duplicité des colons est révélée par la cérémonie de la médaille. »}
Vérification : l'auxiliaire \emph{être} + participe passé accordé avec le nouveau sujet (« révélée », féminin singulier).

\textbf{3. Condensation + terme générique.} Les trois adjectifs « joyeux, fier et plein d'espoir » expriment tous un sentiment positif ; on les remplace par un terme générique :
\emph{« Avant la cérémonie, Meka déborde d'optimisme. »}
On condense l'énumération en un nom unique, ce qui raccourcit la phrase sans perte de sens.

\textbf{Conclusion} : chaque reformulation préserve le sens, change la formulation et évite la copie du texte — c'est exactement ce qu'exige l'épreuve de contraction.
\end{exoresolu}

\courssub{Méthode 4 : Appliquer les notions de l'unité — employer les conjonctions de coordination et de subordination}

\begin{methode}[Choisir la bonne conjonction pour lier les idées]
\begin{enumerate}
\item \textbf{Distinguer les deux catégories.}
\begin{itemize}
\item \textbf{\cle{Coordination}} : relie deux éléments de même fonction. Les 7 conjonctions : \emph{mais, ou, et, donc, or, ni, car} (moyen mnémotechnique : « Mais où est donc Ornicar ? »).
\item \textbf{\cle{Subordination}} : introduit une proposition dépendante. \emph{que, parce que, puisque, bien que, quoique, si, lorsque, afin que, pour que}...
\end{itemize}
\item \textbf{Choisir selon le \cle{rapport logique}} : cause (\emph{parce que, car, puisque}), conséquence (\emph{donc, c'est pourquoi, si bien que}), opposition (\emph{mais, bien que, alors que}), but (\emph{pour que, afin de}), temps (\emph{quand, lorsque, dès que}), condition (\emph{si, à condition que}).
\item \textbf{Surveiller le \cle{mode}} : après \emph{bien que, quoique, afin que, pour que, à condition que} $\rightarrow$ \textbf{subjonctif} ; après \emph{parce que, puisque, lorsque} $\rightarrow$ \textbf{indicatif}.
\item \textbf{Dans le résumé}, utiliser ces conjonctions pour fusionner les phrases et marquer la progression argumentative.
\end{enumerate}
\end{methode}

\begin{exoresolu}[Emploi des conjonctions]
\textbf{Énoncé.} 1) Complétez avec la conjonction qui convient (nature et rapport logique à préciser) : « Meka est fier de sa médaille, \dots\ elle ne représente qu'une duperie. » — « \dots\ il a travaillé toute sa vie pour les colons, il attend une récompense. »

2) Reliez en une seule phrase : « Meka perd sa terre. Meka perd ses illusions. »

\tcblower
\textbf{Solution détaillée.}

\textbf{1. Complétion.}
\begin{itemize}
\item « Meka est fier de sa médaille, \textbf{mais} elle ne représente qu'une duperie. » \emph{Mais} est une conjonction de \textbf{coordination} exprimant l'\textbf{opposition} entre la fierté de Meka et la réalité de la médaille.
\item « \textbf{Puisqu'}il a travaillé toute sa vie pour les colons, il attend une récompense. » \emph{Puisque} est une conjonction de \textbf{subordination} exprimant la \textbf{cause} ; elle est suivie de l'indicatif (« a travaillé »), ce qui est correct.
\end{itemize}

\textbf{2. Fusion en une phrase.} Plusieurs solutions correctes :
\begin{itemize}
\item Coordination : « Meka perd sa terre \textbf{et} ses illusions. » (addition)
\item Subordination : « \textbf{Lorsqu'}il perd sa terre, Meka perd \textbf{aussi} ses illusions. » (simultanéité)
\item Avec valeur de conséquence : « Meka perd sa terre, \textbf{si bien qu'}il perd aussi ses illusions. »
\end{itemize}
La meilleure pour un résumé est la première : la plus courte, donc la plus économique.

\textbf{Conclusion} : le choix de la conjonction dépend du rapport logique (opposition, cause, addition) ; la subordination permet de hiérarchiser les idées, la coordination de les juxtaposer sobrement.
\end{exoresolu}

\courssub{Méthode 5 : Rédiger et justifier une démarche — distinguer contenus manifestes et contenus latents}

\begin{methode}[Dégager le sens manifeste et le sens latent d'un message]
\begin{enumerate}
\item \textbf{\cle{Contenu manifeste}} : ce qui est dit explicitement, le sens littéral, directement observable (les faits, les mots prononcés).
\item \textbf{\cle{Contenu latent}} : ce qui est suggéré, implicite, caché derrière les mots (intentions, ironie, critiques non avouées, non-dits).
\item \textbf{Procédure} : (a) reformuler d'abord le sens littéral ; (b) repérer les indices de l'implicite (ironie, exagération, contraste entre les dires et les faits, silences) ; (c) formuler le sens latent avec des précautions : « l'auteur laisse entendre que... », « derrière cette scène se cache... ».
\item \textbf{Justifier} toujours l'interprétation par un élément précis du texte (un mot, un contraste, une situation).
\end{enumerate}
\textbf{Application type} : dans \emph{Le Vieux nègre et la médaille}, la cérémonie officielle (manifeste : une fête honorifique) cache une humiliation (latent : la médaille récompense la soumission).
\end{methode}

\begin{exoresolu}[Analyse d'un énoncé à double contenu]
\textbf{Énoncé.} Le commandant de cercle déclare lors de la cérémonie : « Meka est un grand ami de la France, un exemple pour tous les indigènes. » Distinguez le contenu manifeste et le contenu latent de cette déclaration, en justifiant votre analyse.

\tcblower
\textbf{Solution détaillée.}

\textbf{Étape 1 — Contenu manifeste.} Littéralement, le commandant fait l'éloge de Meka : il le présente comme un ami de la France et un modèle à suivre. C'est ce qui est dit explicitement, devant l'assistance.

\textbf{Étape 2 — Indices de l'implicite.}
\begin{itemize}
\item Le mot \emph{« indigènes »} : terme méprisant de l'administration coloniale qui contredit le mot \emph{« ami »} ; on ne parle pas ainsi d'un égal.
\item Le contexte : Meka a donné sa terre et travaillé toute sa vie ; la médaille ne coûte rien aux colons. Le contraste entre le sacrifice réel et la récompense symbolique révèle l'hypocrisie.
\item « Un exemple pour tous » : l'éloge sert en réalité à encourager la soumission des autres Africains.
\end{itemize}

\textbf{Étape 3 — Contenu latent.} Derrière l'éloge, l'administration coloniale manipule Meka : la médaille récompense et encourage la docilité, non le mérite. Le discours officiel masque la domination.

\textbf{Étape 4 — Formulation prudente.} « L'auteur laisse entendre, par le contraste entre l'éloge et la réalité de la situation de Meka, que cette cérémonie est une mascarade destinée à entretenir la soumission des colonisés. »

\textbf{Conclusion} : le contenu manifeste est un éloge officiel ; le contenu latent, révélé par le vocabulaire méprisant et le contraste sacrifice/récompense, est une critique de l'hypocrisie coloniale. Toute interprétation du sens latent doit être justifiée par des indices textuels précis.
\end{exoresolu}

\courssub{Méthode 6 : Rédiger et justifier une démarche — préparer et présenter un exposé}

\begin{methode}[Préparer et présenter un exposé oral]
\begin{enumerate}
\item \textbf{Avant} : définir le sujet en une question ; collecter et trier les informations ; construire un plan en 3 parties (introduction, développement en 2 ou 3 points, conclusion) ; rédiger des \cle{fiches} avec des mots-clés (jamais le texte intégral).
\item \textbf{Introduction} : accroche (fait, question), présentation du sujet, annonce du plan.
\item \textbf{Développement} : une idée par partie, annoncée puis expliquée puis illustrée ; transitions explicites (« après avoir vu..., nous allons... »).
\item \textbf{Conclusion} : bilan des idées essentielles, ouverture éventuelle, remerciement.
\item \textbf{Pendant la présentation} : regarder l'auditoire, parler lentement et distinctement, ne pas lire, gérer le temps, utiliser des connecteurs pour guider l'écoute.
\item \textbf{Justifier sa démarche} : être capable d'expliquer le choix du plan et des exemples si le jury interroge.
\end{enumerate}
\end{methode}

\begin{exoresolu}[Plan d'exposé]
\textbf{Énoncé.} Vous devez présenter un exposé de 5 minutes sur : « La dénonciation de la colonisation dans \emph{Le Vieux nègre et la médaille} ». Rédigez l'introduction complète de votre exposé et annoncez le plan.

\tcblower
\textbf{Solution détaillée.}

\textbf{Étape 1 — Accroche.} Partir d'un fait marquant de l'œuvre : « Que vaut une médaille offerte à un homme à qui l'on a tout pris ? »

\textbf{Étape 2 — Présentation du sujet.} « \emph{Le Vieux nègre et la médaille} de Ferdinand Oyono, publié en 1956, raconte l'histoire de Meka, vieillard africain décoré par l'administration coloniale pour sa fidélité. »

\textbf{Étape 3 — Problématique.} « Nous nous demanderons comment Oyono, à travers ce récit, dénonce les mécanismes et les illusions de la colonisation. »

\textbf{Étape 4 — Annonce du plan.} « Nous verrons d'abord que le personnage de Meka incarne la naïveté des colonisés ; ensuite que la cérémonie de la médaille constitue une mascarade révélatrice de l'hypocrisie coloniale ; enfin que l'évolution de Meka, de la fierté à la désillusion, traduit la prise de conscience. »

\textbf{Vérification de la démarche} : le plan suit la progression du roman (avant, pendant, après la cérémonie), chaque partie correspond à un aspect de la dénonciation, et l'ensemble respecte la durée (3 parties $\times$ environ 1 min 30 + introduction et conclusion = 5 minutes).

\textbf{Conclusion} : l'exposé s'ouvre sur une accroche, pose une problématique claire et annonce un plan en trois parties justifié par la chronologie de l'œuvre — la démarche est ainsi explicite et défendable devant le jury.
\end{exoresolu}

\begin{aretenir}[Check-list de la contraction de texte (épreuve du Bac)]
\begin{itemize}
\item \textbf{Longueur} : $1/4$ du texte original $\pm 10\,\%$ (compter les mots au brouillon).
\item \textbf{Structure} : Thème $\rightarrow$ Thèse $\rightarrow$ Arguments principaux (dans l'ordre) $\rightarrow$ Conclusion.
\item \textbf{Style} : Troisième personne, temps du texte (souvent présent de narration ou passé), phrases complètes liées par des connecteurs.
\item \textbf{Reformulation} : Nominalisation, voix passive, substitution lexicale, condensation. \textbf{Zéro citation}, zéro « je ».
\item \textbf{Relecture} : Vérifier le nombre de mots, la cohérence logique, l'orthographe grammaticale (accords, subjonctifs après \emph{bien que, pour que}...).
\end{itemize}
\end{aretenir}

\begin{attention}[Les 5 erreurs qui coûtent des points au Bac]
\begin{enumerate}
\item \textbf{Copier des phrases du texte} dans le résumé au lieu de reformuler : tout passage recopié (même sans guillemets) est sanctionné. Reformulez toujours avec vos propres mots.
\item \textbf{Dépasser ou sous-estimer la longueur imposée} : comptez vos mots et respectez le quart du texte ($\pm 10\,\%$). Un résumé de la moitié du texte n'est pas une contraction.
\item \textbf{Ajouter son opinion personnelle} (« je trouve que... », « à mon avis... ») : le résumé est objectif, à la troisième personne ; il restitue la pensée de l'auteur, pas la vôtre.
\item \textbf{Confondre coordination et subordination} ou employer l'indicatif après \emph{bien que} : rappelez-vous que \emph{bien que, quoique, pour que, afin que} exigent le \textbf{subjonctif}.
\item \textbf{Présenter le sens latent sans le justifier} : toute interprétation de l'implicite doit s'appuyer sur un indice précis du texte (mot, contraste, ironie), sinon elle est gratuite et non créditée.
\end{enumerate}
\end{attention}

\newpage

\coursec{Exercices}

\courssub{Je vérifie mes connaissances}

\begin{exercice}[QCM : les outils de liaison]
Pour chaque question, entoure la bonne réponse.
\begin{enumerate}
\item Dans la phrase « Meka reçut la médaille \textbf{car} il avait donné ses terres », la conjonction \emph{car} exprime :
\begin{enumerate}
\item la conséquence \item la cause \item la concession
\end{enumerate}
\item Laquelle de ces conjonctions est une conjonction de \textbf{subordination} ?
\begin{enumerate}
\item mais \item donc \item bien que
\end{enumerate}
\item Après « bien que », le verbe se met généralement au :
\begin{enumerate}
\item indicatif \item subjonctif \item impératif
\end{enumerate}
\item Un résumé-contraction de 90 mots ± 10 \% peut contenir au maximum :
\begin{enumerate}
\item 90 mots \item 99 mots \item 110 mots
\end{enumerate}
\end{enumerate}
\end{exercice}

\begin{exercice}[Vrai ou faux — justifie ta réponse]
Dis si chaque affirmation est vraie ou fausse, puis justifie en une phrase.
\begin{enumerate}
\item Dans une contraction de texte, on peut ajouter son opinion personnelle pour enrichir le résumé.
\item Les conjonctions de coordination relient deux éléments de même fonction grammaticale.
\item Un contenu latent est une idée explicitement formulée par l'auteur.
\item La nominalisation consiste à transformer un verbe ou une proposition en groupe nominal.
\end{enumerate}
\end{exercice}

\begin{exercice}[Définitions]
Définis en deux ou trois phrases chacun des termes suivants, en donnant pour chacun un exemple tiré de \emph{Le Vieux Nègre et la Médaille} ou d'un texte argumentatif étudié :
\begin{enumerate}
\item la \cle{thèse} d'un texte argumentatif ;
\item la \cle{contraction de texte} ;
\item le \cle{contenu latent} (ou implicite) ;
\item la \cle{conjonction de subordination}.
\end{enumerate}
\end{exercice}

\begin{exercice}[Repérage direct]
Soit la phrase extraite d'un éditorial : « Bien que l'État ait créé des structures d'accompagnement, les jeunes peinent à accéder au financement, car les banques exigent des garanties. »
\begin{enumerate}
\item Relève la conjonction de subordination et précise la relation logique qu'elle exprime.
\item Relève la conjonction de coordination et précise la relation logique qu'elle exprime.
\item Indique le mode du verbe « ait créé » et explique ce choix.
\end{enumerate}
\end{exercice}

\courssub{Je m'entraîne}

\begin{exercice}[Transformation syntaxique]
Transforme les cinq phrases coordonnées suivantes (inspirées de \emph{Le Vieux Nègre et la Médaille}) en phrases complexes comportant une subordonnée. Respecte la concordance des temps et l'emploi du subjonctif quand il est requis. La relation logique attendue est indiquée entre parenthèses.
\begin{enumerate}
\item Meka a donné ses terres à la mission ; il a reçu une médaille. (cause)
\item Meka est vieux et pauvre ; les Blancs le traitent avec mépris. (conséquence)
\item Meka est fier de sa médaille ; il ignore ce qu'elle signifie vraiment. (concession)
\item Les villageois se rassemblent ; ils veulent voir la médaille de Meka. (but)
\item Tu obéiras aux Blancs ; tu seras récompensé. (condition)
\end{enumerate}
\end{exercice}

\begin{exercice}[Nominalisation]
Reformule chacune des phrases suivantes en nominalisant la partie soulignée par le sens (transforme la proposition verbale en groupe nominal). Exemple : « Les jeunes créent des entreprises » $\rightarrow$ « la création d'entreprises par les jeunes ».
\begin{enumerate}
\item Bien que l'État \textbf{ait créé} des structures d'accompagnement, les jeunes peinent à accéder au financement.
\item Quand les jeunes \textbf{entreprennent}, ils créent des emplois.
\item Parce que les banques \textbf{refusent} les prêts, les projets échouent.
\item Si l'on \textbf{forme} les jeunes, le chômage diminuera.
\end{enumerate}
\end{exercice}

\begin{exercice}[Contenus manifestes et latents]
Lis cet extrait adapté du discours du Chef de Poste dans \emph{Le Vieux Nègre et la Médaille} : « Ce brave Meka ! Voilà un homme qui a su comprendre les bienfaits de la civilisation. Sa médaille est la preuve que la France sait reconnaître ses amis. »
\begin{enumerate}
\item Dégage le contenu manifeste (ce qui est dit explicitement) en une phrase.
\item Relève \textbf{trois} contenus latents (implicites) véhiculés par ce discours.
\item Pour chaque implicite, indique l'indice linguistique qui le trahit : ironie, euphémisme, modalisation, choix du vocabulaire.
\end{enumerate}
\end{exercice}

\begin{exercice}[Réduction d'un paragraphe]
Voici un paragraphe de 80 mots : « L'entrepreneuriat est souvent présenté comme la solution miracle au chômage des jeunes au Cameroun. Pourtant, force est de constater que les jeunes entrepreneurs se heurtent à de nombreux obstacles : l'accès difficile au crédit bancaire, la lourdeur des procédures administratives, l'absence de formation adaptée. Malgré ces difficultés, certains réussissent et prouvent que la détermination peut triompher des obstacles. Encore faut-il que l'État accompagne réellement ces volontés. »
\begin{enumerate}
\item Souligne les idées essentielles du paragraphe (trois au maximum).
\item Rédige une contraction de ce paragraphe en 25 mots ± 10 \%.
\item Indique le nombre exact de mots de ta contraction.
\end{enumerate}
\end{exercice}

\courssub{Je me prépare au Bac}

\begin{exercice}[Sujet type Bac Technique : contraction d'un éditorial]
\textbf{Texte} (extrait adapté d'un éditorial du journal \emph{Le Jour}, Yaoundé) : « La jeunesse camerounaise est aujourd'hui sommée d'entreprendre. Dans les discours officiels comme dans les médias, l'entrepreneuriat est érigé en remède universel contre le chômage qui frappe des centaines de milliers de diplômés. Certes, l'initiative est louable : un pays ne peut se développer si sa jeunesse attend tout de l'État. Cependant, on ne peut ignorer les réalités du terrain. Bien que l'État ait créé des structures d'accompagnement et des fonds spéciaux, les jeunes peinent à accéder au financement, car les banques exigent des garanties qu'ils ne possèdent pas. De plus, les procédures administratives demeurent longues et décourageantes. Faut-il pour autant renoncer ? Non. Mais il faut cesser de faire croire aux jeunes que la volonté suffit. Entrepreneur, oui ; mais à condition que les pouvoirs publics créent un environnement réellement favorable. Sinon, l'entrepreneuriat restera un slogan vide, et la jeunesse, désabusée, continuera de rêver de l'émigration. » (environ 160 mots pour l'exercice ; au Bac, le texte complet compte 280 mots)
\begin{enumerate}
\item Quelle est la thèse défendue par l'auteur de cet éditorial ? Justifie ta réponse en relevant deux arguments du texte.
\item Relevez deux conjonctions de subordination marquant la concession ou la condition, et analysez le mode verbal qui les suit.
\item Reformulez la phrase suivante en la nominalisant : « Bien que l'État ait créé des structures d'accompagnement, les jeunes peinent à accéder au financement. »
\item Identifiez dans le texte deux marqueurs de la structure argumentative (connecteurs d'opposition, de conclusion) et précisez leur rôle.
\item Rédigez une contraction du texte en 50 mots ± 10 \% (au Bac : texte de 280 mots $\rightarrow$ 90 mots ± 10 \%). Indiquez le nombre de mots.
\end{enumerate}
\end{exercice}

\begin{exercice}[Analyse textuelle et langue : \emph{Le Vieux Nègre et la Médaille}]
Lis l'extrait suivant (adapté de la lecture méthodique n° 4) : « Meka attendait la médaille comme on attend la pluie après une longue sécheresse. Il avait donné ses terres, il avait donné ses fils à la guerre, et voilà que les Blancs allaient enfin le reconnaître. Le jour de la cérémonie, il mit ses plus beaux habits. Mais quand on lui accrocha la médaille sur la poitrine, il sentit seulement le poids du métal froid. »
\begin{enumerate}
\item Dégagez le contenu manifeste de cet extrait en une phrase, puis relevez deux contenus latents en précisant les indices linguistiques (comparaison, ironie, champ lexical).
\item Relevez dans l'extrait deux conjonctions de coordination et une conjonction de subordination ; classez-les selon la relation logique exprimée.
\item Transformez la phrase « Il avait donné ses terres, il avait donné ses fils à la guerre, et voilà que les Blancs allaient enfin le reconnaître » en une phrase complexe comportant une subordonnée de cause et une subordonnée de conséquence.
\item En vous appuyant sur vos réponses, rédigez un court paragraphe (10 à 15 lignes) montrant comment Ferdinand Oyono utilise l'ironie comme arme critique contre la colonisation dans cet extrait.
\end{enumerate}
\end{exercice}

\begin{exercice}[Épreuve orale blanche : l'exposé]
Dans le cadre de la préparation à l'épreuve orale du Baccalauréat technique, tu présenteras devant la classe un exposé de \textbf{5 minutes} sur le sujet suivant : « L'ironie comme arme anticoloniale chez Ferdinand Oyono dans \emph{Le Vieux Nègre et la Médaille} ».
\begin{enumerate}
\item \textbf{Préparation écrite} (à rendre avant l'exposé) : rédige une fiche de préparation comportant une problématique, un plan en deux ou trois parties annoncé dans une introduction, et une conclusion avec ouverture.
\item \textbf{Support visuel} : conçois un support de 3 diapositives (slides) maximum ; indique pour chacune le titre et les éléments qu'elle contiendra (aucun texte long : mots-clés et exemples uniquement).
\item \textbf{Passation} : ton exposé sera évalué par tes pairs à l'aide de la grille MINESEC ci-dessous ; prépare-toi à répondre à deux questions du jury de classe.
\end{enumerate}
\begin{center}
\begin{tabularx}{\linewidth}{|l|X|c|}
\hline
\rowcolor{popblueL} \textbf{Critère} & \textbf{Attendus} & \textbf{Note} \\
\hline
Contenu & Pertinence des idées, appui sur l'œuvre & /4 \\
\hline
Organisation & Introduction, plan annoncé, transitions, conclusion & /4 \\
\hline
Langue & Correction syntaxique, liaisons logiques, vocabulaire & /4 \\
\hline
Communication & Voix, regard, gestion du temps (5 min), support visuel & /4 \\
\hline
Réponses aux questions & Clarté et justification des réponses & /4 \\
\hline
\end{tabularx}
\end{center}
\end{exercice}

\newpage

\coursec{Corrigés des exercices}

\begin{corrige}[Exercice 1 — QCM : les outils de liaison]
\begin{enumerate}
\item \textbf{Réponse b : la cause.} La conjonction de coordination \emph{car} introduit la raison pour laquelle Meka a reçu la médaille : il avait donné ses terres. \emph{Car} exprime toujours la cause ou la justification.
\item \textbf{Réponse c : bien que.} \emph{Mais} et \emph{donc} sont des conjonctions de coordination ; \emph{bien que} introduit une subordonnée circonstancielle de concession.
\item \textbf{Réponse b : subjonctif.} Après \emph{bien que} (concession), le verbe de la subordonnée se met au subjonctif : « bien qu'il \emph{soit} vieux ».
\item \textbf{Réponse b : 99 mots.} Calcul : 10 \% de 90 = 9 ; le maximum est donc $90 + 9 = 99$ mots (la fourchette va de 81 à 99 mots).
\end{enumerate}
\end{corrige}

\begin{corrige}[Exercice 2 — Vrai ou faux]
\begin{enumerate}
\item \textbf{Faux.} La contraction de texte doit restituer fidèlement la pensée de l'auteur : toute opinion personnelle constitue une faute méthodologique éliminatoire au Bac.
\item \textbf{Vrai.} Les conjonctions de coordination (\emph{mais, ou, et, donc, or, ni, car}) unissent deux éléments de même nature et de même fonction : deux mots, deux groupes ou deux propositions indépendantes.
\item \textbf{Faux.} Le contenu latent est l'idée \emph{sous-entendue}, non formulée explicitement, que le lecteur doit déduire des indices linguistiques ; c'est le contenu manifeste qui est explicite.
\item \textbf{Vrai.} La nominalisation transforme un verbe ou une proposition en groupe nominal : « les jeunes entreprennent » devient « l'entrepreneuriat des jeunes » ; c'est un outil majeur de la contraction.
\end{enumerate}
\end{corrige}

\begin{corrige}[Exercice 3 — Définitions]
\begin{enumerate}
\item \textbf{La thèse} est l'idée principale, le point de vue que l'auteur d'un texte argumentatif défend et cherche à faire admettre au lecteur. Elle est appuyée par des arguments et des exemples. \emph{Exemple :} dans \emph{Le Vieux Nègre et la Médaille}, la thèse implicite d'Oyono est que la colonisation repose sur l'illusion et l'humiliation des colonisés.
\item \textbf{La contraction de texte} est un exercice de réduction qui consiste à restituer, en un nombre de mots imposé (le quart du texte au Bac, ± 10 \%), les idées essentielles d'un texte, sans commentaire ni opinion personnelle. Elle exige reformulation, nominalisation et suppression des exemples. \emph{Exemple :} réduire l'éditorial sur l'entrepreneuriat à 50 mots.
\item \textbf{Le contenu latent} (ou implicite) est ce qu'un énoncé suggère sans le dire ouvertement : sous-entendus, ironie, présupposés. On le repère grâce aux indices linguistiques (choix du vocabulaire, modalisation, comparaisons). \emph{Exemple :} quand le Chef de Poste vante « les bienfaits de la civilisation », il laisse entendre, par ironie chez Oyono, que la colonisation est en réalité une spoliation.
\item \textbf{La conjonction de subordination} est un mot-outil (\emph{que, parce que, bien que, si, quand, pour que}\ldots) qui introduit une proposition subordonnée dépendant d'une proposition principale. Elle précise la relation logique : cause, conséquence, but, concession, condition, temps. \emph{Exemple :} « Meka est fier \textbf{bien qu'}il ignore le sens de sa médaille » (concession).
\end{enumerate}
\end{corrige}

\begin{corrige}[Exercice 4 — Repérage direct]
\begin{enumerate}
\item Conjonction de subordination : \textbf{« bien que »}. Elle exprime la \textbf{concession} (ou opposition) : malgré l'existence des structures, le résultat attendu n'est pas atteint.
\item Conjonction de coordination : \textbf{« car »}. Elle exprime la \textbf{cause} : l'exigence de garanties par les banques explique la difficulté d'accès au financement.
\item « Ait créé » est au \textbf{subjonctif plus-que-parfait}. Ce mode est requis après la conjonction de concession \emph{bien que} ; le plus-que-parfait marque l'antériorité : la création des structures est antérieure à la difficulté actuelle des jeunes.
\end{enumerate}
\end{corrige}

\begin{corrige}[Exercice 5 — Transformation syntaxique]
Corrections possibles (toute formulation respectant la relation logique et les modes est acceptée) :
\begin{enumerate}
\item \textbf{Parce qu'}il avait donné ses terres à la mission, Meka a reçu une médaille. (subordonnée de cause, indicatif plus-que-parfait pour l'antériorité)
\item Meka est si vieux et si pauvre \textbf{que} les Blancs le traitent avec mépris. (subordonnée de conséquence avec \emph{si\ldots que}, indicatif)
\item \textbf{Bien qu'}il soit fier de sa médaille, Meka ignore ce qu'elle signifie vraiment. (concession : \emph{bien que} + subjonctif présent \emph{soit})
\item Les villageois se rassemblent \textbf{pour qu'}on voie la médaille de Meka / \textbf{afin de} voir la médaille de Meka. (but : \emph{pour que} + subjonctif)
\item \textbf{Si} tu obéis aux Blancs, tu seras récompensé. (condition : \emph{si} + présent, futur dans la principale)
\end{enumerate}
\textbf{Points de vigilance :} ne jamais écrire « bien qu'il \emph{est} » ; après \emph{si} exprimant la condition, on n'emploie jamais le futur (« si tu obéiras » est fautif).
\end{corrige}

\begin{corrige}[Exercice 6 — Nominalisation]
\begin{enumerate}
\item « Bien que l'État ait créé des structures d'accompagnement » $\rightarrow$ « \textbf{Malgré la création de structures d'accompagnement par l'État}, les jeunes peinent à accéder au financement. »
\item « Quand les jeunes entreprennent » $\rightarrow$ « \textbf{L'entrepreneuriat des jeunes} crée des emplois. »
\item « Parce que les banques refusent les prêts » $\rightarrow$ « \textbf{Le refus des prêts par les banques} provoque l'échec des projets. »
\item « Si l'on forme les jeunes » $\rightarrow$ « \textbf{La formation des jeunes} ferait diminuer le chômage. »
\end{enumerate}
\textbf{Point de vigilance :} la nominalisation supprime la conjonction ; il faut parfois la remplacer par une préposition adaptée (\emph{malgré} pour la concession, \emph{grâce à / à cause de} pour la cause).
\end{corrige}

\begin{corrige}[Exercice 7 — Contenus manifestes et latents]
\begin{enumerate}
\item \textbf{Contenu manifeste :} le Chef de Poste félicite officiellement Meka et affirme que la médaille prouve la reconnaissance de la France envers ses fidèles.
\item \textbf{Trois contenus latents :}
\begin{itemize}
\item La « civilisation » coloniale est en réalité une domination : l'expression « les bienfaits de la civilisation » est un \textbf{euphémisme} qui masque la spoliation des terres.
\item La médaille est une récompense dérisoire : « sa médaille est la preuve\ldots » fonctionne par \textbf{ironie} (chez Oyono qui fait parler le personnage) : une médaille ne compense ni les terres ni les fils perdus.
\item La « reconnaissance » de la France est intéressée et paternaliste : le possessif « \emph{ses} amis » et le superlatif de façade relèvent d'un \textbf{choix de vocabulaire} qui réduit les colonisés à des obligés.
\end{itemize}
\item \textbf{Indices :} euphémisme (« bienfaits »), ironie d'auteur (l'écart entre le discours et la réalité vécue par Meka), vocabulaire possessif et paternaliste (« ses amis »), modalisation affirmative excessive (« la preuve que ») qui invite le lecteur à douter.
\end{enumerate}
\end{corrige}

\begin{corrige}[Exercice 8 — Réduction d'un paragraphe]
\begin{enumerate}
\item \textbf{Idées essentielles (3) :} (a) l'entrepreneuriat est présenté comme solution au chômage des jeunes ; (b) les jeunes entrepreneurs rencontrent de nombreux obstacles (crédit, administration, formation) ; (c) la réussite exige l'accompagnement réel de l'État.
\item \textbf{Contraction possible :} « Présenté comme remède au chômage des jeunes, l'entrepreneuriat se heurte au crédit difficile, aux lourdeurs administratives et au manque de formation ; il exige donc un véritable accompagnement de l'État. »
\item \textbf{Dénombrement :} Présenté(1) comme(2) remède(3) au(4) chômage(5) des(6) jeunes(7), l'entrepreneuriat(8) se(9) heurte(10) au(11) crédit(12) difficile(13), aux(14) lourdeurs(15) administratives(16) et(17) au(18) manque(19) de(20) formation(21) ; il(22) exige(23) donc(24) un(25) véritable(26) accompagnement(27) de(28) l'État(29). Soit \textbf{29 mots} ; fourchette autorisée : $25 \pm 10\%$, c'est-à-dire de 22 à 27 mots. 29 dépasse : on supprime « donc un véritable » $\rightarrow$ « \ldots manque de formation ; il exige l'accompagnement de l'État. » Nouveau compte : 25 mots. \textbf{Contraction finale conforme : 25 mots.}
\end{enumerate}
\textbf{Point de vigilance :} toujours compter les mots (les mots composés à trait d'union comptent pour un) et annoncer le total ; un dépassement de la fourchette est sanctionné.
\end{corrige}

\begin{corrige}[Exercice 9 — Sujet type Bac : contraction d'un éditorial]
\textbf{1. Thèse et arguments (4 pts).} Thèse : l'entrepreneuriat des jeunes ne peut réussir que si les pouvoirs publics créent un environnement réellement favorable ; sinon il restera un slogan vide. Arguments attendus (deux parmi) : les jeunes n'accèdent pas au financement car les banques exigent des garanties ; les procédures administratives sont longues et décourageantes ; faire croire que « la volonté suffit » est un mensonge. \emph{Justification exigée :} citer ou reformuler les arguments du texte, pas d'opinion personnelle.

\textbf{2. Subordination et modes (4 pts).} « \textbf{Bien que} l'État \emph{ait créé}\ldots » : concession, subjonctif plus-que-parfait (obligatoire après \emph{bien que}). « \textbf{à condition que} les pouvoirs publics \emph{créent}\ldots » : condition, subjonctif présent. Le subjonctif exprime ici l'hypothèque, le non-assuré.

\textbf{3. Nominalisation (3 pts).} « \textbf{Malgré la création de structures d'accompagnement par l'État}, les jeunes peinent à accéder au financement. » (Ou : « La création de structures d'accompagnement par l'État n'empêche pas les jeunes de peiner à accéder au financement. »)

\textbf{4. Marqueurs de la structure argumentative (4 pts).} « \textbf{Certes}\ldots \textbf{Cependant} » : concession puis réfutation, l'auteur admet un point avant de le nuancer ; « \textbf{De plus} » : ajout d'un argument ; « \textbf{Faut-il\ldots ? Non. Mais\ldots} » : question rhétorique et réponse qui font avancer la démonstration ; « \textbf{Sinon} » : annonce de la conséquence négative, proche de la conclusion.

\textbf{5. Contraction en 50 mots ± 10 \% (5 pts).} Modèle : « L'entrepreneuriat, présenté comme remède au chômage des jeunes diplômés, se heurte au refus des banques et aux lourdeurs administratives. La volonté ne suffit pas : sans un environnement réellement favorable créé par les pouvoirs publics, il restera un slogan vide poussant la jeunesse vers l'émigration. » Dénombrement : 45 mots, dans la fourchette [45 ; 55]. \textbf{Annoncer le nombre de mots.}

\textbf{Points de vigilance :} respecter la fourchette (ici 45 à 55 mots ; au Bac, 81 à 99 mots pour 90) ; aucune opinion personnelle ; conserver la thèse et les arguments, supprimer exemples et répétitions ; employer des connecteurs et des nominalisations.
\end{corrige}

\begin{corrige}[Exercice 10 — Analyse textuelle : \emph{Le Vieux Nègre et la Médaille}]
\textbf{1. Contenus.} \emph{Manifeste :} après avoir tout sacrifié, Meka attend avec espoir la médaille, mais sa remise ne lui procure qu'une sensation froide et pesante. \emph{Latents :} (a) l'attente de la médaille est vitale et illusoire — indice : la \textbf{comparaison} « comme on attend la pluie après une longue sécheresse » ; (b) la reconnaissance coloniale est un leurre et une déception — indices : la gradation des sacrifices (« ses terres\ldots ses fils »), l'adverbe « \emph{enfin} » ironique, le \textbf{champ lexical} de la froideur et de la pesanteur (« poids », « métal froid ») opposé à la fête (« plus beaux habits »), et l'adverbe restrictif « \emph{seulement} ».

\textbf{2. Conjonctions.} Coordination : « \textbf{et} voilà que\ldots » (addition, enchaînement) ; « \textbf{Mais} quand on lui accrocha\ldots » (opposition, qui fait basculer le récit de l'espoir à la déception). Subordination : « \textbf{comme} on attend la pluie » (comparaison) ; on accepte aussi « \textbf{quand} on lui accrocha » (temps).

\textbf{3. Transformation.} « \textbf{Parce qu'}il avait donné ses terres et ses fils à la guerre, les Blancs allaient enfin le reconnaître, \textbf{si bien que} Meka attendait la médaille avec confiance. » (Subordonnée de cause avec \emph{parce que} + plus-que-parfait ; subordonnée de conséquence avec \emph{si bien que} ou \emph{de sorte que} + indicatif.)

\textbf{4. Paragraphe argumenté (modèle).} Dans cet extrait, Oyono fait de l'ironie une arme de dénonciation coloniale. D'abord, la comparaison initiale assimile l'attente de la médaille à celle de la pluie : la reconnaissance des Blancs devient un besoin vital, ce qui révèle l'aliénation du colonisé. Ensuite, la gradation « il avait donné ses terres, il avait donné ses fils » mesure l'énormité du sacrifice, que la médaille ne saurait compenser : l'écart entre le prix payé et la récompense constitue l'ironie amère du texte. Le basculement s'opère avec « Mais » : à la fête (« ses plus beaux habits ») succède la désillusion, signifiée par le champ lexical du froid et du poids (« le poids du métal froid »), tandis que « seulement » réduit la cérémonie à une sensation physique désagréable. Ainsi, sans jamais condamner ouvertement, Oyono laisse les faits et les images dénoncer l'imposture coloniale : la médaille, symbole de la « reconnaissance » de la France, n'est qu'un objet de métal qui écrase celui qui la reçoit.

\textbf{Barème indicatif (20 pts) :} contenus manifeste/latents 4 pts ; classement des conjonctions 4 pts ; transformation 4 pts ; paragraphe 8 pts (analyse 5, langue et liaisons 3). \textbf{Points de vigilance :} distinguer coordination et subordination ; citer précisément les indices linguistiques ; dans le paragraphe, articuler les idées avec des connecteurs et rester dans l'analyse du texte.
\end{corrige}

\begin{corrige}[Exercice 11 — Épreuve orale blanche : l'exposé]
\textbf{1. Fiche de préparation (modèle attendu).}
\emph{Problématique :} comment Ferdinand Oyono utilise-t-il l'ironie pour dénoncer l'imposture coloniale dans \emph{Le Vieux Nègre et la Médaille} ?
\emph{Introduction :} accroche (la remise de la médaille), présentation de l'œuvre et de l'auteur, problématique, annonce du plan.
\emph{Plan :} I. Une ironie des situations : les sacrifices de Meka (terres, fils) payés d'une médaille dérisoire ; la cérémonie qui tourne à l'humiliation. II. Une ironie des discours et des personnages : le paternalisme du Chef de Poste (« les bienfaits de la civilisation »), l'écart entre les paroles officielles et la réalité vécue. III. La portée critique de l'ironie : réveiller la conscience du lecteur, démasquer la colonisation sans discours direct.
\emph{Conclusion :} bilan (l'ironie comme arme efficace car elle fait réfléchir) et ouverture (l'ironie chez Mongo Beti, ou la littérature engagée africaine).

\textbf{2. Support visuel (3 diapositives).}
\begin{itemize}
\item \textbf{Diapo 1 — « Meka : le prix de la médaille »} : mots-clés : terres données, fils à la guerre, médaille ; citation courte : « le poids du métal froid ».
\item \textbf{Diapo 2 — « L'ironie des discours »} : mots-clés : paternalisme, euphémisme ; citation : « les bienfaits de la civilisation » ; écart discours/réalité.
\item \textbf{Diapo 3 — « Portée »} : mots-clés : dénonciation, lucidité du lecteur ; ouverture : littérature anticoloniale africaine.
\end{itemize}
Règle : aucun texte long, uniquement mots-clés et exemples.

\textbf{3. Passation et évaluation.} L'exposé est noté sur 20 selon la grille : contenu /4, organisation /4, langue /4, communication /4, réponses /4. \textbf{Points de vigilance :} respecter les 5 minutes (s'entraîner avec un chronomètre) ; annoncer le plan et le rappeler en conclusion ; soigner les liaisons logiques (\emph{d'abord, ensuite, enfin, cependant}) ; regarder l'auditoire ; pour les questions du jury, répondre en justifiant par l'œuvre (exemples précis). Questions probables : « Pourquoi l'ironie est-elle plus efficace qu'une accusation directe ? » ; « Meka est-il conscient de l'ironie de sa situation ? »
\end{corrige}

\newpage

\coursec{Fiche bilan}

\begin{popfigure}
\begin{tikzpicture}[mindmap, grow cyclic, every node/.style={concept, text=white, font=\small},
root concept/.append style={concept color=popblue, font=\bfseries\small},
level 1 concept/.append style={level distance=3.4cm, sibling angle=60, font=\scriptsize},
level 2 concept/.append style={level distance=2.2cm, font=\tiny}]
\node[root concept]{Formuler les idées essentielles d'un texte à résumer}
child[concept color=poporange]{ node {Texte argumentatif}
  child { node {Thèse, arguments, exemples} }
  child { node {Connecteurs logiques} }
}
child[concept color=popgreen]{ node {Lecture méthodique}
  child { node {\emph{Le Vieux nègre et la médaille} (lectures 4 et 5)} }
}
child[concept color=poppurple]{ node {Contraction de texte}
  child { node {Repérer la structure} }
  child { node {Réduire : 1/4 du texte} }
}
child[concept color=poppink]{ node {Reformulation}
  child { node {Ses propres mots} }
  child { node {Fidélité au sens} }
}
child[concept color=popgold]{ node {Outils de liaison}
  child { node {Coordination : mais, ou, et, donc, or, ni, car} }
  child { node {Subordination : que, parce que, bien que, si} }
}
child[concept color=popblue]{ node {Contenus manifeste et latent}
  child { node {Dit et sous-entendu} }
};
\end{tikzpicture}
\legende{Carte mentale de la leçon : les six compétences à mobiliser pour formuler les idées essentielles d'un texte à résumer.}
\end{popfigure}

\begin{bilan}
\textbf{Définitions clés}
\begin{itemize}
\item \cle{Idée essentielle} : idée sans laquelle le sens global du texte disparaît ; elle porte la thèse ou un argument majeur de l'auteur.
\item \cle{Contraction de texte} : exercice qui consiste à réduire un texte au quart de sa longueur (au Bac technique, environ 100 mots pour un texte de 400 mots), sans commentaire personnel ni jugement.
\item \cle{Reformulation} : dire la même chose avec ses propres mots, en gardant le sens, le ton et l'ordre des idées de l'auteur.
\item \cle{Texte argumentatif} : texte qui défend une \cle{thèse} à l'aide d'\cle{arguments} illustrés par des \cle{exemples}.
\item \cle{Contenu manifeste} : ce qui est dit explicitement ; \cle{contenu latent} : ce qui est suggéré, sous-entendu (implicite).
\end{itemize}

\textbf{Repères à connaître par cœur}
\begin{center}
\begin{tabularx}{\linewidth}{|l|X|}\hline
\rowcolor{popblueL} Notion & À retenir \\ \hline
Conjonctions de coordination & mais, ou, et, donc, or, ni, car : relient des éléments de même fonction \\ \hline
Conjonctions de subordination & que, parce que, bien que, si, lorsque, puisque : introduisent une subordonnée \\ \hline
Étapes du résumé & Lire $\rightarrow$ dégager le thème $\rightarrow$ repérer la thèse $\rightarrow$ dégager les idées essentielles $\rightarrow$ reformuler $\rightarrow$ vérifier \\ \hline
Règle du quart & Longueur du résumé = nombre de mots du texte $\div$ 4 (marge de $\pm$ 10\,\%) \\ \hline
Interdits du résumé & Copier des phrases, citer, donner son avis, ajouter des idées, utiliser « je » \\ \hline
\end{tabularx}
\end{center}

\textbf{Méthodes express}
\begin{itemize}
\item \cle{Repérage} : souligner une idée par paragraphe ; éliminer exemples, répétitions, digressions, figures d'insistance.
\item \cle{Réduction} : supprimer les expansions du nom (adjectifs, compléments), remplacer une énumération par un mot générique, condenser une subordonnée en nom (« parce qu'il était pauvre » $\rightarrow$ « sa pauvreté »).
\item \cle{Reformulation} : changer le vocabulaire (synonymes), changer la construction (passif/actif, phrase nominale), mais jamais le sens.
\item \cle{Exposé oral} : préparer un plan (introduction, 2 ou 3 parties, conclusion), des fiches courtes, et soigner voix, regard et gestuelle.
\end{itemize}
\end{bilan}

\begin{aretenir}[Checklist avant le Bac]
\begin{itemize}
\item[$\square$] Je sais définir la contraction de texte et la distinguer du résumé-commentaire.
\item[$\square$] Je sais calculer le quart du texte et compter mes mots.
\item[$\square$] Je repère la thèse de l'auteur avant de rédiger.
\item[$\square$] Je dégage une idée essentielle par unité de sens.
\item[$\square$] J'élimine exemples, citations, répétitions et détails secondaires.
\item[$\square$] Je reformule avec mes propres mots sans trahir le sens.
\item[$\square$] Je relie les idées avec les bons connecteurs (coordination et subordination).
\item[$\square$] Je rédige à la troisième personne, sans donner mon avis.
\item[$\square$] Je distingue contenu manifeste et contenu latent dans un texte.
\item[$\square$] Je connais les grandes étapes de la lecture méthodique d'\emph{un extrait du Vieux nègre et la médaille} (situation, thème, procédés, sens).
\item[$\square$] Je sais préparer et présenter un exposé avec un plan clair.
\item[$\square$] Je relis mon résumé pour vérifier grammaire, orthographe et longueur.
\end{itemize}
\end{aretenir}

\findecours

\end{document}
